\chapter{Binary Search}
\chapterepigraph{Binary search is not a trick for sorted arrays. It is a
trick for any space of candidate answers in which feasibility is monotonic
-- once you see that monotonicity, the array itself becomes optional.}

\section{What Is Binary Search, Really?}
\label{sec:bs-intro}

Most students first meet binary search as a way to find a target value in a
sorted array: compare the middle element to the target, and discard half the
array based on the result. That description is correct, but it is also
dangerously narrow, because it hides the actual reason binary search works
behind the specific setting of ``a sorted array.'' The real requirement is
much more general, and recognizing it is what lets you apply binary search
to problems that do not look anything like array search at first glance.

\begin{ideabox}[The One Property Binary Search Actually Needs]
Binary search applies whenever you are searching over a totally ordered
space of candidates (indices, integers, real numbers, even configurations)
on which some predicate $P$ is \textbf{monotonic} -- that is, $P$ is
\texttt{false} for a prefix of the space and \texttt{true} for the remaining
suffix (or vice versa), with a single crossover point. If such a predicate
exists, binary search finds the crossover point in $O(\log(\text{range}))$
evaluations of $P$, regardless of what $P$ actually computes or how
expensive each evaluation is.
\end{ideabox}

In classical array search, the candidate space is the array's indices
$0, 1, \dots, n-1$, and the predicate is $P(i) = (\textit{arr}[i] \ge
\textit{target})$ -- which is monotonic precisely because the array is
sorted. But the same machinery works just as well when the candidate space
is, say, every integer capacity a ship could have, or every possible eating
speed for Koko's bananas, or every real number that could be a square root --
as long as we can answer ``does this candidate satisfy the requirement, or
not'' in a way that flips exactly once as the candidate increases. This
chapter is organized around exactly that distinction.

\subsection*{Three Flavors of Binary Search in This Chapter}

\begin{patternbox}[How to Tell Which Flavor You Are Facing]
\begin{itemize}
  \item \textbf{Binary search on a 1D array
        (Sections~\ref{sec:bs-find-x}--\ref{sec:bs-peak-1d}).} The keywords
        are usually \emph{``sorted array,''} \emph{``rotated sorted
        array,''} or \emph{``find an index/value such that...''} The
        candidate space is literally the array's indices, and $\textit{lo},
        \textit{hi}$ start as $0$ and $n-1$.
  \item \textbf{Binary search on the answer
        (Sections~\ref{sec:bs-sqrt}--\ref{sec:bs-kth-element}).} The
        keywords are \emph{``minimize the maximum ...''},
        \emph{``maximize the minimum ...''}, or \emph{``find the smallest/
        largest value such that some condition holds.''} Here there may be
        no array to search at all -- the candidate space is a numeric
        range of possible answers (e.g.\ every integer from $1$ to
        $10^9$), and we binary search directly on that range, using a
        helper feasibility check (often itself an $O(n)$ scan) as the
        monotonic predicate.
  \item \textbf{Binary search on a 2D array
        (Sections~\ref{sec:bs-row-max-ones}--\ref{sec:bs-matrix-median}).}
        The keywords are \emph{``row-wise and/or column-wise sorted
        matrix.''} Here the candidate space is two-dimensional, and the
        technique is usually to either flatten the 2D index space into a
        1D one, or to run a 1D binary search independently along each row
        or column, or to binary search on the answer (as above) using a
        2D-aware feasibility check.
\end{itemize}
The single most important habit to build in this chapter is translating a
problem's wording into ``what is the candidate space, and what is the
monotonic predicate that flips exactly once across it'' \emph{before}
writing any code. Once those two things are identified, the
\texttt{while(lo <= hi)} skeleton is almost always the same.
\end{patternbox}

\subsection*{The Skeleton, and Where Bugs Actually Come From}

Nearly every binary search bug comes from getting one of three things
wrong: the loop condition (\texttt{lo <= hi} vs.\ \texttt{lo < hi}), the
update after a comparison (\texttt{mid - 1} / \texttt{mid + 1} vs.\
\texttt{mid} itself), or an integer-overflow in the midpoint computation.
We will use the convention
\[
  \textit{mid} = \textit{lo} + (\textit{hi} - \textit{lo}) / 2
\]
throughout this chapter, rather than $(\textit{lo}+\textit{hi})/2$, because
it avoids overflow when $\textit{lo}$ and $\textit{hi}$ are both close to
the maximum representable integer -- a habit worth having permanently, even
in problems where the bound is small enough that it would not matter today.

Each section in this chapter follows the same structure used in the Greedy
chapter: a precise problem statement and the keywords that identify it;
identification of the candidate space and the monotonic predicate; the
algorithm in words and in formal pseudocode; a hand-traced dry run; complete
compilable C++ code; and a careful time/space complexity analysis.

\newpage

\section{Binary Search to Find X in a Sorted Array}
\label{sec:bs-find-x}

\problemstatement{We are given an array $\textit{arr}$ of $n$ elements,
sorted in increasing order, and a target value $X$. We must determine the
index of $X$ in the array if it is present, and return $-1$ if it is not.
We must do this faster than the obvious $O(n)$ linear scan.}

\begin{patternbox}
\emph{``Sorted array''} appearing anywhere in a problem statement,
combined with a request to \emph{find/locate} a specific value or the
boundary of some property, is the single most reliable trigger for binary
search in this entire chapter. The reason a sorted array enables binary
search is precisely the monotonicity property from
Section~\ref{sec:bs-intro}: define $P(i) = (\textit{arr}[i] \ge X)$; because
the array is sorted, $P$ is \texttt{false} for every index before $X$'s
position and \texttt{true} from $X$'s position onward (or, if $X$ is
absent, from the position where elements first exceed $X$ onward) -- a
single crossover, which is exactly what binary search is built to find.
\end{patternbox}

\subsection*{Understanding the problem carefully}

If the array were unsorted, knowing that $\textit{arr}[\textit{mid}] \ne X$
would tell us nothing about where $X$ might be -- it could be anywhere
among the remaining elements, so we would have no choice but to check them
all. Sortedness changes this completely: if
$\textit{arr}[\textit{mid}] < X$, then because every element to the left of
$\textit{mid}$ is even smaller (the array is sorted in increasing order),
\emph{none} of them can equal $X$ either, so the entire left half can be
discarded in one step. Symmetrically, if $\textit{arr}[\textit{mid}] > X$,
the entire right half can be discarded. Only when
$\textit{arr}[\textit{mid}] = X$ do we stop, having found the answer.

\begin{ideabox}[Candidate Space and Predicate]
The candidate space is the set of indices $\{0, 1, \dots, n-1\}$. We
maintain a search window $[\textit{lo}, \textit{hi}]$, initially the whole
array, and repeatedly look at the midpoint, using the comparison with $X$ to
discard exactly the half of the window that provably cannot contain $X$.
\end{ideabox}

\subsection*{Why halving the search space each time is correct}

Each comparison at $\textit{mid}$ produces one of exactly three outcomes,
and in every outcome we either find the answer immediately or can prove
that an entire half of the remaining window is irrelevant. Since no element
is ever incorrectly discarded (the sortedness guarantee makes every
discard provably safe) and the window strictly shrinks by roughly half on
every iteration, the process must terminate, either by finding $X$ or by
the window becoming empty ($\textit{lo} > \textit{hi}$), at which point we
can safely conclude $X$ is not present.

\subsection*{Algorithm}

\begin{enumerate}
  \item Initialize $\textit{lo} \gets 0$, $\textit{hi} \gets n-1$.
  \item While $\textit{lo} \le \textit{hi}$:
  \begin{enumerate}
    \item $\textit{mid} \gets \textit{lo} + (\textit{hi}-\textit{lo})/2$.
    \item If $\textit{arr}[\textit{mid}] = X$: return $\textit{mid}$.
    \item If $\textit{arr}[\textit{mid}] < X$: $\textit{lo} \gets
          \textit{mid}+1$.
    \item Otherwise: $\textit{hi} \gets \textit{mid}-1$.
  \end{enumerate}
  \item If the loop ends without returning, return $-1$.
\end{enumerate}

\begin{algorithm}[H]
\caption{Binary Search}
\begin{algorithmic}[1]
\Procedure{BinarySearch}{$\textit{arr}, X$}
  \State $\textit{lo} \gets 0$, $\textit{hi} \gets n-1$
  \While{$\textit{lo} \le \textit{hi}$}
    \State $\textit{mid} \gets \textit{lo} + (\textit{hi}-\textit{lo})/2$
    \If{$\textit{arr}[\textit{mid}] = X$}
      \State \Return \textit{mid}
    \ElsIf{$\textit{arr}[\textit{mid}] < X$}
      \State $\textit{lo} \gets \textit{mid} + 1$
    \Else
      \State $\textit{hi} \gets \textit{mid} - 1$
    \EndIf
  \EndWhile
  \State \Return $-1$
\EndProcedure
\end{algorithmic}
\end{algorithm}

\subsection*{Dry run}

\dryrun{Take $\textit{arr} = [-1, 0, 3, 5, 9, 12]$, $X = 9$.}

\begin{itemize}
  \item $\textit{lo}=0,\textit{hi}=5$: $\textit{mid}=2$,
        $\textit{arr}[2]=3 < 9$. $\textit{lo}=3$.
  \item $\textit{lo}=3,\textit{hi}=5$: $\textit{mid}=4$,
        $\textit{arr}[4]=9 = 9$. Return $4$.
\end{itemize}

The answer is index $4$.

Now take $X = 2$ on the same array.

\begin{itemize}
  \item $\textit{lo}=0,\textit{hi}=5$: $\textit{mid}=2$,
        $\textit{arr}[2]=3 > 2$. $\textit{hi}=1$.
  \item $\textit{lo}=0,\textit{hi}=1$: $\textit{mid}=0$,
        $\textit{arr}[0]=-1 < 2$. $\textit{lo}=1$.
  \item $\textit{lo}=1,\textit{hi}=1$: $\textit{mid}=1$,
        $\textit{arr}[1]=0 < 2$. $\textit{lo}=2$.
  \item $\textit{lo}=2 > \textit{hi}=1$: loop ends. Return $-1$.
\end{itemize}

The answer is $-1$: $2$ is not present in the array.

\subsection*{C++ implementation (iterative and recursive)}

\begin{lstlisting}
#include <bits/stdc++.h>
using namespace std;

// Iterative version
int binarySearch(vector<int>& arr, int X) {
    int lo = 0, hi = arr.size() - 1;

    while (lo <= hi) {
        int mid = lo + (hi - lo) / 2;
        if (arr[mid] == X) return mid;
        else if (arr[mid] < X) lo = mid + 1;
        else hi = mid - 1;
    }
    return -1;
}

// Recursive version
int binarySearchRecursive(vector<int>& arr, int lo, int hi, int X) {
    if (lo > hi) return -1;

    int mid = lo + (hi - lo) / 2;
    if (arr[mid] == X) return mid;
    else if (arr[mid] < X) return binarySearchRecursive(arr, mid + 1, hi, X);
    else return binarySearchRecursive(arr, lo, mid - 1, X);
}

int main() {
    vector<int> arr = {-1, 0, 3, 5, 9, 12};
    cout << "Iterative: " << binarySearch(arr, 9) << endl;
    cout << "Recursive: "
         << binarySearchRecursive(arr, 0, arr.size() - 1, 9) << endl;
    return 0;
}
\end{lstlisting}

\begin{complexitybox}
\textbf{Time Complexity.} Each iteration (or recursive call) discards at
least half of the remaining search window, so the window shrinks from size
$n$ to size $0$ in at most $\lceil \log_2 n \rceil + 1$ steps, giving
\[
  O(\log n).
\]
\textbf{Space Complexity.} The iterative version uses only the three
scalar variables \textit{lo}, \textit{hi}, \textit{mid}, giving $O(1)$
space. The recursive version uses $O(\log n)$ space for the call stack,
since the recursion depth equals the number of halvings performed.
\end{complexitybox}

\begin{takeawaybox}
Binary search's correctness rests entirely on being able to \emph{prove},
at each step, that an entire half of the remaining search space cannot
contain the answer. In plain array search this proof comes directly from
sortedness; every other binary search problem in this chapter is really
asking you to find or construct an equivalent monotonicity proof in a
different setting.
\end{takeawaybox}

\section{Implement Lower Bound}
\label{sec:bs-lower-bound}

\problemstatement{We are given a sorted array $\textit{arr}$ of $n$
elements and a value $X$. We must find the \textbf{lower bound} of $X$:
the index of the first element in $\textit{arr}$ that is $\textbf{greater
than or equal to}$ $X$. If no such element exists (every element is smaller
than $X$), the lower bound is $n$ (one past the end of the array).}

\begin{patternbox}
\emph{``Find the first/smallest index such that some condition holds''} --
here, the condition is $\textit{arr}[i] \ge X$ -- is the
\textbf{find-the-boundary} variant of binary search, distinct from the
exact-match search of Section~\ref{sec:bs-find-x}. The keyword is not
``find $X$'' but \textbf{``find the first position where a property
becomes true.''} This phrasing recurs constantly throughout this chapter
(first/last occurrence, search insert position, floor/ceil, and many
``binary search on the answer'' problems), so internalizing the
boundary-search template here pays off many times over.
\end{patternbox}

\subsection*{Understanding the problem carefully}

Define the predicate $P(i) = (\textit{arr}[i] \ge X)$. Because the array is
sorted in increasing order, $P$ is \texttt{false} for some prefix of
indices and \texttt{true} for the remaining suffix, with exactly one
crossover point -- the lower bound is, by definition, that crossover index.
Unlike Section~\ref{sec:bs-find-x}, we are not looking for a single exact
match; we are looking for the \emph{boundary} between ``too small'' and
``big enough,'' which exists regardless of whether $X$ itself is actually
present in the array.

\begin{ideabox}[Candidate Space and Predicate]
The candidate space is indices $0, \dots, n$ (note: $n$ itself is a valid
answer, representing ``no element is large enough''). The predicate is
$P(i) = (\textit{arr}[i] \ge X)$ for $i < n$, and we want the smallest index
where $P$ holds.
\end{ideabox}

\subsection*{Why we keep a separate ``best answer so far'' variable}

In exact-match search, the loop stops the instant we find a match. Here, we
cannot stop early, because finding \emph{one} index where $\textit{arr}[i]
\ge X$ does not tell us it is the \emph{first} such index -- there might be
an even smaller valid index hiding in the left half. So instead of
returning immediately, whenever $\textit{arr}[\textit{mid}] \ge X$ we
record $\textit{mid}$ as a candidate answer and \emph{keep searching to the
left}, hoping to find an even smaller valid index; whenever
$\textit{arr}[\textit{mid}] < X$, we know the entire left portion up to and
including $\textit{mid}$ is too small, so we discard it and search to the
right. This is the general template for every ``find the first/last index
satisfying a monotonic property'' problem: maintain a running best answer,
and keep shrinking the window based on the predicate, rather than stopping
on the first observed success.

\subsection*{Algorithm}

\begin{enumerate}
  \item Initialize $\textit{lo} \gets 0$, $\textit{hi} \gets n-1$,
        $\textit{ans} \gets n$.
  \item While $\textit{lo} \le \textit{hi}$:
  \begin{enumerate}
    \item $\textit{mid} \gets \textit{lo} + (\textit{hi}-\textit{lo})/2$.
    \item If $\textit{arr}[\textit{mid}] \ge X$: record
          $\textit{ans} \gets \textit{mid}$, and search left:
          $\textit{hi} \gets \textit{mid}-1$.
    \item Otherwise: search right: $\textit{lo} \gets \textit{mid}+1$.
  \end{enumerate}
  \item Return \textit{ans}.
\end{enumerate}

\begin{algorithm}[H]
\caption{Lower Bound}
\begin{algorithmic}[1]
\Procedure{LowerBound}{$\textit{arr}, X$}
  \State $\textit{lo} \gets 0$, $\textit{hi} \gets n-1$, $\textit{ans} \gets n$
  \While{$\textit{lo} \le \textit{hi}$}
    \State $\textit{mid} \gets \textit{lo} + (\textit{hi}-\textit{lo})/2$
    \If{$\textit{arr}[\textit{mid}] \ge X$}
      \State $\textit{ans} \gets \textit{mid}$
      \State $\textit{hi} \gets \textit{mid} - 1$
    \Else
      \State $\textit{lo} \gets \textit{mid} + 1$
    \EndIf
  \EndWhile
  \State \Return \textit{ans}
\EndProcedure
\end{algorithmic}
\end{algorithm}

\subsection*{Dry run}

\dryrun{Take $\textit{arr} = [1, 2, 2, 3]$, $X = 2$.}

\begin{itemize}
  \item $\textit{lo}=0,\textit{hi}=3,\textit{ans}=4$: $\textit{mid}=1$,
        $\textit{arr}[1]=2 \ge 2$. $\textit{ans}=1$, $\textit{hi}=0$.
  \item $\textit{lo}=0,\textit{hi}=0$: $\textit{mid}=0$,
        $\textit{arr}[0]=1 \ge 2$? No. $\textit{lo}=1$.
  \item $\textit{lo}=1 > \textit{hi}=0$: loop ends.
\end{itemize}

The lower bound is index $1$ -- the first occurrence of $2$.

Now take $X = 5$ on the same array (larger than every element).

\begin{itemize}
  \item $\textit{lo}=0,\textit{hi}=3,\textit{ans}=4$: $\textit{mid}=1$,
        $\textit{arr}[1]=2 \ge 5$? No. $\textit{lo}=2$.
  \item $\textit{lo}=2,\textit{hi}=3$: $\textit{mid}=2$,
        $\textit{arr}[2]=2 \ge 5$? No. $\textit{lo}=3$.
  \item $\textit{lo}=3,\textit{hi}=3$: $\textit{mid}=3$,
        $\textit{arr}[3]=3 \ge 5$? No. $\textit{lo}=4$.
  \item $\textit{lo}=4 > \textit{hi}=3$: loop ends. $\textit{ans}$ remains
        $4$ (its initial value, never overwritten).
\end{itemize}

The lower bound is $4 = n$, correctly signaling that no element is large
enough.

\subsection*{C++ implementation}

\begin{lstlisting}
#include <bits/stdc++.h>
using namespace std;

int lowerBound(vector<int>& arr, int X) {
    int n = arr.size();
    int lo = 0, hi = n - 1, ans = n;

    while (lo <= hi) {
        int mid = lo + (hi - lo) / 2;
        if (arr[mid] >= X) {
            ans = mid;
            hi = mid - 1;
        } else {
            lo = mid + 1;
        }
    }
    return ans;
}

int main() {
    vector<int> arr = {1, 2, 2, 3};
    cout << "Lower bound of 2: " << lowerBound(arr, 2) << endl;
    cout << "Lower bound of 5: " << lowerBound(arr, 5) << endl;
    return 0;
}
\end{lstlisting}

\begin{complexitybox}
\textbf{Time Complexity.} The search window halves each iteration, giving
\[
  O(\log n).
\]
\textbf{Space Complexity.} Only a constant number of scalar variables are
used, giving $O(1)$ space.
\end{complexitybox}

\begin{takeawaybox}
Whenever a problem asks for the \emph{first} (or \emph{last}) index
satisfying a monotonic condition, do not stop the binary search the moment
you observe the condition holding -- instead, record it as a tentative
answer and keep narrowing the search toward the boundary, in the direction
that could produce an even better (earlier, or later) answer. This
``record and keep narrowing'' template is reused, almost verbatim, in
upper bound, floor/ceil, and first/last occurrence.
\end{takeawaybox}

\section{Implement Upper Bound}
\label{sec:bs-upper-bound}

\problemstatement{We are given a sorted array $\textit{arr}$ of $n$
elements and a value $X$. We must find the \textbf{upper bound} of $X$:
the index of the first element in $\textit{arr}$ that is
\textbf{strictly greater than} $X$. If no such element exists, the upper
bound is $n$.}

\begin{patternbox}
Upper bound is lower bound's twin, with one word changed: ``greater than or
equal to'' becomes ``strictly greater than.'' Whenever you see \emph{``first
index strictly greater than $X$''} as opposed to \emph{``first index at
least $X$,''} you are being asked for an upper bound rather than a lower
bound. The two are easy to confuse under time pressure, so it is worth
fixing the distinction precisely: lower bound finds where $X$ \emph{could
first appear}; upper bound finds where the elements equal to $X$
\emph{end}.
\end{patternbox}

\subsection*{Understanding the problem carefully}

The predicate here is $P(i) = (\textit{arr}[i] > X)$, which, exactly as in
lower bound, is \texttt{false} for a prefix and \texttt{true} for a suffix
of the sorted array, with a single crossover. The only change from
Section~\ref{sec:bs-lower-bound} is the strictness of the comparison: an
element exactly equal to $X$ satisfies the lower-bound predicate
($\ge X$) but \emph{not} the upper-bound predicate ($> X$). Consequently,
if $X$ appears multiple times in the array, the lower bound points to the
first occurrence of $X$, while the upper bound points to the index
\emph{immediately after} the last occurrence of $X$ -- a fact we will use
directly in Section~\ref{sec:bs-first-last} and
Section~\ref{sec:bs-count-occurrences}.

\begin{ideabox}[Candidate Space and Predicate]
The candidate space is indices $0, \dots, n$. The predicate is
$P(i) = (\textit{arr}[i] > X)$ for $i < n$, and we want the smallest index
where $P$ holds.
\end{ideabox}

\subsection*{Why the same template applies, with one flipped comparison}

The reasoning is identical to lower bound: whenever
$\textit{arr}[\textit{mid}] > X$, this index is a valid candidate for the
upper bound, but there might be an even earlier index that is also strictly
greater than $X$, so we record it and search left. Whenever
$\textit{arr}[\textit{mid}] \le X$ (note: this now includes equality, unlike
lower bound), this index and everything before it is disqualified, so we
search right. The only code change relative to lower bound is replacing the
comparison \texttt{arr[mid] >= X} with \texttt{arr[mid] > X}.

\subsection*{Algorithm}

\begin{enumerate}
  \item Initialize $\textit{lo} \gets 0$, $\textit{hi} \gets n-1$,
        $\textit{ans} \gets n$.
  \item While $\textit{lo} \le \textit{hi}$:
  \begin{enumerate}
    \item $\textit{mid} \gets \textit{lo} + (\textit{hi}-\textit{lo})/2$.
    \item If $\textit{arr}[\textit{mid}] > X$: record
          $\textit{ans} \gets \textit{mid}$; $\textit{hi} \gets
          \textit{mid}-1$.
    \item Otherwise: $\textit{lo} \gets \textit{mid}+1$.
  \end{enumerate}
  \item Return \textit{ans}.
\end{enumerate}

\begin{algorithm}[H]
\caption{Upper Bound}
\begin{algorithmic}[1]
\Procedure{UpperBound}{$\textit{arr}, X$}
  \State $\textit{lo} \gets 0$, $\textit{hi} \gets n-1$, $\textit{ans} \gets n$
  \While{$\textit{lo} \le \textit{hi}$}
    \State $\textit{mid} \gets \textit{lo} + (\textit{hi}-\textit{lo})/2$
    \If{$\textit{arr}[\textit{mid}] > X$}
      \State $\textit{ans} \gets \textit{mid}$
      \State $\textit{hi} \gets \textit{mid} - 1$
    \Else
      \State $\textit{lo} \gets \textit{mid} + 1$
    \EndIf
  \EndWhile
  \State \Return \textit{ans}
\EndProcedure
\end{algorithmic}
\end{algorithm}

\subsection*{Dry run}

\dryrun{Take $\textit{arr} = [1, 2, 2, 3]$, $X = 2$.}

\begin{itemize}
  \item $\textit{lo}=0,\textit{hi}=3,\textit{ans}=4$: $\textit{mid}=1$,
        $\textit{arr}[1]=2 > 2$? No. $\textit{lo}=2$.
  \item $\textit{lo}=2,\textit{hi}=3$: $\textit{mid}=2$,
        $\textit{arr}[2]=2 > 2$? No. $\textit{lo}=3$.
  \item $\textit{lo}=3,\textit{hi}=3$: $\textit{mid}=3$,
        $\textit{arr}[3]=3 > 2$? Yes. $\textit{ans}=3$, $\textit{hi}=2$.
  \item $\textit{lo}=3 > \textit{hi}=2$: loop ends.
\end{itemize}

The upper bound is index $3$ -- one past the last occurrence of $2$
(which sits at index $2$), consistent with the lower bound of $2$ found in
Section~\ref{sec:bs-lower-bound} being index $1$: the occurrences of $2$
span indices $[1, 2]$.

\subsection*{C++ implementation}

\begin{lstlisting}
#include <bits/stdc++.h>
using namespace std;

int upperBound(vector<int>& arr, int X) {
    int n = arr.size();
    int lo = 0, hi = n - 1, ans = n;

    while (lo <= hi) {
        int mid = lo + (hi - lo) / 2;
        if (arr[mid] > X) {
            ans = mid;
            hi = mid - 1;
        } else {
            lo = mid + 1;
        }
    }
    return ans;
}

int main() {
    vector<int> arr = {1, 2, 2, 3};
    cout << "Upper bound of 2: " << upperBound(arr, 2) << endl;
    return 0;
}
\end{lstlisting}

\begin{complexitybox}
\textbf{Time Complexity.} $O(\log n)$, identical to lower bound, since
the algorithm structure is unchanged.
\textbf{Space Complexity.} $O(1)$ auxiliary space.
\end{complexitybox}

\begin{takeawaybox}
Lower bound and upper bound differ by exactly one comparison operator
($\ge$ versus $>$), but that single difference is what makes the pair
together so useful: the number of occurrences of $X$ in a sorted array is
always $\textit{upperBound}(X) - \textit{lowerBound}(X)$, and the first and
last occurrence indices of $X$ are $\textit{lowerBound}(X)$ and
$\textit{upperBound}(X) - 1$ respectively (when $X$ is present at all).
\end{takeawaybox}

\section{Search Insert Position}
\label{sec:bs-search-insert}

\problemstatement{We are given a sorted array $\textit{arr}$ of $n$
distinct integers and a target value $X$. If $X$ is found in the array,
return its index. If not, return the index at which $X$ \textbf{would be
inserted} to keep the array sorted.}

\begin{patternbox}
This problem is worth pausing on specifically \emph{because} it looks new
but is not: \emph{``the index where $X$ would be inserted to keep the array
sorted''} is, word for word, a description of the \textbf{lower bound} of
$X$ from Section~\ref{sec:bs-lower-bound}. The position just before the
first element $\ge X$ is exactly where $X$ slots in. Recognizing when a
new-sounding problem statement is a known pattern in disguise is as
important a skill in this chapter as deriving a pattern from scratch.
\end{patternbox}

\subsection*{Understanding the problem carefully}

If $X$ is present in the array, the lower bound of $X$ is precisely the
index of $X$ (since the array has distinct elements, the first index with
$\textit{arr}[i] \ge X$ is the unique index where $\textit{arr}[i] = X$).
If $X$ is absent, the lower bound is the first index whose value exceeds
$X$ -- which is exactly the index where $X$ should be inserted, since
everything before that index is smaller than $X$ and everything from that
index onward is larger.

\begin{ideabox}[Candidate Space and Predicate]
Identical to lower bound: candidate space is indices $0, \dots, n$,
predicate $P(i) = (\textit{arr}[i] \ge X)$, answer is the smallest index
where $P$ holds.
\end{ideabox}

\subsection*{Why no new derivation is needed}

There is nothing new to prove here -- this section exists to train the
habit of \emph{reducing} an unfamiliar-sounding problem to a pattern you
already have a proven, working template for, rather than re-deriving binary
search bounds from scratch every time a problem uses different words for
the same underlying question.

\subsection*{Algorithm}

Identical to Algorithm~\ref{sec:bs-lower-bound}'s lower-bound procedure,
called directly on $X$.

\subsection*{Dry run}

\dryrun{Take $\textit{arr} = [1, 3, 5, 6]$, $X = 5$.}

This is present at index $2$; the lower bound of $5$ is index $2$
(verify: $\textit{arr}[2]=5 \ge 5$, and $\textit{arr}[1]=3 < 5$), so the
answer is $2$.

Now take $X = 2$ on the same array. $2$ is absent. The lower bound of $2$
is index $1$ (the first element $\ge 2$ is $\textit{arr}[1]=3$), so $2$
would be inserted at index $1$, giving $[1,2,3,5,6]$ -- correct.

\subsection*{C++ implementation}

\begin{lstlisting}
#include <bits/stdc++.h>
using namespace std;

int searchInsert(vector<int>& arr, int X) {
    int n = arr.size();
    int lo = 0, hi = n - 1, ans = n;

    while (lo <= hi) {
        int mid = lo + (hi - lo) / 2;
        if (arr[mid] >= X) {
            ans = mid;
            hi = mid - 1;
        } else {
            lo = mid + 1;
        }
    }
    return ans;
}

int main() {
    vector<int> arr = {1, 3, 5, 6};
    cout << "Insert position for 5: " << searchInsert(arr, 5) << endl;
    cout << "Insert position for 2: " << searchInsert(arr, 2) << endl;
    return 0;
}
\end{lstlisting}

\begin{complexitybox}
\textbf{Time Complexity.} $O(\log n)$.
\textbf{Space Complexity.} $O(1)$.
\end{complexitybox}

\begin{takeawaybox}
Before deriving a new binary search predicate from scratch, check whether
the problem is just lower bound, upper bound, or a simple combination of
the two, restated in different words. Many ``new'' binary search problems
on interview sheets are exactly this.
\end{takeawaybox}

\section{Floor and Ceil in a Sorted Array}
\label{sec:bs-floor-ceil}

\problemstatement{We are given a sorted array $\textit{arr}$ of $n$
elements and a value $X$. We must find the \textbf{floor} of $X$ -- the
largest element in $\textit{arr}$ that is $\le X$ -- and the
\textbf{ceil} of $X$ -- the smallest element in $\textit{arr}$ that is
$\ge X$. If no floor (resp.\ ceil) exists, report $-1$.}

\begin{patternbox}
\emph{``Floor''} and \emph{``ceil''} are simply ``largest $\le X$'' and
``smallest $\ge X$'' -- and ``smallest $\ge X$'' is, value-for-value, the
\emph{value} sitting at the lower-bound index from
Section~\ref{sec:bs-lower-bound}. ``Largest $\le X$'' is its mirror image,
and is found with the symmetric boundary-search template, narrowing toward
the right instead of the left. Whenever a problem uses the words
\emph{``floor''} or \emph{``ceil''} (or, equivalently, \emph{``closest
element not exceeding''} / \emph{``closest element not less than''}),
translate it immediately into a boundary search.
\end{patternbox}

\subsection*{Understanding the problem carefully}

\textbf{Ceil} of $X$: define $P(i) = (\textit{arr}[i] \ge X)$. As in lower
bound, find the smallest index where $P$ holds; the ceil is the value at
that index (or $-1$ if no such index exists, i.e.\ $X$ exceeds every
element).

\textbf{Floor} of $X$: define $Q(i) = (\textit{arr}[i] \le X)$. This
predicate is \texttt{true} for a \emph{prefix} of the array and
\texttt{false} for the remaining suffix -- the mirror image of $P$. We want
the \emph{largest} index where $Q$ holds, so instead of narrowing left
after a success (as in lower/upper bound), we narrow \emph{right} after a
success, since we are now searching for the rightmost true position rather
than the leftmost one.

\begin{ideabox}[Candidate Space and Predicate]
For ceil: candidate space is indices $0,\dots,n-1$ (or ``not found''),
predicate $P(i) = (\textit{arr}[i] \ge X)$, want smallest true index.
For floor: same candidate space, predicate $Q(i) = (\textit{arr}[i] \le X)$,
want \emph{largest} true index.
\end{ideabox}

\subsection*{Why floor needs the mirrored template}

In every boundary search so far (lower bound, upper bound), we wanted the
\emph{leftmost} index satisfying a predicate that becomes true only after
some crossover point, so a successful check let us narrow the window to the
\emph{left}, hunting for an even earlier success. Floor flips this: the
predicate $\textit{arr}[i] \le X$ is true \emph{before} the crossover and
false after, so we want the \emph{rightmost} successful index, meaning a
successful check should narrow the window to the \emph{right}, hunting for
an even later success that is still within the true region. This mirrored
narrowing direction is the only structural change; everything else about
the binary search skeleton is identical.

\subsection*{Algorithm}

\textbf{Ceil:}
\begin{enumerate}
  \item $\textit{lo} \gets 0$, $\textit{hi} \gets n-1$,
        $\textit{ceil} \gets -1$.
  \item While $\textit{lo} \le \textit{hi}$: $\textit{mid} \gets
        \textit{lo}+(\textit{hi}-\textit{lo})/2$. If
        $\textit{arr}[\textit{mid}] \ge X$: $\textit{ceil} \gets
        \textit{arr}[\textit{mid}]$, $\textit{hi} \gets \textit{mid}-1$.
        Else: $\textit{lo} \gets \textit{mid}+1$.
  \item Return \textit{ceil}.
\end{enumerate}

\textbf{Floor:}
\begin{enumerate}
  \item $\textit{lo} \gets 0$, $\textit{hi} \gets n-1$,
        $\textit{floor} \gets -1$.
  \item While $\textit{lo} \le \textit{hi}$: $\textit{mid} \gets
        \textit{lo}+(\textit{hi}-\textit{lo})/2$. If
        $\textit{arr}[\textit{mid}] \le X$: $\textit{floor} \gets
        \textit{arr}[\textit{mid}]$, $\textit{lo} \gets \textit{mid}+1$.
        Else: $\textit{hi} \gets \textit{mid}-1$.
  \item Return \textit{floor}.
\end{enumerate}

\begin{algorithm}[H]
\caption{Floor and Ceil}
\begin{algorithmic}[1]
\Procedure{FloorCeil}{$\textit{arr}, X$}
  \State $\textit{lo} \gets 0$, $\textit{hi} \gets n-1$
  \State $\textit{floor} \gets -1$, $\textit{ceil} \gets -1$
  \While{$\textit{lo} \le \textit{hi}$}
    \State $\textit{mid} \gets \textit{lo} + (\textit{hi}-\textit{lo})/2$
    \If{$\textit{arr}[\textit{mid}] \le X$}
      \State $\textit{floor} \gets \textit{arr}[\textit{mid}]$, $\textit{lo} \gets \textit{mid}+1$
    \Else
      \State $\textit{hi} \gets \textit{mid}-1$
    \EndIf
  \EndWhile
  \State (reset $\textit{lo}, \textit{hi}$ and repeat symmetrically for ceil)
  \State \Return $(\textit{floor}, \textit{ceil})$
\EndProcedure
\end{algorithmic}
\end{algorithm}

\subsection*{Dry run}

\dryrun{Take $\textit{arr} = [3, 4, 4, 7, 8, 10]$, $X = 5$.}

\textbf{Ceil} (smallest $\ge 5$):
\begin{itemize}
  \item $\textit{lo}=0,\textit{hi}=5$: $\textit{mid}=2$,
        $\textit{arr}[2]=4 \ge 5$? No. $\textit{lo}=3$.
  \item $\textit{lo}=3,\textit{hi}=5$: $\textit{mid}=4$,
        $\textit{arr}[4]=8 \ge 5$? Yes. $\textit{ceil}=8$, $\textit{hi}=3$.
  \item $\textit{lo}=3,\textit{hi}=3$: $\textit{mid}=3$,
        $\textit{arr}[3]=7 \ge 5$? Yes. $\textit{ceil}=7$, $\textit{hi}=2$.
  \item $\textit{lo}=3 > \textit{hi}=2$: loop ends. $\textit{ceil}=7$.
\end{itemize}

\textbf{Floor} (largest $\le 5$):
\begin{itemize}
  \item $\textit{lo}=0,\textit{hi}=5$: $\textit{mid}=2$,
        $\textit{arr}[2]=4 \le 5$? Yes. $\textit{floor}=4$, $\textit{lo}=3$.
  \item $\textit{lo}=3,\textit{hi}=5$: $\textit{mid}=4$,
        $\textit{arr}[4]=8 \le 5$? No. $\textit{hi}=3$.
  \item $\textit{lo}=3,\textit{hi}=3$: $\textit{mid}=3$,
        $\textit{arr}[3]=7 \le 5$? No. $\textit{hi}=2$.
  \item $\textit{lo}=3 > \textit{hi}=2$: loop ends. $\textit{floor}=4$.
\end{itemize}

So floor$(5) = 4$ and ceil$(5) = 7$.

\subsection*{C++ implementation}

\begin{lstlisting}
#include <bits/stdc++.h>
using namespace std;

pair<int,int> floorCeil(vector<int>& arr, int X) {
    int n = arr.size();
    int floorVal = -1, ceilVal = -1;

    // Floor: largest element <= X
    int lo = 0, hi = n - 1;
    while (lo <= hi) {
        int mid = lo + (hi - lo) / 2;
        if (arr[mid] <= X) {
            floorVal = arr[mid];
            lo = mid + 1;
        } else {
            hi = mid - 1;
        }
    }

    // Ceil: smallest element >= X
    lo = 0; hi = n - 1;
    while (lo <= hi) {
        int mid = lo + (hi - lo) / 2;
        if (arr[mid] >= X) {
            ceilVal = arr[mid];
            hi = mid - 1;
        } else {
            lo = mid + 1;
        }
    }

    return {floorVal, ceilVal};
}

int main() {
    vector<int> arr = {3, 4, 4, 7, 8, 10};
    auto [floorVal, ceilVal] = floorCeil(arr, 5);
    cout << "Floor: " << floorVal << ", Ceil: " << ceilVal << endl;
    return 0;
}
\end{lstlisting}

\begin{complexitybox}
\textbf{Time Complexity.} Two independent binary searches, each
$O(\log n)$, giving an overall time complexity of
\[
  O(\log n).
\]
\textbf{Space Complexity.} $O(1)$ auxiliary space.
\end{complexitybox}

\begin{takeawaybox}
Floor and ceil are the value-returning siblings of lower and upper bound.
Whichever direction the predicate's ``true'' region lies in (a prefix or a
suffix of the sorted array) tells you which direction to narrow the search
window after a successful check: narrow toward the \emph{far} end of the
true region if you want its near boundary, or keep the running answer
updated while narrowing toward the \emph{near} end if you want the boundary
closest to where you started.
\end{takeawaybox}

\section{Find the First and Last Occurrence of an Element}
\label{sec:bs-first-last}

\problemstatement{We are given a sorted array $\textit{arr}$ of $n$
elements, possibly with duplicates, and a target $X$. We must find the
index of the \textbf{first} occurrence of $X$ and the index of the
\textbf{last} occurrence of $X$. If $X$ is not present, both should be
$-1$.}

\begin{patternbox}
\emph{``First occurrence''} and \emph{``last occurrence''} of a value in a
sorted array are, respectively, exactly the lower bound and
(upper bound $- 1$) of that value, \emph{provided the value is actually
present}. This section is mostly about correctly handling the ``what if it
is not present at all'' case, since lower bound and upper bound alone do
not distinguish between ``$X$ is present, starting here'' and ``$X$ is
absent, and this is just where it would be inserted.''
\end{patternbox}

\subsection*{Understanding the problem carefully}

Recall from Section~\ref{sec:bs-upper-bound} that, when $X$ is present,
\[
  \textit{firstOccurrence} = \textit{lowerBound}(X), \qquad
  \textit{lastOccurrence} = \textit{upperBound}(X) - 1.
\]
The subtlety is verifying presence. The lower bound of $X$ gives the first
index whose value is $\ge X$; this is the same as the first occurrence of
$X$ \emph{only if} the value actually stored at that index equals $X$ (and
the index is within bounds). If $\textit{lowerBound}(X) = n$, or if
$\textit{arr}[\textit{lowerBound}(X)] \ne X$, then $X$ is simply not in the
array at all.

\begin{ideabox}[Approach]
Compute $\textit{lb} \gets \textit{lowerBound}(X)$. If $\textit{lb} = n$ or
$\textit{arr}[\textit{lb}] \ne X$, return $(-1,-1)$. Otherwise, the first
occurrence is $\textit{lb}$, and the last occurrence is
$\textit{upperBound}(X) - 1$.
\end{ideabox}

\subsection*{Why checking presence only once is sufficient}

Because the array is sorted, every occurrence of $X$ forms a single
contiguous block. If the very first candidate position (the lower bound)
does not actually hold $X$, no later position can hold $X$ either -- the
lower bound is, by construction, the leftmost position any element $\ge X$
could occupy, so if even that position fails to be exactly $X$, $X$ does
not appear anywhere in the array. This means a single equality check after
computing the lower bound is enough to settle presence, with no need for
an additional independent search.

\subsection*{Algorithm}

\begin{enumerate}
  \item $\textit{lb} \gets \textit{lowerBound}(\textit{arr}, X)$ (using
        Algorithm~\ref{sec:bs-lower-bound}).
  \item If $\textit{lb} = n$ or $\textit{arr}[\textit{lb}] \ne X$: return
        $(-1, -1)$.
  \item Otherwise: $\textit{ub} \gets \textit{upperBound}(\textit{arr}, X)$
        (using Algorithm~\ref{sec:bs-upper-bound}). Return
        $(\textit{lb}, \textit{ub} - 1)$.
\end{enumerate}

\begin{algorithm}[H]
\caption{First and Last Occurrence}
\begin{algorithmic}[1]
\Procedure{FirstAndLast}{$\textit{arr}, X$}
  \State $\textit{lb} \gets \Call{LowerBound}{\textit{arr}, X}$
  \If{$\textit{lb} = n$ \textbf{or} $\textit{arr}[\textit{lb}] \ne X$}
    \State \Return $(-1, -1)$
  \EndIf
  \State $\textit{ub} \gets \Call{UpperBound}{\textit{arr}, X}$
  \State \Return $(\textit{lb}, \textit{ub} - 1)$
\EndProcedure
\end{algorithmic}
\end{algorithm}

\subsection*{Dry run}

\dryrun{Take $\textit{arr} = [5, 7, 7, 8, 8, 10]$, $X = 8$.}

$\textit{lowerBound}(8)$: scanning, the first index with value $\ge 8$ is
index $3$ ($\textit{arr}[3]=8$). Since $\textit{arr}[3]=8=X$, $X$ is
present. $\textit{upperBound}(8)$: the first index with value $>8$ is
index $5$ ($\textit{arr}[5]=10$). So
$\textit{lastOccurrence} = 5 - 1 = 4$.

The answer is (first $=3$, last $=4$).

Now take $X = 6$ on the same array. $\textit{lowerBound}(6) = 1$
($\textit{arr}[1]=7 \ge 6$ is the first such index). But
$\textit{arr}[1] = 7 \ne 6$, so $6$ is not present, and the answer is
$(-1, -1)$.

\subsection*{C++ implementation}

\begin{lstlisting}
#include <bits/stdc++.h>
using namespace std;

int lowerBound(vector<int>& arr, int X) {
    int n = arr.size();
    int lo = 0, hi = n - 1, ans = n;
    while (lo <= hi) {
        int mid = lo + (hi - lo) / 2;
        if (arr[mid] >= X) { ans = mid; hi = mid - 1; }
        else lo = mid + 1;
    }
    return ans;
}

int upperBound(vector<int>& arr, int X) {
    int n = arr.size();
    int lo = 0, hi = n - 1, ans = n;
    while (lo <= hi) {
        int mid = lo + (hi - lo) / 2;
        if (arr[mid] > X) { ans = mid; hi = mid - 1; }
        else lo = mid + 1;
    }
    return ans;
}

pair<int,int> firstAndLast(vector<int>& arr, int X) {
    int n = arr.size();
    int lb = lowerBound(arr, X);
    if (lb == n || arr[lb] != X) return {-1, -1};

    int ub = upperBound(arr, X);
    return {lb, ub - 1};
}

int main() {
    vector<int> arr = {5, 7, 7, 8, 8, 10};
    auto [first, last] = firstAndLast(arr, 8);
    cout << "First: " << first << ", Last: " << last << endl;
    return 0;
}
\end{lstlisting}

\begin{complexitybox}
\textbf{Time Complexity.} Two binary searches (lower bound and, if needed,
upper bound), each $O(\log n)$, giving an overall time complexity of
\[
  O(\log n).
\]
\textbf{Space Complexity.} $O(1)$ auxiliary space.
\end{complexitybox}

\begin{takeawaybox}
Lower bound and upper bound, used together with a single presence check,
solve any ``range of occurrences'' question about a sorted array. There is
rarely a need to write a bespoke binary search for ``first occurrence'' or
``last occurrence'' from scratch -- compose them from the lower/upper bound
primitives instead.
\end{takeawaybox}

\section{Count Occurrences of an Element in a Sorted Array}
\label{sec:bs-count-occurrences}

\problemstatement{We are given a sorted array $\textit{arr}$ of $n$
elements, possibly with duplicates, and a target $X$. We must count how
many times $X$ occurs in $\textit{arr}$.}

\begin{patternbox}
\emph{``Count occurrences of $X$ in a sorted array''} is immediately
solvable once you recall that all occurrences of $X$ occupy a single
contiguous block (Section~\ref{sec:bs-first-last}), bounded by the lower
bound and upper bound of $X$. The keyword to internalize: any time a
sorted-array problem asks ``how many,'' check first whether the answer is
simply the \emph{width} of a contiguous range you can find with two
boundary searches, rather than something requiring an actual scan.
\end{patternbox}

\subsection*{Understanding the problem carefully}

The lower bound of $X$ is the first index where $X$ could occur; the upper
bound of $X$ is the first index strictly past where $X$ can occur. Every
index in between holds a value equal to $X$ (since the array is sorted and
both boundaries are tight), and no index outside that range does. So the
count of occurrences is simply the number of indices in
$[\textit{lowerBound}(X), \textit{upperBound}(X))$, i.e.
\[
  \textit{count} = \textit{upperBound}(X) - \textit{lowerBound}(X).
\]
This formula is correct even when $X$ is entirely absent: in that case
$\textit{lowerBound}(X) = \textit{upperBound}(X)$ (both point to the same
insertion position), giving a count of $0$ automatically, with no special
case needed.

\begin{ideabox}[Why No Presence Check Is Needed Here]
Unlike Section~\ref{sec:bs-first-last}, where we needed an explicit
presence check before reporting indices (an absent element has no valid
index to report), here the quantity we want is a \emph{count}, and the
subtraction $\textit{upperBound}(X) - \textit{lowerBound}(X)$ naturally
evaluates to $0$ when $X$ is absent, without requiring a separate branch.
\end{ideabox}

\subsection*{Algorithm}

\begin{enumerate}
  \item $\textit{lb} \gets \textit{lowerBound}(\textit{arr}, X)$.
  \item $\textit{ub} \gets \textit{upperBound}(\textit{arr}, X)$.
  \item Return $\textit{ub} - \textit{lb}$.
\end{enumerate}

\begin{algorithm}[H]
\caption{Count Occurrences}
\begin{algorithmic}[1]
\Procedure{CountOccurrences}{$\textit{arr}, X$}
  \State $\textit{lb} \gets \Call{LowerBound}{\textit{arr}, X}$
  \State $\textit{ub} \gets \Call{UpperBound}{\textit{arr}, X}$
  \State \Return $\textit{ub} - \textit{lb}$
\EndProcedure
\end{algorithmic}
\end{algorithm}

\subsection*{Dry run}

\dryrun{Take $\textit{arr} = [2, 4, 4, 4, 4, 8, 10]$, $X = 4$.}

$\textit{lowerBound}(4)$: first index with value $\ge 4$ is index $1$.
$\textit{upperBound}(4)$: first index with value $> 4$ is index $5$
($\textit{arr}[5]=8$). So $\textit{count} = 5 - 1 = 4$, correctly matching
the four occurrences of $4$ at indices $1,2,3,4$.

Now take $X = 6$ on the same array (absent). Both
$\textit{lowerBound}(6)$ and $\textit{upperBound}(6)$ equal $5$ (the first
index with value $\ge 6$, and also the first index with value $>6$, is
index $5$, where $\textit{arr}[5]=8$). So $\textit{count} = 5-5=0$,
correctly reporting absence.

\subsection*{C++ implementation}

\begin{lstlisting}
#include <bits/stdc++.h>
using namespace std;

int lowerBound(vector<int>& arr, int X) {
    int n = arr.size();
    int lo = 0, hi = n - 1, ans = n;
    while (lo <= hi) {
        int mid = lo + (hi - lo) / 2;
        if (arr[mid] >= X) { ans = mid; hi = mid - 1; }
        else lo = mid + 1;
    }
    return ans;
}

int upperBound(vector<int>& arr, int X) {
    int n = arr.size();
    int lo = 0, hi = n - 1, ans = n;
    while (lo <= hi) {
        int mid = lo + (hi - lo) / 2;
        if (arr[mid] > X) { ans = mid; hi = mid - 1; }
        else lo = mid + 1;
    }
    return ans;
}

int countOccurrences(vector<int>& arr, int X) {
    return upperBound(arr, X) - lowerBound(arr, X);
}

int main() {
    vector<int> arr = {2, 4, 4, 4, 4, 8, 10};
    cout << "Occurrences of 4: " << countOccurrences(arr, 4) << endl;
    cout << "Occurrences of 6: " << countOccurrences(arr, 6) << endl;
    return 0;
}
\end{lstlisting}

\begin{complexitybox}
\textbf{Time Complexity.} Two binary searches, each $O(\log n)$, giving
\[
  O(\log n),
\]
a dramatic improvement over the $O(n)$ linear scan that counting
occurrences would otherwise require.
\textbf{Space Complexity.} $O(1)$ auxiliary space.
\end{complexitybox}

\begin{takeawaybox}
``How many elements satisfy property $P$'' on a sorted array, where $P$
itself is monotonic (true on a contiguous range), almost always reduces to
``find the two boundaries of the true range and subtract.'' This is a
strictly more general restatement of the lower/upper-bound subtraction
trick, worth recognizing beyond just counting exact value matches.
\end{takeawaybox}

\section{Search in a Rotated Sorted Array I}
\label{sec:bs-rotated-1}

\problemstatement{An array that was originally sorted in increasing order
has been \textbf{rotated} at some unknown pivot -- for example,
$[0,1,2,4,5,6,7]$ rotated at index $3$ becomes $[4,5,6,7,0,1,2]$. All
elements are distinct. Given such a rotated array $\textit{arr}$ and a
target $X$, find the index of $X$, or return $-1$ if it is not present.
We must still achieve $O(\log n)$ time.}

\begin{patternbox}
\emph{``Rotated sorted array''} is a clear signal that plain binary search
(Section~\ref{sec:bs-find-x}) cannot be applied directly, because the array
as a whole is no longer globally sorted -- comparing
$\textit{arr}[\textit{mid}]$ to $X$ no longer tells us which half to
discard. But the keyword \emph{``sorted, then rotated''} carries a crucial
piece of remaining structure: although the \emph{whole} array is not
sorted, \textbf{at least one of the two halves split at any midpoint always
is}. Whenever a problem mentions rotation, immediately look for this
``one half is always sorted'' structure, since it is what restores binary
search's monotonic-discard guarantee.
\end{patternbox}

\subsection*{Understanding the problem carefully}

Pick any midpoint $\textit{mid}$ in a rotated sorted array. The two halves
$[\textit{lo}, \textit{mid}]$ and $[\textit{mid}, \textit{hi}]$ cannot both
contain the rotation pivot (the single place where a larger element is
immediately followed by a smaller one) -- the pivot lies in at most one of
the two halves, so the other half is guaranteed to be a genuinely sorted,
unrotated run of increasing values. We can detect which half is sorted with
one comparison: if $\textit{arr}[\textit{lo}] \le
\textit{arr}[\textit{mid}]$, the left half is sorted; otherwise, the right
half must be sorted.

\begin{ideabox}[Greedy Decision at Each Step]
At each step, first determine which half is sorted using
$\textit{arr}[\textit{lo}] \le \textit{arr}[\textit{mid}]$. Then, check
whether $X$ could possibly lie within that sorted half's value range (using
ordinary comparisons, since that half really is sorted). If $X$'s value
falls within the sorted half's range, discard the other half and recurse
into the sorted half; otherwise, $X$ (if present at all) must be in the
\emph{other} half, so discard the sorted half and recurse into the
unsorted-looking half instead -- which, on the next iteration, will itself
be split into one sorted and one possibly-unsorted piece, maintaining the
invariant.
\end{ideabox}

\subsection*{Why this discard rule is correct}

The key invariant maintained throughout is that the current search window
$[\textit{lo}, \textit{hi}]$ is itself a rotated sorted array (a
sub-array of the original, which inherits the ``rotated sorted'' structure
since removing elements from one end of a rotated sorted run leaves a
smaller rotated sorted run). At every step we identify a genuinely sorted
half by comparing its endpoints directly -- a valid comparison because that
half truly has no rotation inside it. If $X$ lies within that sorted half's
numeric range (i.e.\ between its smallest and largest values), it can only
be found there, since the other half's values lie entirely outside this
range (otherwise the sortedness of the identified half, combined with the
single-rotation-point structure, would be violated). So discarding the
half that provably cannot contain $X$ is always safe, exactly mirroring the
discard argument from ordinary binary search, just with one extra
case-split to first identify which half is safe to reason about directly.

\subsection*{Algorithm}

\begin{enumerate}
  \item Initialize $\textit{lo} \gets 0$, $\textit{hi} \gets n-1$.
  \item While $\textit{lo} \le \textit{hi}$:
  \begin{enumerate}
    \item $\textit{mid} \gets \textit{lo}+(\textit{hi}-\textit{lo})/2$.
    \item If $\textit{arr}[\textit{mid}] = X$: return \textit{mid}.
    \item If $\textit{arr}[\textit{lo}] \le \textit{arr}[\textit{mid}]$
          (left half is sorted):
    \begin{itemize}
      \item If $\textit{arr}[\textit{lo}] \le X < \textit{arr}[\textit{mid}]$:
            search left, $\textit{hi} \gets \textit{mid}-1$.
      \item Else: search right, $\textit{lo} \gets \textit{mid}+1$.
    \end{itemize}
    \item Otherwise (right half is sorted):
    \begin{itemize}
      \item If $\textit{arr}[\textit{mid}] < X \le \textit{arr}[\textit{hi}]$:
            search right, $\textit{lo} \gets \textit{mid}+1$.
      \item Else: search left, $\textit{hi} \gets \textit{mid}-1$.
    \end{itemize}
  \end{enumerate}
  \item If the loop ends without returning, return $-1$.
\end{enumerate}

\begin{algorithm}[H]
\caption{Search in Rotated Sorted Array (Distinct Elements)}
\begin{algorithmic}[1]
\Procedure{SearchRotated}{$\textit{arr}, X$}
  \State $\textit{lo} \gets 0$, $\textit{hi} \gets n-1$
  \While{$\textit{lo} \le \textit{hi}$}
    \State $\textit{mid} \gets \textit{lo}+(\textit{hi}-\textit{lo})/2$
    \If{$\textit{arr}[\textit{mid}] = X$} \Return \textit{mid} \EndIf
    \If{$\textit{arr}[\textit{lo}] \le \textit{arr}[\textit{mid}]$}
      \If{$\textit{arr}[\textit{lo}] \le X$ \textbf{and} $X < \textit{arr}[\textit{mid}]$}
        \State $\textit{hi} \gets \textit{mid} - 1$
      \Else
        \State $\textit{lo} \gets \textit{mid} + 1$
      \EndIf
    \Else
      \If{$\textit{arr}[\textit{mid}] < X$ \textbf{and} $X \le \textit{arr}[\textit{hi}]$}
        \State $\textit{lo} \gets \textit{mid} + 1$
      \Else
        \State $\textit{hi} \gets \textit{mid} - 1$
      \EndIf
    \EndIf
  \EndWhile
  \State \Return $-1$
\EndProcedure
\end{algorithmic}
\end{algorithm}

\subsection*{Dry run}

\dryrun{Take $\textit{arr} = [4,5,6,7,0,1,2]$, $X = 0$.}

\begin{itemize}
  \item $\textit{lo}=0,\textit{hi}=6$: $\textit{mid}=3$,
        $\textit{arr}[3]=7 \ne 0$. Is $\textit{arr}[0]=4 \le
        \textit{arr}[3]=7$? Yes, left half sorted. Is
        $4 \le 0 < 7$? No ($0 < 4$). So search right:
        $\textit{lo}=4$.
  \item $\textit{lo}=4,\textit{hi}=6$: $\textit{mid}=5$,
        $\textit{arr}[5]=1 \ne 0$. Is $\textit{arr}[4]=0 \le
        \textit{arr}[5]=1$? Yes, left half sorted. Is
        $0 \le 0 < 1$? Yes. Search left: $\textit{hi}=4$.
  \item $\textit{lo}=4,\textit{hi}=4$: $\textit{mid}=4$,
        $\textit{arr}[4]=0 = 0$. Return $4$.
\end{itemize}

The answer is index $4$.

\subsection*{C++ implementation}

\begin{lstlisting}
#include <bits/stdc++.h>
using namespace std;

int searchRotated(vector<int>& arr, int X) {
    int n = arr.size();
    int lo = 0, hi = n - 1;

    while (lo <= hi) {
        int mid = lo + (hi - lo) / 2;
        if (arr[mid] == X) return mid;

        if (arr[lo] <= arr[mid]) { // left half sorted
            if (arr[lo] <= X && X < arr[mid]) {
                hi = mid - 1;
            } else {
                lo = mid + 1;
            }
        } else { // right half sorted
            if (arr[mid] < X && X <= arr[hi]) {
                lo = mid + 1;
            } else {
                hi = mid - 1;
            }
        }
    }
    return -1;
}

int main() {
    vector<int> arr = {4, 5, 6, 7, 0, 1, 2};
    cout << "Index of 0: " << searchRotated(arr, 0) << endl;
    return 0;
}
\end{lstlisting}

\begin{complexitybox}
\textbf{Time Complexity.} Each iteration discards at least one full half
of the current window, exactly as in ordinary binary search, giving
\[
  O(\log n).
\]
\textbf{Space Complexity.} $O(1)$ auxiliary space.
\end{complexitybox}

\begin{takeawaybox}
When an array's global sortedness is broken in one specific, structured way
(here, a single rotation point), look for a local property that still
holds at every step of binary search -- here, ``at least one of the two
halves around any midpoint is genuinely sorted.'' Identifying and exploiting
that local invariant is what lets binary search survive structural
violations of global sortedness.
\end{takeawaybox}

\section{Search in a Rotated Sorted Array II}
\label{sec:bs-rotated-2}

\problemstatement{Same setting as Section~\ref{sec:bs-rotated-1} -- a
sorted array rotated at an unknown pivot -- except now the array
\textbf{may contain duplicate values}. Given the rotated array
$\textit{arr}$ and a target $X$, determine whether $X$ is present
(\texttt{true}/\texttt{false}).}

\begin{patternbox}
The keyword change from Section~\ref{sec:bs-rotated-1} is just one word:
\emph{``may contain duplicates.''} This single word breaks the comparison
$\textit{arr}[\textit{lo}] \le \textit{arr}[\textit{mid}]$ that we relied on
to decide which half is sorted -- with duplicates, this comparison can be
\texttt{true} even when $\textit{lo}$, $\textit{mid}$, and the rotation
pivot are tangled together in a way that makes \emph{neither} half safely
identifiable as sorted. Whenever a rotated-array problem mentions
duplicates, expect to need an escape hatch for this ambiguous case, and
expect the worst-case time complexity to degrade as a result.
\end{patternbox}

\subsection*{Understanding the problem carefully}

Consider $\textit{arr} = [3,3,3,1,2,3]$ with $\textit{lo}=0$ pointing to
value $3$, and suppose $\textit{mid}$ also happens to land on a $3$ (say
$\textit{mid}=2$, $\textit{arr}[2]=3$). Then
$\textit{arr}[\textit{lo}] = \textit{arr}[\textit{mid}] = 3$, so the
inequality $\textit{arr}[\textit{lo}] \le \textit{arr}[\textit{mid}]$ holds,
which the Section~\ref{sec:bs-rotated-1} logic would interpret as ``the
left half is sorted.'' But the rotation pivot (where $1$ follows $3$)
actually sits \emph{inside} this very range, so the left half is not
genuinely a clean sorted run with respect to deciding where $X$ could be --
the equal values at the boundary hide the pivot's location. With distinct
elements this ambiguity cannot arise (any two unequal boundary values
immediately tell us which side the pivot is on), but duplicates can make
$\textit{arr}[\textit{lo}] = \textit{arr}[\textit{mid}]$ true regardless of
where the pivot actually is.

\begin{ideabox}[{Greedy Decision, with an Escape Hatch}]
Use the exact same logic as Section~\ref{sec:bs-rotated-1}, \emph{except}:
whenever $\textit{arr}[\textit{lo}] = \textit{arr}[\textit{mid}]$ \emph{and}
$\textit{arr}[\textit{mid}] = \textit{arr}[\textit{hi}]$ (or, more simply,
whenever the usual ``which half is sorted'' check is ambiguous because of
equal boundary values), we cannot safely discard either half based on
comparisons alone. In that situation, shrink the window by one element from
each side ($\textit{lo} \gets \textit{lo}+1$, $\textit{hi} \gets
\textit{hi}-1$) and try again -- this never discards a position that could
uniquely hold $X$ in a way that loses correctness, but it does cost us the
usual halving speed in the worst case.
\end{ideabox}

\subsection*{Why shrinking by one is the right (and only) safe fallback}

When $\textit{arr}[\textit{lo}] = \textit{arr}[\textit{mid}]$, we genuinely
cannot tell, from this comparison alone, which half contains the rotation
pivot -- both possibilities are consistent with the observed equality. Since
we cannot identify a provably safe half to discard, we are not entitled to
discard a full half at all without risking skipping over $X$. The one
adjustment we \emph{can} make safely is to shrink the boundary by a single
element: if $\textit{arr}[\textit{lo}]$ duplicates the value at
$\textit{mid}$, removing just $\textit{lo}$ from consideration cannot lose
$X$, because if $X = \textit{arr}[\textit{lo}]$, the duplicate value
persists at $\textit{mid}$ and remains discoverable; if $X \ne
\textit{arr}[\textit{lo}]$, removing $\textit{lo}$ loses nothing relevant.
This is why the worst case degrades to $O(n)$: an adversarial array of
nearly all-equal elements (e.g. $[2,2,2,2,2,2,2,1,2]$) can force this
one-at-a-time shrinking on almost every step.

\subsection*{Algorithm}

\begin{enumerate}
  \item Initialize $\textit{lo} \gets 0$, $\textit{hi} \gets n-1$.
  \item While $\textit{lo} \le \textit{hi}$:
  \begin{enumerate}
    \item $\textit{mid} \gets \textit{lo}+(\textit{hi}-\textit{lo})/2$.
    \item If $\textit{arr}[\textit{mid}] = X$: return \texttt{true}.
    \item If $\textit{arr}[\textit{lo}] = \textit{arr}[\textit{mid}]$ and
          $\textit{arr}[\textit{mid}] = \textit{arr}[\textit{hi}]$:
          $\textit{lo} \gets \textit{lo}+1$, $\textit{hi} \gets
          \textit{hi}-1$ (ambiguous; shrink by one from each side).
    \item Else if $\textit{arr}[\textit{lo}] \le \textit{arr}[\textit{mid}]$
          (left half safely sorted): apply the same range check as
          Section~\ref{sec:bs-rotated-1}.
    \item Else (right half safely sorted): apply the same range check as
          Section~\ref{sec:bs-rotated-1}.
  \end{enumerate}
  \item If the loop ends without returning, return \texttt{false}.
\end{enumerate}

\begin{algorithm}[H]
\caption{Search in Rotated Sorted Array (With Duplicates)}
\begin{algorithmic}[1]
\Procedure{SearchRotatedDup}{$\textit{arr}, X$}
  \State $\textit{lo} \gets 0$, $\textit{hi} \gets n-1$
  \While{$\textit{lo} \le \textit{hi}$}
    \State $\textit{mid} \gets \textit{lo}+(\textit{hi}-\textit{lo})/2$
    \If{$\textit{arr}[\textit{mid}] = X$} \Return \textbf{true} \EndIf
    \If{$\textit{arr}[\textit{lo}] = \textit{arr}[\textit{mid}]$ \textbf{and} $\textit{arr}[\textit{mid}] = \textit{arr}[\textit{hi}]$}
      \State $\textit{lo} \gets \textit{lo}+1$, $\textit{hi} \gets \textit{hi}-1$
    \ElsIf{$\textit{arr}[\textit{lo}] \le \textit{arr}[\textit{mid}]$}
      \If{$\textit{arr}[\textit{lo}] \le X$ \textbf{and} $X < \textit{arr}[\textit{mid}]$}
        \State $\textit{hi} \gets \textit{mid} - 1$
      \Else
        \State $\textit{lo} \gets \textit{mid} + 1$
      \EndIf
    \Else
      \If{$\textit{arr}[\textit{mid}] < X$ \textbf{and} $X \le \textit{arr}[\textit{hi}]$}
        \State $\textit{lo} \gets \textit{mid} + 1$
      \Else
        \State $\textit{hi} \gets \textit{mid} - 1$
      \EndIf
    \EndIf
  \EndWhile
  \State \Return \textbf{false}
\EndProcedure
\end{algorithmic}
\end{algorithm}

\subsection*{Dry run}

\dryrun{Take $\textit{arr} = [3,1,2,3,3,3,3]$, $X = 1$.}

\begin{itemize}
  \item $\textit{lo}=0,\textit{hi}=6$: $\textit{mid}=3$,
        $\textit{arr}[3]=3 \ne 1$. Is $\textit{arr}[0]=3=\textit{arr}[3]=3$
        and $\textit{arr}[3]=3=\textit{arr}[6]=3$? Yes -- ambiguous. Shrink:
        $\textit{lo}=1$, $\textit{hi}=5$.
  \item $\textit{lo}=1,\textit{hi}=5$: $\textit{mid}=3$,
        $\textit{arr}[3]=3 \ne 1$. Is $\textit{arr}[1]=1=\textit{arr}[3]=3$?
        No. Is $\textit{arr}[1]=1 \le \textit{arr}[3]=3$? Yes, left half
        sorted. Is $1 \le 1 < 3$? Yes. Search left: $\textit{hi}=2$.
  \item $\textit{lo}=1,\textit{hi}=2$: $\textit{mid}=1$,
        $\textit{arr}[1]=1=1$. Return \texttt{true}.
\end{itemize}

The answer is \texttt{true}.

\subsection*{C++ implementation}

\begin{lstlisting}
#include <bits/stdc++.h>
using namespace std;

bool searchRotatedDup(vector<int>& arr, int X) {
    int n = arr.size();
    int lo = 0, hi = n - 1;

    while (lo <= hi) {
        int mid = lo + (hi - lo) / 2;
        if (arr[mid] == X) return true;

        if (arr[lo] == arr[mid] && arr[mid] == arr[hi]) {
            lo++;
            hi--;
        } else if (arr[lo] <= arr[mid]) {
            if (arr[lo] <= X && X < arr[mid]) {
                hi = mid - 1;
            } else {
                lo = mid + 1;
            }
        } else {
            if (arr[mid] < X && X <= arr[hi]) {
                lo = mid + 1;
            } else {
                hi = mid - 1;
            }
        }
    }
    return false;
}

int main() {
    vector<int> arr = {3, 1, 2, 3, 3, 3, 3};
    cout << (searchRotatedDup(arr, 1) ? "true" : "false") << endl;
    return 0;
}
\end{lstlisting}

\begin{complexitybox}
\textbf{Time Complexity.} In the best/average case, the same halving as
Section~\ref{sec:bs-rotated-1} applies, giving $O(\log n)$. But in the
worst case -- an array of almost entirely duplicate values, e.g.
$[2,2,2,2,2,2,1,2]$ -- the ambiguous case forces shrinking by one element
at a time, giving a worst-case time complexity of
\[
  O(n).
\]
\textbf{Space Complexity.} $O(1)$ auxiliary space.
\end{complexitybox}

\begin{takeawaybox}
Duplicates are binary search's natural enemy: they tend to destroy the
clean monotonicity or structural invariants that comparisons rely on,
forcing a fallback to linear-time shrinking in the worst case. Whenever a
sheet problem adds ``with duplicates'' to a previously clean binary search
problem, expect to need an explicit ambiguous-case branch, and expect to
state a degraded worst-case complexity honestly rather than claiming
$O(\log n)$ that the duplicates actually prevent.
\end{takeawaybox}

\section{Find Minimum in a Rotated Sorted Array}
\label{sec:bs-find-min-rotated}

\problemstatement{We are given a sorted array (with all distinct elements)
that has been rotated at some unknown pivot. We must find the
\textbf{minimum element} of the array in $O(\log n)$ time.}

\begin{patternbox}
\emph{``Minimum element of a rotated sorted array''} is really asking us to
\textbf{find the rotation pivot itself} -- the minimum element is exactly
the element sitting at the pivot (the one place where the rotated array's
increasing run restarts). Once you see ``rotated sorted array'' combined
with ``find the minimum/find the rotation point,'' reach for the same
``one half is always sorted'' invariant used in
Sections~\ref{sec:bs-rotated-1}--\ref{sec:bs-rotated-2}, but this time the
goal is to home in on the pivot's location directly, rather than to locate
a separately given target value.
\end{patternbox}

\subsection*{Understanding the problem carefully}

In a rotated sorted array, exactly one of the following is true at any
candidate window $[\textit{lo}, \textit{hi}]$: either the window itself is
already fully sorted (no rotation point inside it), in which case its
minimum is simply $\textit{arr}[\textit{lo}]$; or the rotation point lies
strictly inside the window, in which case we can use the same
left-half/right-half sortedness check from Section~\ref{sec:bs-rotated-1}
to decide which half the rotation point falls into, and discard the half
that is provably free of it.

\begin{ideabox}[Greedy Choice]
If $\textit{arr}[\textit{lo}] \le \textit{arr}[\textit{hi}]$, the current
window is already fully sorted (no rotation inside it), so its minimum is
$\textit{arr}[\textit{lo}]$ -- record it and stop. Otherwise, compare
$\textit{arr}[\textit{mid}]$ to $\textit{arr}[\textit{lo}]$: if
$\textit{arr}[\textit{mid}] \ge \textit{arr}[\textit{lo}]$, the left half
$[\textit{lo}, \textit{mid}]$ is sorted (so its minimum is
$\textit{arr}[\textit{lo}]$, not better than what is achievable on the
right, since the right half must contain the actual rotation point and
hence a smaller value); discard the left half by moving $\textit{lo}
\gets \textit{mid}+1$. Otherwise, the rotation point lies in
$[\textit{lo}, \textit{mid}]$ itself, so we keep $\textit{mid}$ as a
candidate minimum and search left: $\textit{hi} \gets \textit{mid}$ (not
$\textit{mid}-1$, since $\textit{mid}$ itself might be the true minimum).
\end{ideabox}

\subsection*{Why discarding the sorted half is always correct here}

If $\textit{arr}[\textit{mid}] \ge \textit{arr}[\textit{lo}]$, the segment
$[\textit{lo}, \textit{mid}]$ has no internal rotation (it is increasing
throughout), so its smallest value is $\textit{arr}[\textit{lo}]$, which is
\emph{not smaller} than the global minimum unless the global minimum
happens to equal $\textit{arr}[\textit{lo}]$ itself -- but that can only
happen if the whole array (or at least everything from \textit{lo} onward
that matters) is sorted with no rotation remaining, a case that is only
reachable by continuing to shrink toward the actual pivot, which must then
lie at or to the right of \textit{mid}. So it is always safe to move
$\textit{lo}$ past this sorted segment; the true minimum, if not already
found, must lie strictly within $(\textit{mid}, \textit{hi}]$. Conversely,
if $\textit{arr}[\textit{mid}] < \textit{arr}[\textit{lo}]$, the array has
``wrapped around'' somewhere in $[\textit{lo}, \textit{mid}]$, so the
rotation pivot -- and hence the true minimum -- is guaranteed to lie within
that left segment, including possibly at $\textit{mid}$ itself, which is
why we narrow to $\textit{hi} \gets \textit{mid}$ rather than excluding it.

\subsection*{Algorithm}

\begin{enumerate}
  \item Initialize $\textit{lo} \gets 0$, $\textit{hi} \gets n-1$.
  \item While $\textit{lo} < \textit{hi}$:
  \begin{enumerate}
    \item If $\textit{arr}[\textit{lo}] \le \textit{arr}[\textit{hi}]$:
          the window is sorted; break (the minimum is
          $\textit{arr}[\textit{lo}]$).
    \item $\textit{mid} \gets \textit{lo}+(\textit{hi}-\textit{lo})/2$.
    \item If $\textit{arr}[\textit{mid}] \ge \textit{arr}[\textit{lo}]$:
          $\textit{lo} \gets \textit{mid}+1$.
    \item Else: $\textit{hi} \gets \textit{mid}$.
  \end{enumerate}
  \item Return $\textit{arr}[\textit{lo}]$.
\end{enumerate}

\begin{algorithm}[H]
\caption{Find Minimum in Rotated Sorted Array}
\begin{algorithmic}[1]
\Procedure{FindMin}{\textit{arr}}
  \State $\textit{lo} \gets 0$, $\textit{hi} \gets n-1$
  \While{$\textit{lo} < \textit{hi}$}
    \If{$\textit{arr}[\textit{lo}] \le \textit{arr}[\textit{hi}]$}
      \State \textbf{break}
    \EndIf
    \State $\textit{mid} \gets \textit{lo}+(\textit{hi}-\textit{lo})/2$
    \If{$\textit{arr}[\textit{mid}] \ge \textit{arr}[\textit{lo}]$}
      \State $\textit{lo} \gets \textit{mid} + 1$
    \Else
      \State $\textit{hi} \gets \textit{mid}$
    \EndIf
  \EndWhile
  \State \Return $\textit{arr}[\textit{lo}]$
\EndProcedure
\end{algorithmic}
\end{algorithm}

\subsection*{Dry run}

\dryrun{Take $\textit{arr} = [4,5,6,7,0,1,2]$.}

\begin{itemize}
  \item $\textit{lo}=0,\textit{hi}=6$: $\textit{arr}[0]=4 \le
        \textit{arr}[6]=2$? No. $\textit{mid}=3$,
        $\textit{arr}[3]=7 \ge \textit{arr}[0]=4$? Yes. $\textit{lo}=4$.
  \item $\textit{lo}=4,\textit{hi}=6$: $\textit{arr}[4]=0 \le
        \textit{arr}[6]=2$? Yes -- window is sorted. Break.
\end{itemize}

The minimum is $\textit{arr}[4] = 0$.

\subsection*{C++ implementation}

\begin{lstlisting}
#include <bits/stdc++.h>
using namespace std;

int findMin(vector<int>& arr) {
    int n = arr.size();
    int lo = 0, hi = n - 1;

    while (lo < hi) {
        if (arr[lo] <= arr[hi]) break; // already sorted

        int mid = lo + (hi - lo) / 2;
        if (arr[mid] >= arr[lo]) {
            lo = mid + 1;
        } else {
            hi = mid;
        }
    }
    return arr[lo];
}

int main() {
    vector<int> arr = {4, 5, 6, 7, 0, 1, 2};
    cout << "Minimum: " << findMin(arr) << endl;
    return 0;
}
\end{lstlisting}

\begin{complexitybox}
\textbf{Time Complexity.} Each iteration halves the search window (or
terminates early once the window is detected to be sorted), giving
\[
  O(\log n).
\]
\textbf{Space Complexity.} $O(1)$ auxiliary space.
\end{complexitybox}

\begin{takeawaybox}
Finding the rotation pivot and searching for a target in a rotated array
(Section~\ref{sec:bs-rotated-1}) rely on the exact same underlying
invariant -- one half around any midpoint is always genuinely sorted. The
pivot itself is simply the boundary that this repeated halving converges
to when there is no separate target value to compare against, only the
structure of the rotation itself.
\end{takeawaybox}

\section{Find How Many Times an Array Has Been Rotated}
\label{sec:bs-how-many-rotations}

\problemstatement{Given a sorted array that has been rotated some number of
times (all elements distinct), determine \textbf{how many times} it has
been rotated -- equivalently, find the number of left rotations applied to
the original sorted array.}

\begin{patternbox}
This problem is solved by noticing that it is asking for almost the same
quantity as Section~\ref{sec:bs-find-min-rotated}, phrased differently: the
number of left rotations applied to a sorted array equals exactly the
\textbf{index} of the minimum element. If the original sorted array is
rotated left by $k$ positions, the element that used to be at index $0$
(the global minimum) ends up sitting at index $k$. Whenever a problem
about a rotated array asks for ``how many rotations,'' translate it
immediately into ``find the index of the minimum.''
\end{patternbox}

\subsection*{Understanding the problem carefully}

Picture the original sorted array $[v_0 < v_1 < \cdots < v_{n-1}]$. Rotating
it left by $k$ positions produces
$[v_k, v_{k+1}, \dots, v_{n-1}, v_0, v_1, \dots, v_{k-1}]$. The smallest
value $v_0$ now sits at position $k$ in the rotated array. So finding the
number of rotations is exactly finding the position of the minimum element
-- the same quantity computed in Section~\ref{sec:bs-find-min-rotated},
except we now want the \emph{index} rather than the \emph{value}.

\begin{ideabox}[Reuse, Don't Re-derive]
Apply the exact same binary search as
Section~\ref{sec:bs-find-min-rotated} to locate the minimum, but track and
return $\textit{lo}$ (the index) instead of $\textit{arr}[\textit{lo}]$
(the value).
\end{ideabox}

\subsection*{Algorithm}

Identical to Algorithm~\ref{sec:bs-find-min-rotated}'s \textsc{FindMin}
procedure, returning $\textit{lo}$ instead of $\textit{arr}[\textit{lo}]$.

\subsection*{Dry run}

\dryrun{Take $\textit{arr} = [4,5,6,7,0,1,2]$ (same array as
Section~\ref{sec:bs-find-min-rotated}).}

The binary search trace is identical to before, terminating with
$\textit{lo} = 4$. So the array has been rotated $4$ times -- indeed, the
original sorted array would have been $[0,1,2,4,5,6,7]$, and rotating it
left by $4$ positions gives exactly $[4,5,6,7,0,1,2]$.

\subsection*{C++ implementation}

\begin{lstlisting}
#include <bits/stdc++.h>
using namespace std;

int countRotations(vector<int>& arr) {
    int n = arr.size();
    int lo = 0, hi = n - 1;

    while (lo < hi) {
        if (arr[lo] <= arr[hi]) break; // already sorted

        int mid = lo + (hi - lo) / 2;
        if (arr[mid] >= arr[lo]) {
            lo = mid + 1;
        } else {
            hi = mid;
        }
    }
    return lo; // index of the minimum = number of rotations
}

int main() {
    vector<int> arr = {4, 5, 6, 7, 0, 1, 2};
    cout << "Number of rotations: " << countRotations(arr) << endl;
    return 0;
}
\end{lstlisting}

\begin{complexitybox}
\textbf{Time Complexity.} $O(\log n)$, identical to
Section~\ref{sec:bs-find-min-rotated}.
\textbf{Space Complexity.} $O(1)$ auxiliary space.
\end{complexitybox}

\begin{takeawaybox}
When a new problem statement turns out to want the \emph{index} of a
quantity you already know how to binary search for (rather than its
\emph{value}), there is usually nothing new to derive -- just track the
index alongside, or instead of, the value throughout the same search.
\end{takeawaybox}

\section{Single Element in a Sorted Array}
\label{sec:bs-single-element}

\problemstatement{We are given a sorted array in which \textbf{every
element appears exactly twice}, except for one element that appears
\textbf{exactly once}. Find that single element, in $O(\log n)$ time.}

\begin{patternbox}
At first glance this looks like it has nothing to do with sortedness --
``find the element that appears once'' sounds like a job for hashing or
XOR, both $O(n)$. The keyword that should stop you is \emph{``sorted
array''} appearing in the same breath as a uniqueness condition: whenever a
problem about counts, pairs, or uniqueness also guarantees sortedness,
suspect that sortedness is providing exploitable \emph{positional}
structure -- here, specifically, that pairs of equal elements occupy
predictable (even, odd) index pairs \emph{before} the single element, and
that this pairing pattern shifts by one position immediately
\emph{after} it.
\end{patternbox}

\subsection*{Understanding the problem carefully}

Suppose the single (unpaired) element sits at index $s$. For every pair
of equal elements before $s$, since each pair occupies two consecutive
array slots and there is an even number of elements total before $s$ (every
element before $s$ is properly paired), the first element of each such pair
lands on an \textbf{even} index, and its duplicate lands on the very next
(\textbf{odd}) index. That is: for $i < s$, pairs look like
$(\textit{arr}[2k], \textit{arr}[2k+1])$ with $\textit{arr}[2k] =
\textit{arr}[2k+1]$. After $s$, the single element has shifted the parity
of every subsequent pair by one position, so pairs there look like
$(\textit{arr}[2k+1], \textit{arr}[2k+2])$ instead.

\begin{ideabox}[Candidate Space and Predicate]
The candidate space is indices $0, \dots, n-1$ (restricted to even indices,
by construction). The predicate is
$P(\textit{even index } i) = (\textit{arr}[i] = \textit{arr}[i+1])$ --
\texttt{true} for every even index strictly before the single element, and
\texttt{false} from the single element's (even-aligned) position onward.
We binary search for the first even index where $P$ becomes \texttt{false}
-- that index holds the single element.
\end{ideabox}

\subsection*{Why forcing \textit{mid} to be even is necessary and correct}

The predicate $P$ above is only meaningful when evaluated at \emph{even}
indices, because the pairing pattern is defined relative to even/odd
positions, not relative to $\textit{mid}$'s raw value. So at each step, if
the computed midpoint $\textit{mid}$ happens to be odd, we adjust it down
by one (to the preceding even index) before testing $P$. After this
adjustment, the same monotonicity argument as ordinary binary search
applies directly: $P$ is true for an initial contiguous run of even
indices and false afterward, with the crossover landing exactly on the
single element's position.

\subsection*{Algorithm}

\begin{enumerate}
  \item Initialize $\textit{lo} \gets 0$, $\textit{hi} \gets n-2$
        (we only need to compare each candidate to its right neighbor).
  \item While $\textit{lo} < \textit{hi}$:
  \begin{enumerate}
    \item $\textit{mid} \gets \textit{lo}+(\textit{hi}-\textit{lo})/2$.
    \item If \textit{mid} is odd: $\textit{mid} \gets \textit{mid}-1$
          (force to even).
    \item If $\textit{arr}[\textit{mid}] = \textit{arr}[\textit{mid}+1]$:
          the single element is to the right; $\textit{lo} \gets
          \textit{mid}+2$.
    \item Else: the single element is at \textit{mid} or to the left;
          $\textit{hi} \gets \textit{mid}$.
  \end{enumerate}
  \item Return $\textit{arr}[\textit{lo}]$.
\end{enumerate}

\begin{algorithm}[H]
\caption{Single Element in a Sorted Array}
\begin{algorithmic}[1]
\Procedure{SingleElement}{\textit{arr}}
  \State $\textit{lo} \gets 0$, $\textit{hi} \gets n-2$
  \While{$\textit{lo} < \textit{hi}$}
    \State $\textit{mid} \gets \textit{lo}+(\textit{hi}-\textit{lo})/2$
    \If{\textit{mid} is odd}
      \State $\textit{mid} \gets \textit{mid} - 1$
    \EndIf
    \If{$\textit{arr}[\textit{mid}] = \textit{arr}[\textit{mid}+1]$}
      \State $\textit{lo} \gets \textit{mid} + 2$
    \Else
      \State $\textit{hi} \gets \textit{mid}$
    \EndIf
  \EndWhile
  \State \Return $\textit{arr}[\textit{lo}]$
\EndProcedure
\end{algorithmic}
\end{algorithm}

\subsection*{Dry run}

\dryrun{Take $\textit{arr} = [1,1,2,3,3,4,4,8,8]$ ($n=9$).}

\begin{itemize}
  \item $\textit{lo}=0,\textit{hi}=7$: $\textit{mid}=3$ (odd) $\to 2$.
        $\textit{arr}[2]=2$, $\textit{arr}[3]=3$. Not equal. $\textit{hi}=2$.
  \item $\textit{lo}=0,\textit{hi}=2$: $\textit{mid}=1$ (odd) $\to 0$.
        $\textit{arr}[0]=1$, $\textit{arr}[1]=1$. Equal. $\textit{lo}=2$.
  \item $\textit{lo}=2,\textit{hi}=2$: loop ends ($\textit{lo}=\textit{hi}$).
\end{itemize}

The single element is $\textit{arr}[2] = 2$, which is correct.

\subsection*{C++ implementation}

\begin{lstlisting}
#include <bits/stdc++.h>
using namespace std;

int singleElement(vector<int>& arr) {
    int n = arr.size();
    int lo = 0, hi = n - 2;

    while (lo < hi) {
        int mid = lo + (hi - lo) / 2;
        if (mid % 2 == 1) mid--;

        if (arr[mid] == arr[mid + 1]) {
            lo = mid + 2;
        } else {
            hi = mid;
        }
    }
    return arr[lo];
}

int main() {
    vector<int> arr = {1, 1, 2, 3, 3, 4, 4, 8, 8};
    cout << "Single element: " << singleElement(arr) << endl;
    return 0;
}
\end{lstlisting}

\begin{complexitybox}
\textbf{Time Complexity.} The search window halves each iteration (the
parity adjustment only ever shifts \textit{mid} by one, not affecting the
overall halving rate), giving
\[
  O(\log n).
\]
\textbf{Space Complexity.} $O(1)$ auxiliary space.
\end{complexitybox}

\begin{takeawaybox}
Sortedness can encode positional structure beyond simple ordering --
here, sortedness plus a pairing guarantee creates an exploitable
\emph{parity} pattern. Whenever a sorted-array problem also involves
counts, pairs, or duplicates, look for whether index parity (even/odd
position) is doing some of the work that the array's values alone do in
simpler binary search problems.
\end{takeawaybox}

\section{Find Peak Element}
\label{sec:bs-peak-1d}

\problemstatement{We are given an array $\textit{arr}$ of $n$ distinct
integers (no two adjacent elements are equal), where the array is
conceptually bordered by $-\infty$ on both ends (i.e.\ $\textit{arr}[-1] =
\textit{arr}[n] = -\infty$). An index $i$ is a \textbf{peak} if
$\textit{arr}[i]$ is strictly greater than both of its neighbors. We must
find the index of \textbf{any one} peak element, in $O(\log n)$ time.}

\begin{patternbox}
This problem has \textbf{no sortedness at all} -- the array can be
completely jumbled. This is the most important lesson in this section:
the keyword to look for is not ``sorted,'' but the guarantee of a
\textbf{boundary condition combined with a local comparison rule} that
forces a peak to exist and forces a specific monotonic structure on the
direction of ``ascent.'' Whenever you see ``array bordered by $-\infty$''
or an equivalent guarantee (e.g.\ ``the array is strictly increasing then
strictly decreasing''), suspect a peak-finding binary search even though no
sortedness is mentioned.
\end{patternbox}

\subsection*{Understanding the problem carefully}

Consider the midpoint $\textit{mid}$ and compare $\textit{arr}[\textit{mid}]$
to its right neighbor $\textit{arr}[\textit{mid}+1]$ (treating an
out-of-range neighbor as $-\infty$ per the problem's border convention).

\begin{itemize}
  \item If $\textit{arr}[\textit{mid}] < \textit{arr}[\textit{mid}+1]$:
        the sequence is ``climbing'' at \textit{mid}. Because the array is
        bordered by $-\infty$ on the right end too, \emph{somewhere} to the
        right of \textit{mid} the climbing must eventually stop and turn
        into a peak (it cannot climb forever, since it must come back down
        to $-\infty$ at the far border) -- so a peak is guaranteed to exist
        in $(\textit{mid}, \textit{hi}]$.
  \item If $\textit{arr}[\textit{mid}] > \textit{arr}[\textit{mid}+1]$:
        the sequence is ``descending'' at \textit{mid}, meaning
        \textit{mid} is locally higher than its right neighbor. By a
        symmetric argument applied leftward (the array is bordered by
        $-\infty$ on the left too), a peak is guaranteed to exist somewhere
        in $[\textit{lo}, \textit{mid}]$, \emph{including possibly
        \textit{mid} itself}.
\end{itemize}

\begin{ideabox}[Greedy Choice]
If $\textit{arr}[\textit{mid}] < \textit{arr}[\textit{mid}+1]$, discard the
left half including \textit{mid}, since a peak is guaranteed further
right: $\textit{lo} \gets \textit{mid}+1$. Otherwise, discard the right
half excluding \textit{mid}, since a peak is guaranteed at \textit{mid} or
further left: $\textit{hi} \gets \textit{mid}$.
\end{ideabox}

\subsection*{Why a peak is always guaranteed to exist in the kept half}

This is an existence argument, not a uniqueness or value-comparison
argument, which makes it a structurally different kind of binary search
correctness proof from the rest of this chapter. The guarantee rests
entirely on the $-\infty$ borders: if we always move \emph{toward}
increasing values, we are guaranteed to eventually hit a local maximum
before falling off either border, because the function value cannot
increase forever in a finite array, and any point where increase changes
to decrease (or where increase meets the $-\infty$ border with the array
ascending all the way to its last real element) is, by definition, a peak.
This is the same idea used to prove that a continuous function climbing
from $-\infty$ must have a maximum -- a discrete, array-based analogue.

\subsection*{Algorithm}

\begin{enumerate}
  \item Initialize $\textit{lo} \gets 0$, $\textit{hi} \gets n-1$.
  \item While $\textit{lo} < \textit{hi}$:
  \begin{enumerate}
    \item $\textit{mid} \gets \textit{lo}+(\textit{hi}-\textit{lo})/2$.
    \item If $\textit{arr}[\textit{mid}] < \textit{arr}[\textit{mid}+1]$:
          $\textit{lo} \gets \textit{mid}+1$.
    \item Else: $\textit{hi} \gets \textit{mid}$.
  \end{enumerate}
  \item Return \textit{lo} (equivalently \textit{hi}; they have converged).
\end{enumerate}

\begin{algorithm}[H]
\caption{Find Peak Element}
\begin{algorithmic}[1]
\Procedure{FindPeak}{\textit{arr}}
  \State $\textit{lo} \gets 0$, $\textit{hi} \gets n-1$
  \While{$\textit{lo} < \textit{hi}$}
    \State $\textit{mid} \gets \textit{lo}+(\textit{hi}-\textit{lo})/2$
    \If{$\textit{arr}[\textit{mid}] < \textit{arr}[\textit{mid}+1]$}
      \State $\textit{lo} \gets \textit{mid} + 1$
    \Else
      \State $\textit{hi} \gets \textit{mid}$
    \EndIf
  \EndWhile
  \State \Return \textit{lo}
\EndProcedure
\end{algorithmic}
\end{algorithm}

\subsection*{Dry run}

\dryrun{Take $\textit{arr} = [1, 2, 1, 3, 5, 6, 4]$.}

\begin{itemize}
  \item $\textit{lo}=0,\textit{hi}=6$: $\textit{mid}=3$,
        $\textit{arr}[3]=3 < \textit{arr}[4]=5$. $\textit{lo}=4$.
  \item $\textit{lo}=4,\textit{hi}=6$: $\textit{mid}=5$,
        $\textit{arr}[5]=6 < \textit{arr}[6]=4$? No. $\textit{hi}=5$.
  \item $\textit{lo}=4,\textit{hi}=5$: $\textit{mid}=4$,
        $\textit{arr}[4]=5 < \textit{arr}[5]=6$. $\textit{lo}=5$.
  \item $\textit{lo}=5=\textit{hi}$: loop ends.
\end{itemize}

The peak index is $5$, where $\textit{arr}[5]=6$, which is indeed greater
than both neighbors $\textit{arr}[4]=5$ and $\textit{arr}[6]=4$. (Index $1$,
where $\textit{arr}[1]=2$, is also a valid peak; the problem only requires
\emph{any} one peak, and this algorithm happens to converge to the one on
the side it explored.)

\subsection*{C++ implementation}

\begin{lstlisting}
#include <bits/stdc++.h>
using namespace std;

int findPeak(vector<int>& arr) {
    int n = arr.size();
    int lo = 0, hi = n - 1;

    while (lo < hi) {
        int mid = lo + (hi - lo) / 2;
        if (arr[mid] < arr[mid + 1]) {
            lo = mid + 1;
        } else {
            hi = mid;
        }
    }
    return lo;
}

int main() {
    vector<int> arr = {1, 2, 1, 3, 5, 6, 4};
    cout << "Peak index: " << findPeak(arr) << endl;
    return 0;
}
\end{lstlisting}

\begin{complexitybox}
\textbf{Time Complexity.} The search window halves each iteration, giving
\[
  O(\log n).
\]
\textbf{Space Complexity.} $O(1)$ auxiliary space.
\end{complexitybox}

\begin{takeawaybox}
Binary search does not require sortedness -- it requires a monotonic,
provable rule for discarding half the search space. Here, the rule comes
from a local ascent/descent comparison plus a boundary guarantee, not from
any global ordering of the array's values. Whenever you see ``find a local
maximum/minimum,'' check whether boundary conditions and a single local
comparison can replace the sortedness argument used elsewhere in this
chapter.
\end{takeawaybox}

\section{Find Square Root of a Number using Binary Search}
\label{sec:bs-sqrt}

\noindent We now begin the second major group of problems in this chapter:
\textbf{binary search on the answer}. In every problem so far, the
candidate space was a set of array indices. From here on, there is often no
array at all -- the candidate space is a range of numbers that could be the
\emph{answer itself}, and we binary search directly over that range,
checking feasibility of each candidate with a helper function.

\problemstatement{Given a non-negative integer $n$, find $\lfloor
\sqrt{n} \rfloor$, the floor of its square root, without using a built-in
square root function -- and ideally without testing every integer from $1$
to $n$ one at a time.}

\begin{patternbox}
\emph{``Find the floor of the square root''} doesn't mention an array at
all, which is exactly the point: the phrase to watch for is \emph{``find
the largest/smallest value such that some computable condition holds.''}
Here, the condition is $\textit{candidate}^2 \le n$. This condition is
monotonic in \textit{candidate} -- true for small candidates, false for
large ones, with exactly one crossover -- which is all binary search ever
needed in the first place. The array, in earlier sections, was just one
particular way of presenting a monotonic candidate space; an integer range
works just as well.
\end{patternbox}

\subsection*{Understanding the problem carefully}

Define $P(\textit{candidate}) = (\textit{candidate}^2 \le n)$ for
candidates ranging over $0, 1, 2, \dots, n$. As \textit{candidate}
increases, $\textit{candidate}^2$ increases too (squaring is monotonic for
non-negative numbers), so $P$ is \texttt{true} for every candidate up to
$\lfloor \sqrt{n} \rfloor$ and \texttt{false} for every candidate beyond
it. The floor of the square root is exactly the \emph{largest} candidate
for which $P$ holds.

\begin{ideabox}[Candidate Space and Predicate]
The candidate space is the integer range $[0, n]$ (or, more tightly,
$[1, n]$ for $n \ge 1$). The predicate is
$P(\textit{candidate}) = (\textit{candidate}^2 \le n)$, and we want the
largest candidate where $P$ holds -- the same ``find the rightmost true
position'' template used for floor in Section~\ref{sec:bs-floor-ceil},
just applied to a numeric range instead of array values.
\end{ideabox}

\subsection*{Why binary searching the answer directly is valid}

There is nothing conceptually new here beyond recognizing that the
candidate space need not be array indices. The monotonicity of $P$ is
guaranteed by basic algebra ($x \le y \Rightarrow x^2 \le y^2$ for
non-negative $x, y$), which plays exactly the role that array sortedness
played in earlier sections: it guarantees a single, well-defined crossover
point that binary search can locate in logarithmically many evaluations of
$P$, each evaluation costing only $O(1)$ (a single multiplication).

\subsection*{Algorithm}

\begin{enumerate}
  \item Initialize $\textit{lo} \gets 0$, $\textit{hi} \gets n$,
        $\textit{ans} \gets 0$.
  \item While $\textit{lo} \le \textit{hi}$:
  \begin{enumerate}
    \item $\textit{mid} \gets \textit{lo}+(\textit{hi}-\textit{lo})/2$.
    \item If $\textit{mid} \times \textit{mid} \le n$: record
          $\textit{ans} \gets \textit{mid}$; search right:
          $\textit{lo} \gets \textit{mid}+1$.
    \item Else: search left: $\textit{hi} \gets \textit{mid}-1$.
  \end{enumerate}
  \item Return \textit{ans}.
\end{enumerate}

\begin{algorithm}[H]
\caption{Floor of Square Root}
\begin{algorithmic}[1]
\Procedure{FloorSqrt}{$n$}
  \State $\textit{lo} \gets 0$, $\textit{hi} \gets n$, $\textit{ans} \gets 0$
  \While{$\textit{lo} \le \textit{hi}$}
    \State $\textit{mid} \gets \textit{lo}+(\textit{hi}-\textit{lo})/2$
    \If{$\textit{mid} \times \textit{mid} \le n$}
      \State $\textit{ans} \gets \textit{mid}$
      \State $\textit{lo} \gets \textit{mid} + 1$
    \Else
      \State $\textit{hi} \gets \textit{mid} - 1$
    \EndIf
  \EndWhile
  \State \Return \textit{ans}
\EndProcedure
\end{algorithmic}
\end{algorithm}

\subsection*{Dry run}

\dryrun{Take $n = 28$.}

\begin{itemize}
  \item $\textit{lo}=0,\textit{hi}=28$: $\textit{mid}=14$, $14^2=196 \le
        28$? No. $\textit{hi}=13$.
  \item $\textit{lo}=0,\textit{hi}=13$: $\textit{mid}=6$, $6^2=36 \le 28$?
        No. $\textit{hi}=5$.
  \item $\textit{lo}=0,\textit{hi}=5$: $\textit{mid}=2$, $2^2=4 \le 28$?
        Yes. $\textit{ans}=2$, $\textit{lo}=3$.
  \item $\textit{lo}=3,\textit{hi}=5$: $\textit{mid}=4$, $4^2=16 \le 28$?
        Yes. $\textit{ans}=4$, $\textit{lo}=5$.
  \item $\textit{lo}=5,\textit{hi}=5$: $\textit{mid}=5$, $5^2=25 \le 28$?
        Yes. $\textit{ans}=5$, $\textit{lo}=6$.
  \item $\textit{lo}=6 > \textit{hi}=5$: loop ends.
\end{itemize}

The answer is $5$, and indeed $5^2 = 25 \le 28 < 36 = 6^2$.

\subsection*{C++ implementation}

\begin{lstlisting}
#include <bits/stdc++.h>
using namespace std;

int floorSqrt(int n) {
    int lo = 0, hi = n, ans = 0;

    while (lo <= hi) {
        long long mid = lo + (hi - lo) / 2;
        if (mid * mid <= n) {
            ans = mid;
            lo = mid + 1;
        } else {
            hi = mid - 1;
        }
    }
    return ans;
}

int main() {
    cout << "Floor sqrt of 28: " << floorSqrt(28) << endl;
    return 0;
}
\end{lstlisting}

\begin{complexitybox}
\textbf{Time Complexity.} The candidate range $[0, n]$ halves each
iteration, giving
\[
  O(\log n).
\]
\textbf{Space Complexity.} $O(1)$ auxiliary space.
\end{complexitybox}

\begin{takeawaybox}
Binary search on the answer requires exactly two ingredients: a numeric
candidate range, and a feasibility predicate that is monotonic across that
range. Once both exist, the same \texttt{while(lo <= hi)} skeleton from
array-based binary search applies verbatim -- only the meaning of
\textit{lo}, \textit{hi}, and the predicate check change.
\end{takeawaybox}

\section{Find the Nth Root of a Number using Binary Search}
\label{sec:bs-nth-root}

\problemstatement{Given two positive integers $n$ and $m$, find $x$ such
that $x^n = m$, where $x$ is a positive integer. If no such integer $x$
exists, report $-1$.}

\begin{patternbox}
This is the direct generalization of Section~\ref{sec:bs-sqrt} from
``square root'' to ``$n$-th root.'' The keyword pattern is identical:
\emph{``find the value $x$ such that $x$ raised to some power equals (or
is bounded by) a target.''} The only change is the predicate's exponent,
from $2$ to a general $n$ -- everything else about the binary-search-on-
the-answer template carries over unchanged.
\end{patternbox}

\subsection*{Understanding the problem carefully}

Define $P(\textit{candidate}) = (\textit{candidate}^n \le m)$ for
candidates in $[1, m]$. Just as squaring is monotonic for non-negative
integers, raising to any fixed positive power $n$ is monotonic too:
$x \le y \Rightarrow x^n \le y^n$. So $P$ is \texttt{true} for a prefix of
candidates and \texttt{false} afterward, with a single crossover -- the
largest candidate satisfying $P$ is a perfect $n$-th root only if
$\textit{candidate}^n$ exactly equals $m$ at that point; otherwise, no
exact integer $n$-th root exists.

\begin{ideabox}[Candidate Space and Predicate]
The candidate space is integers $[1, m]$. The predicate is
$P(\textit{candidate}) = (\textit{candidate}^n \le m)$. Track the largest
candidate for which $P$ holds; at the end, check whether that candidate,
raised to the $n$-th power, exactly equals $m$.
\end{ideabox}

\subsection*{Why we must verify exactness at the end}

Unlike the square-root problem, which explicitly asked for the
\emph{floor} of the root (so any final value of \textit{ans} is a valid
answer), this problem demands an \emph{exact} integer root, or $-1$ if
none exists. So the binary search itself is unchanged -- it always finds
the largest candidate whose $n$-th power does not exceed $m$ -- but we add
one final check: if that candidate's $n$-th power is not exactly $m$,
there is no integer solution, and we must report $-1$ rather than the
nearest approximation.

\subsection*{Algorithm}

\begin{enumerate}
  \item Initialize $\textit{lo} \gets 1$, $\textit{hi} \gets m$.
  \item While $\textit{lo} \le \textit{hi}$:
  \begin{enumerate}
    \item $\textit{mid} \gets \textit{lo}+(\textit{hi}-\textit{lo})/2$.
    \item Compute $p \gets \textit{mid}^n$ (stop early/clamp if it exceeds
          $m$, to avoid overflow for large $n$).
    \item If $p = m$: return \textit{mid}.
    \item If $p < m$: $\textit{lo} \gets \textit{mid}+1$.
    \item Else: $\textit{hi} \gets \textit{mid}-1$.
  \end{enumerate}
  \item If the loop ends without returning, return $-1$.
\end{enumerate}

\begin{algorithm}[H]
\caption{Nth Root}
\begin{algorithmic}[1]
\Procedure{NthRoot}{$n, m$}
  \State $\textit{lo} \gets 1$, $\textit{hi} \gets m$
  \While{$\textit{lo} \le \textit{hi}$}
    \State $\textit{mid} \gets \textit{lo}+(\textit{hi}-\textit{lo})/2$
    \State $p \gets \textit{mid}^n$ \Comment{computed with early-exit on overflow}
    \If{$p = m$} \Return \textit{mid} \EndIf
    \If{$p < m$}
      \State $\textit{lo} \gets \textit{mid} + 1$
    \Else
      \State $\textit{hi} \gets \textit{mid} - 1$
    \EndIf
  \EndWhile
  \State \Return $-1$
\EndProcedure
\end{algorithmic}
\end{algorithm}

\subsection*{Dry run}

\dryrun{Take $n = 3$, $m = 27$ (find the cube root).}

\begin{itemize}
  \item $\textit{lo}=1,\textit{hi}=27$: $\textit{mid}=14$, $14^3=2744 \le
        27$? No, $p > m$. $\textit{hi}=13$.
  \item $\textit{lo}=1,\textit{hi}=13$: $\textit{mid}=7$, $7^3=343 > 27$.
        $\textit{hi}=6$.
  \item $\textit{lo}=1,\textit{hi}=6$: $\textit{mid}=3$, $3^3=27=27$.
        Return $3$.
\end{itemize}

The answer is $3$, since $3^3 = 27$.

Now take $n=4$, $m=20$. No integer fourth root of $20$ exists ($2^4=16$,
$3^4=81$), so the binary search will narrow down without ever finding
$p=m$, and the final answer will correctly be $-1$.

\subsection*{C++ implementation}

\begin{lstlisting}
#include <bits/stdc++.h>
using namespace std;

// Computes base^exp, capped to indicate ">m" without overflow.
long long power(long long base, int exp, long long cap) {
    long long result = 1;
    for (int i = 0; i < exp; i++) {
        result *= base;
        if (result > cap) return cap + 1; // overflow guard
    }
    return result;
}

int nthRoot(int n, int m) {
    int lo = 1, hi = m;

    while (lo <= hi) {
        int mid = lo + (hi - lo) / 2;
        long long p = power(mid, n, m);

        if (p == m) return mid;
        else if (p < m) lo = mid + 1;
        else hi = mid - 1;
    }
    return -1;
}

int main() {
    cout << "Cube root of 27: " << nthRoot(3, 27) << endl;
    cout << "4th root of 20: " << nthRoot(4, 20) << endl;
    return 0;
}
\end{lstlisting}

\begin{complexitybox}
\textbf{Time Complexity.} The candidate range $[1, m]$ halves each
iteration, and each predicate evaluation costs $O(n)$ to compute the
power (or $O(\log n)$ with fast exponentiation), giving an overall time
complexity of
\[
  O(n \log m)
\]
(or $O(\log n \log m)$ with fast exponentiation).
\textbf{Space Complexity.} $O(1)$ auxiliary space.
\end{complexitybox}

\begin{takeawaybox}
Generalizing a binary-search-on-the-answer problem from a fixed exponent
(square root, exponent $2$) to a variable exponent ($n$-th root) changes
nothing about the search itself -- only the cost of evaluating the
predicate changes, since computing $\textit{candidate}^n$ is more expensive
than $\textit{candidate}^2$. Always separate ``how do I search the
candidate space'' from ``how expensive is one feasibility check,'' since
they contribute independently to the total complexity.
\end{takeawaybox}

\section{Koko Eating Bananas}
\label{sec:bs-koko}

\problemstatement{There are $n$ piles of bananas, where pile $i$ has
$\textit{piles}[i]$ bananas. Koko can eat at most $k$ bananas from a single
pile in one hour (if a pile has fewer than $k$ bananas, she finishes that
pile in that hour and cannot start another pile in the same hour). Koko
wants to eat all the bananas within $h$ hours. Find the
\textbf{minimum} integer eating speed $k$ that allows her to finish within
$h$ hours.}

\begin{patternbox}
\emph{``Find the minimum speed/rate/capacity such that some task completes
within a time/resource limit''} is the archetypal \emph{binary search on
the answer} problem. The two keywords to look for together are
\textbf{``minimize the (speed/capacity/rate)''} and a \textbf{feasibility
constraint} (here, a deadline of $h$ hours) that becomes easier to satisfy
as the candidate answer increases. This combination -- minimize $x$ subject
to ``feasible($x$)'' becoming \texttt{true} once $x$ is large enough -- is
the single most common shape of binary-search-on-the-answer problems on
this sheet, and Koko Eating Bananas is its canonical introduction.
\end{patternbox}

\subsection*{Understanding the problem carefully}

For a fixed candidate speed $k$, the number of hours Koko needs to finish
pile $i$ is $\lceil \textit{piles}[i] / k \rceil$ (she needs a whole extra
hour for any remainder), so the total hours needed across all piles is
\[
  \textit{hoursNeeded}(k) = \sum_{i=1}^{n} \left\lceil
  \frac{\textit{piles}[i]}{k} \right\rceil.
\]
As $k$ increases, each term $\lceil \textit{piles}[i]/k \rceil$ can only
decrease or stay the same (eating faster never requires more hours for the
same pile), so $\textit{hoursNeeded}(k)$ is a non-increasing function of
$k$. Define the feasibility predicate
$P(k) = (\textit{hoursNeeded}(k) \le h)$. Since $\textit{hoursNeeded}$ is
non-increasing in $k$, $P$ is \texttt{false} for small $k$ and
\texttt{true} from some crossover point onward -- exactly the monotonic
structure binary search needs, with the answer being the smallest $k$
where $P$ becomes \texttt{true}.

\begin{ideabox}[Candidate Space and Predicate]
The candidate space is integer eating speeds $k \in [1, \max(\textit{piles})]$
(eating faster than the largest pile is never useful, since any pile is
finished in exactly one hour once $k$ reaches its size). The predicate is
$P(k) = (\textit{hoursNeeded}(k) \le h)$, and we want the smallest $k$
where $P$ holds.
\end{ideabox}

\subsection*{Why this reduces feasibility-checking to a simple scan}

The crucial design decision is to separate two concerns: \emph{deciding
which $k$ to try next} (binary search's job) and \emph{checking whether a
specific $k$ works} (a simple $O(n)$ linear scan over the piles, summing
up $\lceil \textit{piles}[i]/k \rceil$ for the fixed candidate $k$). This
separation is the heart of every binary-search-on-the-answer problem: we
never need to reason about the feasibility check's internal monotonicity
in $i$ -- only about the feasibility \emph{function as a whole}'s
monotonicity in the single candidate parameter $k$, which we established
above using a much simpler, purely algebraic argument (more bananas per
hour can never increase the hours needed).

\subsection*{Algorithm}

\begin{enumerate}
  \item Initialize $\textit{lo} \gets 1$,
        $\textit{hi} \gets \max(\textit{piles})$.
  \item While $\textit{lo} \le \textit{hi}$:
  \begin{enumerate}
    \item $\textit{mid} \gets \textit{lo}+(\textit{hi}-\textit{lo})/2$.
    \item Compute $\textit{hoursNeeded}(\textit{mid})$ by summing
          $\lceil \textit{piles}[i] / \textit{mid} \rceil$ over all piles.
    \item If $\textit{hoursNeeded}(\textit{mid}) \le h$: this speed works;
          record it and try slower: $\textit{hi} \gets \textit{mid}-1$.
    \item Else: too slow; speed up: $\textit{lo} \gets \textit{mid}+1$.
  \end{enumerate}
  \item Return the recorded answer (equivalently, $\textit{lo}$ once the
        loop ends).
\end{enumerate}

\begin{algorithm}[H]
\caption{Koko Eating Bananas}
\begin{algorithmic}[1]
\Procedure{MinEatingSpeed}{\textit{piles}, $h$}
  \State $\textit{lo} \gets 1$, $\textit{hi} \gets \max(\textit{piles})$
  \While{$\textit{lo} \le \textit{hi}$}
    \State $\textit{mid} \gets \textit{lo}+(\textit{hi}-\textit{lo})/2$
    \State $\textit{hours} \gets \sum_i \lceil \textit{piles}[i] / \textit{mid} \rceil$
    \If{$\textit{hours} \le h$}
      \State $\textit{hi} \gets \textit{mid} - 1$
    \Else
      \State $\textit{lo} \gets \textit{mid} + 1$
    \EndIf
  \EndWhile
  \State \Return \textit{lo}
\EndProcedure
\end{algorithmic}
\end{algorithm}

\subsection*{Dry run}

\dryrun{Take $\textit{piles} = [3, 6, 7, 11]$, $h = 8$.}

$\max(\textit{piles}) = 11$.

\begin{itemize}
  \item $\textit{lo}=1,\textit{hi}=11$: $\textit{mid}=6$. Hours:
        $\lceil3/6\rceil+\lceil6/6\rceil+\lceil7/6\rceil+\lceil11/6\rceil
        = 1+1+2+2=6 \le 8$. Feasible. $\textit{hi}=5$.
  \item $\textit{lo}=1,\textit{hi}=5$: $\textit{mid}=3$. Hours:
        $\lceil3/3\rceil+\lceil6/3\rceil+\lceil7/3\rceil+\lceil11/3\rceil
        =1+2+3+4=10 > 8$. Infeasible. $\textit{lo}=4$.
  \item $\textit{lo}=4,\textit{hi}=5$: $\textit{mid}=4$. Hours:
        $\lceil3/4\rceil+\lceil6/4\rceil+\lceil7/4\rceil+\lceil11/4\rceil
        =1+2+2+3=8 \le 8$. Feasible. $\textit{hi}=3$.
  \item $\textit{lo}=4 > \textit{hi}=3$: loop ends.
\end{itemize}

The answer is $\textit{lo} = 4$: the minimum eating speed is $4$ bananas
per hour.

\subsection*{C++ implementation}

\begin{lstlisting}
#include <bits/stdc++.h>
using namespace std;

long long hoursNeeded(vector<int>& piles, int k) {
    long long hours = 0;
    for (int p : piles) {
        hours += (p + k - 1) / k; // ceil(p / k)
    }
    return hours;
}

int minEatingSpeed(vector<int>& piles, int h) {
    int lo = 1, hi = *max_element(piles.begin(), piles.end());

    while (lo <= hi) {
        int mid = lo + (hi - lo) / 2;
        if (hoursNeeded(piles, mid) <= h) {
            hi = mid - 1;
        } else {
            lo = mid + 1;
        }
    }
    return lo;
}

int main() {
    vector<int> piles = {3, 6, 7, 11};
    cout << "Minimum eating speed: "
         << minEatingSpeed(piles, 8) << endl;
    return 0;
}
\end{lstlisting}

\begin{complexitybox}
\textbf{Time Complexity.} The candidate speed range $[1, \max(\textit{piles})]$
halves each iteration ($O(\log(\max(\textit{piles})))$ iterations), and
each feasibility check costs $O(n)$ to scan all piles, giving an overall
time complexity of
\[
  O(n \log(\max(\textit{piles}))).
\]
\textbf{Space Complexity.} $O(1)$ auxiliary space.
\end{complexitybox}

\begin{takeawaybox}
The general recipe for binary search on the answer: (1) identify the
quantity being minimized or maximized as the candidate; (2) write a
feasibility function that, for a fixed candidate, answers a yes/no
question with a simple (often linear) computation; (3) prove that
feasibility is monotonic in the candidate using simple algebraic reasoning
about the feasibility function's components; (4) binary search the
candidate range using that feasibility check as the predicate. Every
remaining ``binary search on the answer'' problem in this chapter follows
exactly this four-step recipe.
\end{takeawaybox}

\section{Minimum Number of Days to Make M Bouquets}
\label{sec:bs-bouquets}

\problemstatement{We are given an array $\textit{bloomDay}$ of $n$ flowers,
where $\textit{bloomDay}[i]$ is the day on which flower $i$ blooms. We need
to make $m$ bouquets, each requiring exactly $k$ \textbf{adjacent}
(consecutive in the array) flowers that have \emph{already bloomed}. Find
the \textbf{minimum day} on which it becomes possible to make all $m$
bouquets, or $-1$ if it is never possible (e.g.\ if $m \times k > n$).}

\begin{patternbox}
\emph{``Minimum day on which some condition becomes achievable''} is
binary search on the answer again, with the candidate now being a
\emph{day} rather than a speed (Section~\ref{sec:bs-koko}) or a numeric
root (Sections~\ref{sec:bs-sqrt}--\ref{sec:bs-nth-root}). The keyword
combination \emph{``minimum day/time such that a task becomes possible''}
is one of the most common phrasings of this pattern, and the underlying
monotonicity is intuitive: waiting longer can only help (more flowers have
bloomed), never hurt.
\end{patternbox}

\subsection*{Understanding the problem carefully}

For a fixed candidate day $d$, a flower $i$ is usable if
$\textit{bloomDay}[i] \le d$. Walking through the array in order, every
maximal run of consecutive usable flowers of length $L$ can be split into
$\lfloor L/k \rfloor$ bouquets (using groups of $k$ adjacent usable flowers
within that run; any leftover flowers in the run, fewer than $k$, cannot
form another bouquet). Summing $\lfloor L/k \rfloor$ over every maximal
usable run gives the total number of bouquets makeable by day $d$. Define
$\textit{bouquetsPossible}(d)$ as this total. As $d$ increases, no flower
that was previously usable becomes unusable, and some additional flowers
may become usable, so usable runs can only grow or merge -- meaning
$\textit{bouquetsPossible}(d)$ is non-decreasing in $d$.

\begin{ideabox}[Candidate Space and Predicate]
The candidate space is days in $[\min(\textit{bloomDay}),
\max(\textit{bloomDay})]$ (waiting before the earliest bloom or past the
latest bloom is never useful). The predicate is
$P(d) = (\textit{bouquetsPossible}(d) \ge m)$, and we want the smallest
$d$ where $P$ holds.
\end{ideabox}

\subsection*{Why merging consecutive usable flowers correctly counts bouquets}

The requirement that a bouquet's $k$ flowers be \emph{adjacent} in the
original array is what makes this feasibility check more than a simple
count of usable flowers divided by $k$: scattered usable flowers separated
by an unbloomed flower cannot be combined into one bouquet, even if their
total count would otherwise suffice. By scanning left to right and
resetting a running streak counter to $0$ whenever an unusable flower is
encountered (and incrementing it, and adding to the bouquet count
whenever the streak count reaches a multiple of $k$, otherwise just
incrementing), we correctly respect the adjacency constraint while still
computing $\textit{bouquetsPossible}(d)$ in a single $O(n)$ pass.

\subsection*{Algorithm}

\begin{enumerate}
  \item If $m \times k > n$: return $-1$ immediately (not enough flowers
        could ever exist).
  \item Initialize $\textit{lo} \gets \min(\textit{bloomDay})$,
        $\textit{hi} \gets \max(\textit{bloomDay})$.
  \item While $\textit{lo} \le \textit{hi}$:
  \begin{enumerate}
    \item $\textit{mid} \gets \textit{lo}+(\textit{hi}-\textit{lo})/2$.
    \item Compute $\textit{bouquetsPossible}(\textit{mid})$ via a single
          linear scan with a streak counter.
    \item If $\textit{bouquetsPossible}(\textit{mid}) \ge m$: feasible;
          try earlier: $\textit{hi} \gets \textit{mid}-1$.
    \item Else: infeasible; try later: $\textit{lo} \gets \textit{mid}+1$.
  \end{enumerate}
  \item Return \textit{lo}.
\end{enumerate}

\begin{algorithm}[H]
\caption{Minimum Day for M Bouquets}
\begin{algorithmic}[1]
\Procedure{MinDays}{\textit{bloomDay}, $m$, $k$}
  \If{$m \times k > n$} \Return $-1$ \EndIf
  \State $\textit{lo} \gets \min(\textit{bloomDay})$, $\textit{hi} \gets \max(\textit{bloomDay})$
  \While{$\textit{lo} \le \textit{hi}$}
    \State $\textit{mid} \gets \textit{lo}+(\textit{hi}-\textit{lo})/2$
    \State $\textit{bouquets} \gets \Call{BouquetsPossible}{\textit{bloomDay}, \textit{mid}, k}$
    \If{$\textit{bouquets} \ge m$}
      \State $\textit{hi} \gets \textit{mid} - 1$
    \Else
      \State $\textit{lo} \gets \textit{mid} + 1$
    \EndIf
  \EndWhile
  \State \Return \textit{lo}
\EndProcedure
\end{algorithmic}
\end{algorithm}

\subsection*{Dry run}

\dryrun{Take $\textit{bloomDay} = [7,7,7,7,12,7,7]$, $m=2$, $k=3$.}

\begin{itemize}
  \item $m \times k = 6 \le n = 7$, so an answer may exist.
  \item Try $d=12$ (the maximum bloom day, as an early sanity check): every
        flower is usable, giving one run of length $7$, so
        $\textit{bouquetsPossible}(12) = \lfloor 7/3 \rfloor = 2 \ge 2$.
        Feasible.
  \item Try $d=7$ (binary search would reach this candidate): flowers at
        indices $0,1,2,3,5,6$ are usable (bloom day $7 \le 7$), but index
        $4$ (bloom day $12$) is not. This splits the array into two runs:
        indices $[0,3]$ (length $4$) and $[5,6]$ (length $2$). Bouquets:
        $\lfloor 4/3 \rfloor + \lfloor 2/3 \rfloor = 1 + 0 = 1 < 2$.
        Infeasible.
\end{itemize}

The binary search narrows between $7$ and $12$; checking the remaining
candidates confirms that day $12$ is the earliest day at which $2$
bouquets become possible (the run of $4$ usable flowers before index $4$
cannot be split to free up a third group without waiting for the last
flower to bloom).

\subsection*{C++ implementation}

\begin{lstlisting}
#include <bits/stdc++.h>
using namespace std;

int bouquetsPossible(vector<int>& bloomDay, int day, int k) {
    int n = bloomDay.size();
    int streak = 0, bouquets = 0;

    for (int i = 0; i < n; i++) {
        if (bloomDay[i] <= day) {
            streak++;
            if (streak == k) {
                bouquets++;
                streak = 0;
            }
        } else {
            streak = 0;
        }
    }
    return bouquets;
}

int minDays(vector<int>& bloomDay, int m, int k) {
    int n = bloomDay.size();
    if ((long long)m * k > n) return -1;

    int lo = *min_element(bloomDay.begin(), bloomDay.end());
    int hi = *max_element(bloomDay.begin(), bloomDay.end());

    while (lo <= hi) {
        int mid = lo + (hi - lo) / 2;
        if (bouquetsPossible(bloomDay, mid, k) >= m) {
            hi = mid - 1;
        } else {
            lo = mid + 1;
        }
    }
    return lo;
}

int main() {
    vector<int> bloomDay = {7, 7, 7, 7, 12, 7, 7};
    cout << "Minimum day: " << minDays(bloomDay, 2, 3) << endl;
    return 0;
}
\end{lstlisting}

\begin{complexitybox}
\textbf{Time Complexity.} The candidate day range halves each iteration,
costing $O(\log(\max(\textit{bloomDay}) - \min(\textit{bloomDay})))$
iterations, and each feasibility check costs $O(n)$, giving an overall
time complexity of
\[
  O(n \log(\max(\textit{bloomDay}))).
\]
\textbf{Space Complexity.} $O(1)$ auxiliary space.
\end{complexitybox}

\begin{takeawaybox}
When a feasibility check involves an \emph{adjacency} or
\emph{contiguity} requirement (here, $k$ consecutive bloomed flowers),
the linear-scan feasibility function typically needs a running streak
counter that resets on disqualification, rather than a simple global
count. This adjacency-aware counting pattern recurs in
Section~\ref{sec:bs-painters-partition} and is worth recognizing as its
own small sub-template within the larger binary-search-on-the-answer
recipe.
\end{takeawaybox}

\section{Find the Smallest Divisor Given a Threshold}
\label{sec:bs-smallest-divisor}

\problemstatement{We are given an array $\textit{nums}$ of $n$ positive
integers and an integer $\textit{threshold}$. We must choose a positive
integer \textit{divisor}, and the resulting sum
$\sum_i \lceil \textit{nums}[i] / \textit{divisor} \rceil$ must be
$\le \textit{threshold}$. Find the \textbf{smallest} such divisor.}

\begin{patternbox}
This is structurally identical to Koko Eating Bananas
(Section~\ref{sec:bs-koko}): the keywords \emph{``smallest divisor''} and
\emph{``sum of ceiling divisions must not exceed a threshold''} map
directly onto \emph{``smallest eating speed''} and \emph{``total hours
must not exceed the limit.''} Once you notice the structural identity, the
entire derivation transfers without modification -- only the variable
names change.
\end{patternbox}

\subsection*{Understanding the problem carefully}

Define $\textit{sumDivisions}(\textit{divisor}) = \sum_i \lceil
\textit{nums}[i] / \textit{divisor} \rceil$. As \textit{divisor} increases,
each term $\lceil \textit{nums}[i]/\textit{divisor} \rceil$ can only
decrease or stay the same, so $\textit{sumDivisions}$ is non-increasing in
\textit{divisor} -- the same monotonicity structure as
$\textit{hoursNeeded}$ in Section~\ref{sec:bs-koko}.

\begin{ideabox}[Candidate Space and Predicate]
The candidate space is integer divisors in $[1, \max(\textit{nums})]$
(any divisor at least as large as the largest element reduces every term
to $1$, so nothing larger is ever needed). The predicate is
$P(\textit{divisor}) = (\textit{sumDivisions}(\textit{divisor}) \le
\textit{threshold})$, and we want the smallest divisor where $P$ holds.
\end{ideabox}

\subsection*{Why the derivation transfers without modification}

There is no new reasoning required: this is a direct relabeling of
Koko Eating Bananas, with $\textit{piles} \to \textit{nums}$,
$k \to \textit{divisor}$, and $h \to \textit{threshold}$. This section
exists specifically to demonstrate that once you recognize the abstract
shape (\emph{minimize a divisor/rate subject to a sum-of-ceiling-divisions
budget}), you should reuse the exact same code template rather than
rederiving it -- pattern recognition, not just algorithmic cleverness, is
the actual skill being exercised here.

\subsection*{Algorithm}

Identical in structure to Algorithm~\ref{sec:bs-koko}'s
\textsc{MinEatingSpeed}, with the feasibility check replaced by
$\textit{sumDivisions}(\textit{mid}) \le \textit{threshold}$.

\subsection*{Dry run}

\dryrun{Take $\textit{nums} = [1, 2, 5, 9]$, $\textit{threshold} = 6$.}

\begin{itemize}
  \item $\textit{lo}=1,\textit{hi}=9$: $\textit{mid}=5$. Sum:
        $\lceil1/5\rceil+\lceil2/5\rceil+\lceil5/5\rceil+\lceil9/5\rceil
        =1+1+1+2=5 \le 6$. Feasible. $\textit{hi}=4$.
  \item $\textit{lo}=1,\textit{hi}=4$: $\textit{mid}=2$. Sum:
        $\lceil1/2\rceil+\lceil2/2\rceil+\lceil5/2\rceil+\lceil9/2\rceil
        =1+1+3+5=10 > 6$. Infeasible. $\textit{lo}=3$.
  \item $\textit{lo}=3,\textit{hi}=4$: $\textit{mid}=3$. Sum:
        $\lceil1/3\rceil+\lceil2/3\rceil+\lceil5/3\rceil+\lceil9/3\rceil
        =1+1+2+3=7 > 6$. Infeasible. $\textit{lo}=4$.
  \item $\textit{lo}=4,\textit{hi}=4$: $\textit{mid}=4$. Sum:
        $\lceil1/4\rceil+\lceil2/4\rceil+\lceil5/4\rceil+\lceil9/4\rceil
        =1+1+2+3=7 > 6$. Infeasible. $\textit{lo}=5$.
  \item $\textit{lo}=5 > \textit{hi}=4$: loop ends.
\end{itemize}

The answer is $5$.

\subsection*{C++ implementation}

\begin{lstlisting}
#include <bits/stdc++.h>
using namespace std;

long long sumDivisions(vector<int>& nums, int divisor) {
    long long sum = 0;
    for (int x : nums) {
        sum += (x + divisor - 1) / divisor; // ceil(x / divisor)
    }
    return sum;
}

int smallestDivisor(vector<int>& nums, int threshold) {
    int lo = 1, hi = *max_element(nums.begin(), nums.end());

    while (lo <= hi) {
        int mid = lo + (hi - lo) / 2;
        if (sumDivisions(nums, mid) <= threshold) {
            hi = mid - 1;
        } else {
            lo = mid + 1;
        }
    }
    return lo;
}

int main() {
    vector<int> nums = {1, 2, 5, 9};
    cout << "Smallest divisor: "
         << smallestDivisor(nums, 6) << endl;
    return 0;
}
\end{lstlisting}

\begin{complexitybox}
\textbf{Time Complexity.}
\[
  O(n \log(\max(\textit{nums}))),
\]
identical in structure to Section~\ref{sec:bs-koko}.
\textbf{Space Complexity.} $O(1)$ auxiliary space.
\end{complexitybox}

\begin{takeawaybox}
Recognizing that two differently-worded problems are the same template
underneath is just as valuable as deriving a template from scratch.
Maintain a small mental (or literal) catalogue of solved
binary-search-on-the-answer templates -- ``minimize rate subject to a
ceiling-division sum budget'' is now one entry in that catalogue, covering
both this problem and Koko Eating Bananas.
\end{takeawaybox}

\section{Capacity to Ship Packages within D Days}
\label{sec:bs-ship-packages}

\problemstatement{We are given an array $\textit{weights}$ of $n$ package
weights, which must be shipped \textbf{in the given order} (no
reordering). Each day, the ship loads a contiguous run of packages, in
order, whose total weight does not exceed the ship's capacity, then sails;
the next day starts loading from the next unshipped package. We must find
the \textbf{minimum ship capacity} that allows all packages to be shipped
within $D$ days.}

\begin{patternbox}
\emph{``Minimum capacity such that all items can be processed within a day
limit, items must stay in order''} is another binary-search-on-the-answer
problem, but the feasibility check itself is structurally different from
Sections~\ref{sec:bs-koko}--\ref{sec:bs-smallest-divisor}: rather than
summing independent per-item costs, we must \textbf{greedily simulate} a
day-by-day loading process, because the order constraint means packages
cannot be considered independently of their neighbors. Whenever a
binary-search-on-the-answer problem also mentions \emph{order} or
\emph{contiguity} constraints on how items are grouped, expect the
feasibility check itself to require a small greedy simulation rather than
a single summation formula.
\end{patternbox}

\subsection*{Understanding the problem carefully}

For a fixed candidate capacity $c$, simulate the shipping process greedily:
start a new day, and keep adding the next package's weight to today's load
as long as it does not exceed $c$; the moment adding the next package
would exceed $c$, end the day and start a new one with that package. This
greedy simulation is itself optimal for \emph{counting the minimum days
needed for a fixed capacity} -- packing each day as full as possible before
moving to the next day can never use \emph{more} days than any other valid
loading strategy, because deliberately under-loading a day only defers
weight to later days without ever reducing the total weight that must
still be shipped.

Define $\textit{daysNeeded}(c)$ as the number of days this greedy
simulation takes. As $c$ increases, each day can carry at least as much as
before, so $\textit{daysNeeded}(c)$ is non-increasing in $c$.

\begin{ideabox}[Candidate Space and Predicate]
The candidate space is integer capacities in
$[\max(\textit{weights}), \sum(\textit{weights})]$: the capacity must be
at least the heaviest single package (otherwise that package alone could
never be shipped), and capacity equal to the total weight always finishes
in exactly $1$ day. The predicate is
$P(c) = (\textit{daysNeeded}(c) \le D)$, and we want the smallest $c$
where $P$ holds.
\end{ideabox}

\subsection*{Why the greedy simulation inside the predicate is valid}

This is a small but important instance of \emph{nesting} a greedy
algorithm inside a binary search predicate: the binary search handles
\emph{which} capacity to test, while a greedy simulation -- packing each
day as full as the fixed capacity allows before moving on -- handles
\emph{whether} that capacity is feasible. The greedy simulation's
correctness rests on the exchange argument that ending a day early
(loading less than the maximum the capacity allows) can never help and can
only push more weight into later days, since the in-order constraint means
no future package can be moved earlier to compensate. This nested
structure -- binary search over a monotonic outer property, with a greedy
or simulation-based feasibility check inside -- is common throughout the
``binary search on the answer'' family and is exactly what makes capacity
$c$ feasible or infeasible a well-defined, monotonic question.

\subsection*{Algorithm}

\begin{enumerate}
  \item Initialize $\textit{lo} \gets \max(\textit{weights})$,
        $\textit{hi} \gets \sum(\textit{weights})$.
  \item While $\textit{lo} \le \textit{hi}$:
  \begin{enumerate}
    \item $\textit{mid} \gets \textit{lo}+(\textit{hi}-\textit{lo})/2$.
    \item Compute $\textit{daysNeeded}(\textit{mid})$ by greedily
          simulating loading days.
    \item If $\textit{daysNeeded}(\textit{mid}) \le D$: feasible; try
          smaller: $\textit{hi} \gets \textit{mid}-1$.
    \item Else: infeasible; try larger: $\textit{lo} \gets \textit{mid}+1$.
  \end{enumerate}
  \item Return \textit{lo}.
\end{enumerate}

\begin{algorithm}[H]
\caption{Minimum Capacity to Ship within D Days}
\begin{algorithmic}[1]
\Procedure{ShipWithinDays}{\textit{weights}, $D$}
  \State $\textit{lo} \gets \max(\textit{weights})$, $\textit{hi} \gets \sum(\textit{weights})$
  \While{$\textit{lo} \le \textit{hi}$}
    \State $\textit{mid} \gets \textit{lo}+(\textit{hi}-\textit{lo})/2$
    \State $\textit{days} \gets \Call{DaysNeeded}{\textit{weights}, \textit{mid}}$
    \If{$\textit{days} \le D$}
      \State $\textit{hi} \gets \textit{mid} - 1$
    \Else
      \State $\textit{lo} \gets \textit{mid} + 1$
    \EndIf
  \EndWhile
  \State \Return \textit{lo}
\EndProcedure
\end{algorithmic}
\end{algorithm}

\subsection*{Dry run}

\dryrun{Take $\textit{weights} = [1,2,3,4,5,6,7,8,9,10]$, $D = 5$.}

$\max = 10$, $\sum = 55$.

\begin{itemize}
  \item $\textit{lo}=10,\textit{hi}=55$: $\textit{mid}=32$. Simulating: day
        loads accumulate $1{+}2{+}3{+}4{+}5{+}6{+}7=28$ (next, $8$, would
        make $36>32$, so day ends); day 2 loads $8{+}9{+}10=27$. Total
        $\textit{daysNeeded}=2 \le 5$. Feasible. $\textit{hi}=31$.
  \item Binary search continues narrowing downward; eventually it
        converges to $\textit{lo}=15$, the minimum feasible capacity (with
        capacity $15$: day loads $1{+}2{+}3{+}4{+}5=15$; day 2 loads
        $6{+}7=13$ (next, $8$, would make $21>15$); day 3 loads $8$ alone
        (next, $9$, would make $17>15$); day 4 loads $9$ alone (next,
        $10$, would make $19>15$); day 5 loads $10$. Total
        $\textit{daysNeeded}=5 \le 5$ -- exactly feasible, and capacity
        $14$ can be checked to need $6$ days, confirming $15$ is minimal).
\end{itemize}

The answer is $15$.

\subsection*{C++ implementation}

\begin{lstlisting}
#include <bits/stdc++.h>
using namespace std;

int daysNeeded(vector<int>& weights, int capacity) {
    int days = 1;
    long long currentLoad = 0;

    for (int w : weights) {
        if (currentLoad + w > capacity) {
            days++;
            currentLoad = 0;
        }
        currentLoad += w;
    }
    return days;
}

int shipWithinDays(vector<int>& weights, int D) {
    int lo = *max_element(weights.begin(), weights.end());
    int hi = accumulate(weights.begin(), weights.end(), 0);

    while (lo <= hi) {
        int mid = lo + (hi - lo) / 2;
        if (daysNeeded(weights, mid) <= D) {
            hi = mid - 1;
        } else {
            lo = mid + 1;
        }
    }
    return lo;
}

int main() {
    vector<int> weights = {1,2,3,4,5,6,7,8,9,10};
    cout << "Minimum capacity: "
         << shipWithinDays(weights, 5) << endl;
    return 0;
}
\end{lstlisting}

\begin{complexitybox}
\textbf{Time Complexity.} The candidate capacity range halves each
iteration, costing $O(\log(\sum(\textit{weights})))$ iterations, and each
feasibility check costs $O(n)$ for the greedy simulation, giving
\[
  O(n \log(\sum(\textit{weights}))).
\]
\textbf{Space Complexity.} $O(1)$ auxiliary space.
\end{complexitybox}

\begin{takeawaybox}
When the items being processed must stay in a fixed order, the
feasibility check for binary search on the answer is usually a greedy
simulation (pack as much as possible into each ``bucket'' before starting
a new one) rather than a simple per-item summation. Verify, separately,
that the greedy simulation itself is optimal for the fixed-candidate
sub-problem -- here, via the exchange argument that under-loading a day
never helps -- before trusting it as the feasibility predicate.
\end{takeawaybox}

\section{Kth Missing Positive Number}
\label{sec:bs-kth-missing}

\problemstatement{We are given a sorted array $\textit{arr}$ of $n$
distinct positive integers, and an integer $k$. We must find the $k$-th
positive integer that is \textbf{missing} from $\textit{arr}$.}

\begin{patternbox}
There is no array index to ``find $X$ at'' here -- we are counting
\emph{missing} numbers, which is a different flavor of monotonic quantity
than anything in Sections~\ref{sec:bs-find-x}--\ref{sec:bs-peak-1d}.
Whenever a sorted-array problem asks about \emph{missing} elements (gaps in
an otherwise-expected sequence) and wants the $k$-th such gap, look for a
formula that counts ``how many integers are missing up to this point in
the array,'' and binary search for the boundary where that count first
reaches $k$.
\end{patternbox}

\subsection*{Understanding the problem carefully}

Consider index $i$ (0-indexed) in the array. If none of $\textit{arr}[0],
\dots, \textit{arr}[i]$ were missing from the positive integers
$1, \dots, \textit{arr}[i]$, there would be exactly $\textit{arr}[i]$ of
them present, occupying $\textit{arr}[i]$ slots; but only $i+1$ of them
\emph{are} actually present (the first $i+1$ elements of the array), so
\[
  \textit{missingCount}(i) = \textit{arr}[i] - (i+1)
\]
counts exactly how many positive integers from $1$ to $\textit{arr}[i]$ are
missing from the array. Because $\textit{arr}$ is sorted and strictly
increasing, $\textit{arr}[i]$ grows at least as fast as $i$ does, so
$\textit{missingCount}(i)$ is non-decreasing in $i$ -- exactly the
monotonic structure we need.

\begin{ideabox}[Candidate Space and Predicate]
The candidate space is indices $0, \dots, n-1$ (with an implicit
``beyond the array'' fallback). The predicate is
$P(i) = (\textit{missingCount}(i) \ge k)$, and we want the smallest index
where $P$ holds.
\end{ideabox}

\subsection*{Why $k + \textit{lo}$ is the final answer}

Let $\textit{lo}$ be the smallest index where $\textit{missingCount}
(\textit{lo}) \ge k$ (or $n$, if even the last element's missing count is
still below $k$, meaning the answer lies beyond the array entirely). Just
before this index, $\textit{missingCount}(\textit{lo}-1) < k$, meaning
fewer than $k$ numbers are missing among $1, \dots, \textit{arr}
[\textit{lo}-1]$; from $\textit{arr}[\textit{lo}-1]+1$ onward (which is not
yet present in the array, since $\textit{arr}[\textit{lo}]$ is strictly
larger), missing numbers resume counting up consecutively until reaching
$\textit{missingCount}(\textit{lo}) \ge k$. A clean way to see the final
formula: exactly $\textit{lo}$ array elements (indices $0, \dots,
\textit{lo}-1$) are known to be \emph{smaller} than the answer (since the
crossover into ``enough missing numbers'' has not yet happened at index
$\textit{lo}-1$), so among the first $k + \textit{lo}$ positive integers,
exactly $\textit{lo}$ are present in the array and the remaining $k$ are
missing -- meaning the $k$-th missing positive integer is precisely
$k + \textit{lo}$.

\subsection*{Algorithm}

\begin{enumerate}
  \item Initialize $\textit{lo} \gets 0$, $\textit{hi} \gets n-1$.
  \item While $\textit{lo} \le \textit{hi}$:
  \begin{enumerate}
    \item $\textit{mid} \gets \textit{lo}+(\textit{hi}-\textit{lo})/2$.
    \item Compute $\textit{missingCount} \gets \textit{arr}[\textit{mid}] -
          (\textit{mid}+1)$.
    \item If $\textit{missingCount} < k$: $\textit{lo} \gets
          \textit{mid}+1$.
    \item Else: $\textit{hi} \gets \textit{mid}-1$.
  \end{enumerate}
  \item Return $k + \textit{lo}$.
\end{enumerate}

\begin{algorithm}[H]
\caption{Kth Missing Positive Number}
\begin{algorithmic}[1]
\Procedure{KthMissing}{$\textit{arr}, k$}
  \State $\textit{lo} \gets 0$, $\textit{hi} \gets n-1$
  \While{$\textit{lo} \le \textit{hi}$}
    \State $\textit{mid} \gets \textit{lo}+(\textit{hi}-\textit{lo})/2$
    \State $\textit{missing} \gets \textit{arr}[\textit{mid}] - (\textit{mid}+1)$
    \If{$\textit{missing} < k$}
      \State $\textit{lo} \gets \textit{mid} + 1$
    \Else
      \State $\textit{hi} \gets \textit{mid} - 1$
    \EndIf
  \EndWhile
  \State \Return $k + \textit{lo}$
\EndProcedure
\end{algorithmic}
\end{algorithm}

\subsection*{Dry run}

\dryrun{Take $\textit{arr} = [2,3,4,7,11]$, $k = 5$.}

\begin{itemize}
  \item $\textit{lo}=0,\textit{hi}=4$: $\textit{mid}=2$,
        $\textit{missing}=\textit{arr}[2]-3=4-3=1 < 5$. $\textit{lo}=3$.
  \item $\textit{lo}=3,\textit{hi}=4$: $\textit{mid}=3$,
        $\textit{missing}=7-4=3 < 5$. $\textit{lo}=4$.
  \item $\textit{lo}=4,\textit{hi}=4$: $\textit{mid}=4$,
        $\textit{missing}=11-5=6 \ge 5$. $\textit{hi}=3$.
  \item $\textit{lo}=4 > \textit{hi}=3$: loop ends.
\end{itemize}

The answer is $k+\textit{lo} = 5+4 = 9$. Checking by hand: the missing
positive integers, in order, are $1, 5, 6, 8, 9, 10, 12, \dots$, and the
$5$-th one is indeed $9$.

\subsection*{C++ implementation}

\begin{lstlisting}
#include <bits/stdc++.h>
using namespace std;

int findKthMissing(vector<int>& arr, int k) {
    int n = arr.size();
    int lo = 0, hi = n - 1;

    while (lo <= hi) {
        int mid = lo + (hi - lo) / 2;
        int missing = arr[mid] - (mid + 1);
        if (missing < k) {
            lo = mid + 1;
        } else {
            hi = mid - 1;
        }
    }
    return k + lo;
}

int main() {
    vector<int> arr = {2, 3, 4, 7, 11};
    cout << "5th missing positive number: "
         << findKthMissing(arr, 5) << endl;
    return 0;
}
\end{lstlisting}

\begin{complexitybox}
\textbf{Time Complexity.} The search window halves each iteration, giving
\[
  O(\log n).
\]
This is a substantial improvement over the natural $O(n)$ linear-scan
approach of walking through the array while counting missing numbers one
at a time.
\textbf{Space Complexity.} $O(1)$ auxiliary space.
\end{complexitybox}

\begin{takeawaybox}
When a problem asks about ``missing'' elements relative to an expected
dense sequence (here, the positive integers), look for a closed-form
counting function -- ``expected count minus actual count up to this
point'' -- that is monotonic, and binary search on that function's
crossover rather than scanning and counting one gap at a time.
\end{takeawaybox}

\section{Aggressive Cows Problem}
\label{sec:bs-aggressive-cows}

\problemstatement{We are given $n$ stall positions $\textit{stalls}[0..n-1]$
along a number line, and $k$ cows to place into these stalls (at most one
cow per stall). The cows are aggressive, so we want to place them to
\textbf{maximize the minimum distance} between any two cows. Find this
maximum possible minimum distance.}

\begin{patternbox}
\emph{``Maximize the minimum distance/gap between placed items''} is the
mirror image of every problem so far in this section: until now we
\emph{minimized} a candidate (speed, day, divisor, capacity) subject to a
feasibility constraint that became \emph{easier} as the candidate grew.
Here, feasibility (``can we place all $k$ cows with at least this much
spacing'') becomes \textbf{harder} as the candidate distance grows, so we
are finding the \textbf{largest} feasible candidate, not the smallest.
Whenever you see \emph{``maximize the minimum ...''}, expect to flip the
binary search's narrowing direction relative to the ``minimize subject to
feasibility'' problems seen earlier.
\end{patternbox}

\subsection*{Understanding the problem carefully}

Sort the stall positions first -- only relative order matters for spacing,
and an unsorted input gives no useful structure to exploit. For a fixed
candidate minimum distance $d$, ask: can we place all $k$ cows in the
sorted stalls such that every pair of adjacent placed cows is at least $d$
apart? This is checked greedily: place the first cow at the first stall,
then repeatedly place the next cow at the first stall that is at least $d$
away from the previously placed cow. If we manage to place all $k$ cows
this way, $d$ is feasible.

As $d$ increases, this greedy placement can only place fewer or equal
cows (every placement decision becomes at least as restrictive), so the
number of cows placeable is non-increasing in $d$ -- meaning feasibility
($\ge k$ cows placed) is \texttt{true} for small $d$ and \texttt{false}
for large $d$, the mirror image of the earlier sections' monotonicity
direction.

\begin{ideabox}[Candidate Space and Predicate]
The candidate space is distances $d \in [1, \max(\textit{stalls}) -
\min(\textit{stalls})]$. The predicate is
$P(d) = (\text{greedy placement with spacing} \ge d \text{ fits all } k
\text{ cows})$, and we want the \textbf{largest} $d$ where $P$ holds --
since $P$ is true for small $d$ and false for large $d$, we record a
successful candidate and narrow \emph{right} (search for an even larger
feasible $d$), the mirrored version of the ``narrow left after success''
template used for minimize-type problems.
\end{ideabox}

\subsection*{Why the greedy feasibility check is optimal}

For a \emph{fixed} target spacing $d$, greedily placing each next cow as
early as possible (at the first stall that is at least $d$ from the
previous placement) maximizes the number of cows we can fit using that
spacing: placing a cow any later than necessary only makes it harder to
fit subsequent cows, never easier, since stalls are processed in
increasing order and no stall can be revisited. This is the same kind of
exchange argument used for interval scheduling in the Greedy chapter --
the earliest valid placement is always at least as good as any later
valid placement for the purpose of fitting more items afterward.

\subsection*{Algorithm}

\begin{enumerate}
  \item Sort $\textit{stalls}$.
  \item Initialize $\textit{lo} \gets 1$,
        $\textit{hi} \gets \textit{stalls}[n-1] - \textit{stalls}[0]$.
  \item While $\textit{lo} \le \textit{hi}$:
  \begin{enumerate}
    \item $\textit{mid} \gets \textit{lo}+(\textit{hi}-\textit{lo})/2$.
    \item Greedily count how many cows fit with minimum spacing
          $\textit{mid}$.
    \item If $\textit{count} \ge k$: feasible; record \textit{mid}, try
          larger: $\textit{lo} \gets \textit{mid}+1$.
    \item Else: infeasible; try smaller: $\textit{hi} \gets \textit{mid}-1$.
  \end{enumerate}
  \item Return the recorded answer (equivalently \textit{hi} once the loop
        ends).
\end{enumerate}

\begin{algorithm}[H]
\caption{Aggressive Cows}
\begin{algorithmic}[1]
\Procedure{AggressiveCows}{\textit{stalls}, $k$}
  \State sort \textit{stalls}
  \State $\textit{lo} \gets 1$, $\textit{hi} \gets \textit{stalls}[n-1] - \textit{stalls}[0]$
  \While{$\textit{lo} \le \textit{hi}$}
    \State $\textit{mid} \gets \textit{lo}+(\textit{hi}-\textit{lo})/2$
    \State $\textit{count} \gets \Call{CowsPlaceable}{\textit{stalls}, \textit{mid}}$
    \If{$\textit{count} \ge k$}
      \State $\textit{lo} \gets \textit{mid} + 1$
    \Else
      \State $\textit{hi} \gets \textit{mid} - 1$
    \EndIf
  \EndWhile
  \State \Return \textit{hi}
\EndProcedure
\end{algorithmic}
\end{algorithm}

\subsection*{Dry run}

\dryrun{Take $\textit{stalls} = [0, 3, 4, 7, 10, 9]$, $k = 4$. Sorted:
$[0,3,4,7,9,10]$.}

\begin{itemize}
  \item $\textit{lo}=1,\textit{hi}=10$: $\textit{mid}=5$. Greedy: place at
        $0$; next $\ge 5$ is $7$, place; next $\ge 12$ -- none. Only $2$
        cows placed $< 4$. Infeasible. $\textit{hi}=4$.
  \item $\textit{lo}=1,\textit{hi}=4$: $\textit{mid}=2$. Greedy: place at
        $0$; next $\ge 2$ is $3$, place; next $\ge 5$ is $7$, place; next
        $\ge 9$ is $9$, place. $4$ cows placed $\ge 4$. Feasible.
        $\textit{lo}=3$.
  \item $\textit{lo}=3,\textit{hi}=4$: $\textit{mid}=3$. Greedy: place at
        $0$; next $\ge 3$ is $3$, place; next $\ge 6$ is $7$, place; next
        $\ge 10$ is $10$, place. $4$ cows placed $\ge 4$. Feasible.
        $\textit{lo}=4$.
  \item $\textit{lo}=4,\textit{hi}=4$: $\textit{mid}=4$. Greedy: place at
        $0$; next $\ge 4$ is $4$, place; next $\ge 8$ is $9$, place; next
        $\ge 13$ -- none. Only $3$ cows placed $< 4$. Infeasible.
        $\textit{hi}=3$.
  \item $\textit{lo}=4 > \textit{hi}=3$: loop ends.
\end{itemize}

The answer is $\textit{hi} = 3$: the maximum minimum distance is $3$.

\subsection*{C++ implementation}

\begin{lstlisting}
#include <bits/stdc++.h>
using namespace std;

int cowsPlaceable(vector<int>& stalls, int d) {
    int count = 1;
    int lastPos = stalls[0];

    for (int i = 1; i < (int)stalls.size(); i++) {
        if (stalls[i] - lastPos >= d) {
            count++;
            lastPos = stalls[i];
        }
    }
    return count;
}

int aggressiveCows(vector<int>& stalls, int k) {
    sort(stalls.begin(), stalls.end());
    int n = stalls.size();
    int lo = 1, hi = stalls[n - 1] - stalls[0];

    while (lo <= hi) {
        int mid = lo + (hi - lo) / 2;
        if (cowsPlaceable(stalls, mid) >= k) {
            lo = mid + 1;
        } else {
            hi = mid - 1;
        }
    }
    return hi;
}

int main() {
    vector<int> stalls = {0, 3, 4, 7, 10, 9};
    cout << "Maximum minimum distance: "
         << aggressiveCows(stalls, 4) << endl;
    return 0;
}
\end{lstlisting}

\begin{complexitybox}
\textbf{Time Complexity.} Sorting costs $O(n \log n)$. The candidate
distance range halves each iteration, costing
$O(\log(\max(\textit{stalls})-\min(\textit{stalls})))$ iterations, and
each feasibility check costs $O(n)$, giving an overall time complexity of
\[
  O(n \log n + n \log(\max(\textit{stalls}))).
\]
\textbf{Space Complexity.} $O(1)$ auxiliary space beyond sorting.
\end{complexitybox}

\begin{takeawaybox}
``Maximize the minimum'' and ``minimize the maximum'' problems both binary
search on the answer, but with the narrowing direction flipped relative to
each other: for ``maximize the minimum,'' record a successful candidate
and search for a \emph{larger} one; for ``minimize the maximum'' (as in
Sections~\ref{sec:bs-koko}--\ref{sec:bs-ship-packages}), record a
successful candidate and search for a \emph{smaller} one. Always pause to
identify which of the two flavors you are facing before writing the
narrowing logic.
\end{takeawaybox}

\section{Book Allocation Problem}
\label{sec:bs-book-allocation}

\problemstatement{We are given $n$ books, where book $i$ has
$\textit{pages}[i]$ pages, and $m$ students. Each student must be allocated
a \textbf{contiguous} range of books (in the given order), every book must
be allocated to exactly one student, and every student must receive at
least one book. We must choose the allocation to \textbf{minimize the
maximum} number of pages assigned to any single student.}

\begin{patternbox}
\emph{``Minimize the maximum load, items assigned in contiguous order,
divided among a fixed number of workers/students''} is exactly the same
shape as Capacity to Ship Packages within D Days
(Section~\ref{sec:bs-ship-packages}): there, we fixed the number of
``buckets'' indirectly through a day limit $D$ and a per-bucket capacity
candidate; here, the number of buckets is fixed directly as $m$ students.
Once you notice that ``contiguous partition into a fixed number of groups,
minimize the maximum group sum'' is the same template regardless of
whether it is phrased as books, packages, or paint sections (see
Section~\ref{sec:bs-painters-partition}), you can reuse the identical
greedy feasibility check every time.
\end{patternbox}

\subsection*{Understanding the problem carefully}

For a fixed candidate maximum $\textit{maxPages}$, greedily simulate
assigning books to the current student until adding the next book's pages
would exceed $\textit{maxPages}$; at that point, move to the next student.
This yields $\textit{studentsNeeded}(\textit{maxPages})$, the minimum
number of students required if no student may hold more than
$\textit{maxPages}$ total pages. As $\textit{maxPages}$ increases, each
student can absorb more books before being forced to stop, so
$\textit{studentsNeeded}$ is non-increasing in $\textit{maxPages}$.

\begin{ideabox}[Candidate Space and Predicate]
The candidate space is integer page totals in
$[\max(\textit{pages}), \sum(\textit{pages})]$ (a single student must be
able to hold at least the largest individual book, and a single student
holding everything always works). The predicate is
$P(\textit{maxPages}) = (\textit{studentsNeeded}(\textit{maxPages}) \le
m)$, and we want the smallest $\textit{maxPages}$ where $P$ holds -- this
is a ``minimize subject to feasibility'' problem, narrowing the same
direction as Sections~\ref{sec:bs-koko}--\ref{sec:bs-ship-packages}, not
the flipped direction of Aggressive Cows.
\end{ideabox}

\subsection*{Why this is identical in structure to ship packages}

The greedy feasibility check here is justified by precisely the same
exchange argument as Section~\ref{sec:bs-ship-packages}: greedily filling
each student with as many books as the candidate maximum allows, before
moving to the next student, never uses more students than any other valid
contiguous allocation respecting that same maximum, because deferring a
book to a later student only ever makes that later student's job harder,
never easier, and books cannot be reordered.

\subsection*{Algorithm}

\begin{enumerate}
  \item Initialize $\textit{lo} \gets \max(\textit{pages})$,
        $\textit{hi} \gets \sum(\textit{pages})$.
  \item While $\textit{lo} \le \textit{hi}$:
  \begin{enumerate}
    \item $\textit{mid} \gets \textit{lo}+(\textit{hi}-\textit{lo})/2$.
    \item Compute $\textit{studentsNeeded}(\textit{mid})$ via greedy
          simulation.
    \item If $\textit{studentsNeeded}(\textit{mid}) \le m$: feasible; try
          smaller: $\textit{hi} \gets \textit{mid}-1$.
    \item Else: infeasible; try larger: $\textit{lo} \gets \textit{mid}+1$.
  \end{enumerate}
  \item Return \textit{lo}.
\end{enumerate}

\begin{algorithm}[H]
\caption{Book Allocation}
\begin{algorithmic}[1]
\Procedure{BookAllocation}{\textit{pages}, $m$}
  \State $\textit{lo} \gets \max(\textit{pages})$, $\textit{hi} \gets \sum(\textit{pages})$
  \While{$\textit{lo} \le \textit{hi}$}
    \State $\textit{mid} \gets \textit{lo}+(\textit{hi}-\textit{lo})/2$
    \State $\textit{students} \gets \Call{StudentsNeeded}{\textit{pages}, \textit{mid}}$
    \If{$\textit{students} \le m$}
      \State $\textit{hi} \gets \textit{mid} - 1$
    \Else
      \State $\textit{lo} \gets \textit{mid} + 1$
    \EndIf
  \EndWhile
  \State \Return \textit{lo}
\EndProcedure
\end{algorithmic}
\end{algorithm}

\subsection*{Dry run}

\dryrun{Take $\textit{pages} = [12, 34, 67, 90]$, $m = 2$.}

$\max = 90$, $\sum = 203$.

\begin{itemize}
  \item $\textit{lo}=90,\textit{hi}=203$: $\textit{mid}=146$. Greedy:
        student 1 takes $12{+}34{+}67=113$ (next, $90$, would make
        $203>146$, stop); student 2 takes $90$. $\textit{studentsNeeded}=2
        \le 2$. Feasible. $\textit{hi}=145$.
  \item Narrowing continues; eventually the search converges to
        $\textit{lo}=113$: with maximum $113$, student 1 takes
        $12{+}34{+}67=113$ (next, $90$, would make $203>113$, stop);
        student 2 takes $90$. Still $2$ students. Checking $112$: student
        1 can only take $12{+}34=46$ (next, $67$, would make $113>112$,
        stop); student 2 takes $67$ (next, $90$, would make $157>112$,
        stop); student 3 takes $90$ -- that needs $3$ students, exceeding
        $m=2$. So $113$ is indeed minimal.
\end{itemize}

The answer is $113$.

\subsection*{C++ implementation}

\begin{lstlisting}
#include <bits/stdc++.h>
using namespace std;

int studentsNeeded(vector<int>& pages, int maxPages) {
    int students = 1;
    long long currentLoad = 0;

    for (int p : pages) {
        if (currentLoad + p > maxPages) {
            students++;
            currentLoad = 0;
        }
        currentLoad += p;
    }
    return students;
}

int bookAllocation(vector<int>& pages, int m) {
    int lo = *max_element(pages.begin(), pages.end());
    int hi = accumulate(pages.begin(), pages.end(), 0);

    while (lo <= hi) {
        int mid = lo + (hi - lo) / 2;
        if (studentsNeeded(pages, mid) <= m) {
            hi = mid - 1;
        } else {
            lo = mid + 1;
        }
    }
    return lo;
}

int main() {
    vector<int> pages = {12, 34, 67, 90};
    cout << "Minimum maximum pages: "
         << bookAllocation(pages, 2) << endl;
    return 0;
}
\end{lstlisting}

\begin{complexitybox}
\textbf{Time Complexity.}
\[
  O(n \log(\sum(\textit{pages}))),
\]
identical in structure to Section~\ref{sec:bs-ship-packages}.
\textbf{Space Complexity.} $O(1)$ auxiliary space.
\end{complexitybox}

\begin{takeawaybox}
``Partition a sequence into a fixed number of contiguous groups to
minimize the maximum group sum'' is a single template with at least three
names on this sheet: Book Allocation, Capacity to Ship Packages, and
Painter's Partition (Section~\ref{sec:bs-painters-partition}). Learn the
template once -- greedy bucket-filling feasibility check nested inside a
minimize-subject-to-feasibility binary search -- and recognize its
disguises rather than re-deriving it from scratch each time.
\end{takeawaybox}

\section{Split Array -- Largest Sum}
\label{sec:bs-split-array}

\problemstatement{We are given an array $\textit{nums}$ of $n$ non-negative
integers and an integer $m$. We must split $\textit{nums}$ into $m$
\textbf{contiguous, non-empty} subarrays so as to \textbf{minimize the
largest sum} among the $m$ subarrays.}

\begin{patternbox}
This is, word for word, the Book Allocation Problem
(Section~\ref{sec:bs-book-allocation}) with ``books/pages/students''
relabeled as ``array/elements/subarrays.'' There is no new idea to extract
here beyond confirming, explicitly, that you recognize the relabeling --
which is itself the point of including this problem on the sheet
immediately after book allocation: to test whether the pattern has truly
been internalized, or only memorized under one specific cover story.
\end{patternbox}

\subsection*{Understanding the problem carefully}

Exactly as in Section~\ref{sec:bs-book-allocation}, define
$\textit{partsNeeded}(\textit{maxSum})$ as the minimum number of
contiguous parts needed if no part's sum may exceed $\textit{maxSum}$,
computed by the same greedy bucket-filling simulation. This is
non-increasing in $\textit{maxSum}$, giving the same monotonic structure.

\begin{ideabox}[Candidate Space and Predicate]
Candidate space: integer sums in $[\max(\textit{nums}),
\sum(\textit{nums})]$. Predicate:
$P(\textit{maxSum}) = (\textit{partsNeeded}(\textit{maxSum}) \le m)$. Want
the smallest $\textit{maxSum}$ where $P$ holds.
\end{ideabox}

\subsection*{Algorithm}

Identical to Algorithm~\ref{sec:bs-book-allocation}'s
\textsc{BookAllocation}, with \textit{pages} relabeled \textit{nums} and
$m$ meaning the number of parts directly.

\subsection*{Dry run}

\dryrun{Take $\textit{nums} = [7,2,5,10,8]$, $m = 2$.}

$\max=10$, $\sum=32$. Binary search converges to $\textit{maxSum}=18$: one
valid split is $[7,2,5]$ (sum $14$) and $[10,8]$ (sum $18$), giving largest
sum $18$; checking $17$ shows it forces a third part
($[7,2,5]=14$, then $[10]=10$ would leave $[8]$ separately, or
$[7,2,5,10]=24>17$ forces an earlier break, ultimately requiring $3$
parts), confirming $18$ is minimal.

\subsection*{C++ implementation}

\begin{lstlisting}
#include <bits/stdc++.h>
using namespace std;

int partsNeeded(vector<int>& nums, int maxSum) {
    int parts = 1;
    long long currentSum = 0;

    for (int x : nums) {
        if (currentSum + x > maxSum) {
            parts++;
            currentSum = 0;
        }
        currentSum += x;
    }
    return parts;
}

int splitArray(vector<int>& nums, int m) {
    int lo = *max_element(nums.begin(), nums.end());
    int hi = accumulate(nums.begin(), nums.end(), 0);

    while (lo <= hi) {
        int mid = lo + (hi - lo) / 2;
        if (partsNeeded(nums, mid) <= m) {
            hi = mid - 1;
        } else {
            lo = mid + 1;
        }
    }
    return lo;
}

int main() {
    vector<int> nums = {7, 2, 5, 10, 8};
    cout << "Minimized largest sum: "
         << splitArray(nums, 2) << endl;
    return 0;
}
\end{lstlisting}

\begin{complexitybox}
\textbf{Time Complexity.}
$O(n \log(\sum(\textit{nums})))$, identical to
Section~\ref{sec:bs-book-allocation}.
\textbf{Space Complexity.} $O(1)$ auxiliary space.
\end{complexitybox}

\begin{takeawaybox}
The fastest way to solve this problem is to recognize it as Book
Allocation in disguise. Building a small catalogue of templates -- rather
than a catalogue of individual problems -- is what lets binary search on
the answer scale to dozens of seemingly distinct sheet problems with only
a handful of underlying ideas.
\end{takeawaybox}

\section{Painter's Partition Problem}
\label{sec:bs-painters-partition}

\problemstatement{We are given $n$ boards with lengths $\textit{boards}[0..n-1]$
and $k$ painters. Each painter paints a \textbf{contiguous} set of boards,
and painting one unit of board length takes one unit of time (painters
work in parallel, each on their own assigned contiguous range). We must
assign boards to painters so as to \textbf{minimize the time the slowest
painter takes} (equivalently, minimize the maximum total board length
assigned to any one painter).}

\begin{patternbox}
A third disguise of the same template as
Sections~\ref{sec:bs-book-allocation}--\ref{sec:bs-split-array}:
``boards/painters'' plays the role of ``pages/students'' or
``array elements/subarray parts.'' By this third repetition, the goal is
no longer to re-derive anything, but to confirm that the pattern --
\emph{partition a sequence into a fixed number of contiguous groups,
minimize the maximum group sum, via binary search on the answer with a
greedy bucket-filling feasibility check} -- has become an immediately
recognizable, reusable unit, regardless of the cover story a problem uses
to disguise it.
\end{patternbox}

\subsection*{Understanding the problem carefully}

Define $\textit{paintersNeeded}(\textit{maxTime})$ exactly as in
Section~\ref{sec:bs-book-allocation}: greedily assign boards to the
current painter until adding the next board would exceed
$\textit{maxTime}$, then move to the next painter. This count is
non-increasing in $\textit{maxTime}$, by the same reasoning as the prior
two sections.

\begin{ideabox}[Candidate Space and Predicate]
Candidate space: integer times in $[\max(\textit{boards}),
\sum(\textit{boards})]$. Predicate:
$P(\textit{maxTime}) = (\textit{paintersNeeded}(\textit{maxTime}) \le k)$.
Want the smallest $\textit{maxTime}$ where $P$ holds.
\end{ideabox}

\subsection*{Algorithm}

Identical to Algorithm~\ref{sec:bs-book-allocation}'s
\textsc{BookAllocation}, with \textit{pages} relabeled \textit{boards} and
$m$ relabeled $k$.

\subsection*{Dry run}

\dryrun{Take $\textit{boards} = [10, 20, 30, 40]$, $k = 2$.}

$\max=40$, $\sum=100$. Binary search converges to $\textit{maxTime}=60$:
one painter takes $[10,20,30]$ (sum $60$), the other takes $[40]$ (sum
$40$); checking $59$ would force the first painter to stop after
$[10,20]=30$ (since adding $30$ would make $60>59$), needing a third
painter for $[30]$ and $[40]$ separately exceeding $k=2$, so $60$ is
minimal.

\subsection*{C++ implementation}

\begin{lstlisting}
#include <bits/stdc++.h>
using namespace std;

int paintersNeeded(vector<int>& boards, int maxTime) {
    int painters = 1;
    long long currentTime = 0;

    for (int len : boards) {
        if (currentTime + len > maxTime) {
            painters++;
            currentTime = 0;
        }
        currentTime += len;
    }
    return painters;
}

int paintersPartition(vector<int>& boards, int k) {
    int lo = *max_element(boards.begin(), boards.end());
    int hi = accumulate(boards.begin(), boards.end(), 0);

    while (lo <= hi) {
        int mid = lo + (hi - lo) / 2;
        if (paintersNeeded(boards, mid) <= k) {
            hi = mid - 1;
        } else {
            lo = mid + 1;
        }
    }
    return lo;
}

int main() {
    vector<int> boards = {10, 20, 30, 40};
    cout << "Minimum time: " << paintersPartition(boards, 2) << endl;
    return 0;
}
\end{lstlisting}

\begin{complexitybox}
\textbf{Time Complexity.}
$O(n \log(\sum(\textit{boards})))$.
\textbf{Space Complexity.} $O(1)$ auxiliary space.
\end{complexitybox}

\begin{takeawaybox}
Three problems, three cover stories, one template. If you find yourself
about to write a fresh derivation for a ``partition into $k$ contiguous
groups, minimize the maximum group sum'' problem, stop and check the
catalogue of Sections~\ref{sec:bs-book-allocation}--\ref{sec:bs-painters-partition}
first.
\end{takeawaybox}

\section{Minimize the Maximum Distance Between Gas Stations}
\label{sec:bs-gas-stations}

\problemstatement{We are given $n$ gas station positions, sorted along a
number line, and an integer $k$. We may insert $k$ additional gas stations
\textbf{anywhere} (not restricted to integer positions) among the existing
ones. We must choose where to insert them so as to \textbf{minimize the
largest gap} between any two consecutive gas stations (existing or newly
inserted), and report that minimized largest gap.}

\begin{patternbox}
This problem keeps the ``minimize the maximum, subject to a budget''
shape of Sections~\ref{sec:bs-book-allocation}--\ref{sec:bs-painters-partition},
but with one crucial new wrinkle: the answer is a \textbf{real number},
not an integer, because stations may be placed at fractional positions.
Whenever a binary-search-on-the-answer problem allows continuous
(real-valued) placement rather than discrete/integer choices, the
\texttt{while(lo <= hi)} integer loop must be replaced with a
fixed-precision loop that runs until \textit{hi} and \textit{lo} are
within some small tolerance $\varepsilon$ of each other.
\end{patternbox}

\subsection*{Understanding the problem carefully}

Consider the existing gaps between consecutive original stations:
$\textit{gap}_1, \textit{gap}_2, \dots, \textit{gap}_{n-1}$. For a fixed
candidate maximum allowed gap $d$, the minimum number of new stations
needed to bring gap $\textit{gap}_i$ down to segments no larger than $d$
is
\[
  \left\lceil \frac{\textit{gap}_i}{d} \right\rceil - 1
\]
(inserting $c$ points into a single gap of length $L$ divides it into
$c+1$ equal-or-smaller segments; we need the smallest $c$ such that
$L/(c+1) \le d$, which rearranges to $c \ge L/d - 1$, i.e.\ $c = \lceil
L/d \rceil - 1$ for integer $c$). Summing this over every existing gap
gives $\textit{stationsNeeded}(d)$, the total new stations required to
guarantee no gap exceeds $d$. As $d$ increases, fewer stations are needed
per gap, so $\textit{stationsNeeded}(d)$ is non-increasing in $d$.

\begin{ideabox}[Candidate Space and Predicate]
The candidate space is real numbers $d \in (0, \max(\textit{gap}_i)]$.
The predicate is $P(d) = (\textit{stationsNeeded}(d) \le k)$, and we want
the smallest $d$ where $P$ holds.
\end{ideabox}

\subsection*{Why a fixed-precision loop replaces the integer loop}

With a continuous candidate space, there is no notion of ``the next
integer candidate,'' so the loop cannot terminate by \textit{lo} and
\textit{hi} crossing as integers do. Instead, we run the loop a fixed
number of times (or until $\textit{hi} - \textit{lo}$ shrinks below a
chosen precision threshold $\varepsilon$, e.g.\ $10^{-6}$), narrowing
\textit{lo} or \textit{hi} to \textit{mid} (not $\textit{mid} \pm 1$, since
there is no discrete step to add or subtract) at each iteration. Each
iteration still halves the remaining interval, so a fixed number of
iterations (around $50$--$100$, since $2^{-100}$ is far smaller than any
reasonable $\varepsilon$) is always enough regardless of the initial
interval width.

\subsection*{Algorithm}

\begin{enumerate}
  \item Compute the existing gaps; $\textit{lo} \gets 0$,
        $\textit{hi} \gets \max(\textit{gap}_i)$.
  \item Repeat a fixed number of times (or while
        $\textit{hi} - \textit{lo} > \varepsilon$):
  \begin{enumerate}
    \item $\textit{mid} \gets (\textit{lo}+\textit{hi})/2$.
    \item Compute $\textit{stationsNeeded}(\textit{mid}) =
          \sum_i \left( \lceil \textit{gap}_i / \textit{mid} \rceil - 1
          \right)$.
    \item If $\textit{stationsNeeded}(\textit{mid}) \le k$:
          $\textit{hi} \gets \textit{mid}$ (feasible; try smaller).
    \item Else: $\textit{lo} \gets \textit{mid}$ (infeasible; try larger).
  \end{enumerate}
  \item Return $\textit{hi}$ (or $\textit{lo}$; they have converged to
        within $\varepsilon$).
\end{enumerate}

\begin{algorithm}[H]
\caption{Minimize Maximum Gas Station Distance}
\begin{algorithmic}[1]
\Procedure{MinimizeMaxDistance}{\textit{stations}, $k$}
  \State compute gaps $\textit{gap}_1, \dots, \textit{gap}_{n-1}$
  \State $\textit{lo} \gets 0$, $\textit{hi} \gets \max_i \textit{gap}_i$
  \For{$\textit{iter} \gets 1$ to $100$}
    \State $\textit{mid} \gets (\textit{lo}+\textit{hi})/2$
    \State $\textit{needed} \gets \sum_i (\lceil \textit{gap}_i / \textit{mid} \rceil - 1)$
    \If{$\textit{needed} \le k$}
      \State $\textit{hi} \gets \textit{mid}$
    \Else
      \State $\textit{lo} \gets \textit{mid}$
    \EndIf
  \EndFor
  \State \Return \textit{hi}
\EndProcedure
\end{algorithmic}
\end{algorithm}

\subsection*{Dry run}

\dryrun{Take stations at $[1, 2, 3, 4, 5]$ (gaps all $=1$), $k = 4$.}

There are $4$ gaps, each of length $1$. With $k=4$ new stations to
distribute (one per gap), each gap of length $1$ becomes two segments of
length $0.5$. So $\textit{stationsNeeded}(0.5) = \sum (\lceil 1/0.5
\rceil - 1) = \sum (2-1) = 4 \cdot 1 = 4 \le 4$ -- feasible. Checking a
slightly smaller candidate, say $d=0.49$:
$\lceil 1/0.49 \rceil - 1 = \lceil 2.04 \rceil - 1 = 3-1=2$ per gap, times
$4$ gaps $=8 > 4$ -- infeasible. So the binary search converges to
$d = 0.5$, the minimized maximum distance.

\subsection*{C++ implementation}

\begin{lstlisting}
#include <bits/stdc++.h>
using namespace std;

double minimizeMaxDistance(vector<int>& stations, int k) {
    int n = stations.size();
    vector<double> gaps(n - 1);
    for (int i = 0; i + 1 < n; i++) {
        gaps[i] = stations[i + 1] - stations[i];
    }

    double lo = 0, hi = *max_element(gaps.begin(), gaps.end());

    for (int iter = 0; iter < 100; iter++) {
        double mid = (lo + hi) / 2.0;
        long long needed = 0;
        for (double g : gaps) {
            needed += (long long)ceil(g / mid) - 1;
        }

        if (needed <= k) {
            hi = mid;
        } else {
            lo = mid;
        }
    }
    return hi;
}

int main() {
    vector<int> stations = {1, 2, 3, 4, 5};
    cout << fixed << setprecision(4);
    cout << "Minimized maximum distance: "
         << minimizeMaxDistance(stations, 4) << endl;
    return 0;
}
\end{lstlisting}

\begin{complexitybox}
\textbf{Time Complexity.} The fixed-precision loop runs a constant number
of iterations (e.g.\ $100$), and each feasibility check costs $O(n)$,
giving
\[
  O(n \cdot \text{iterations}) = O(n),
\]
treating the iteration count as a constant determined by the desired
precision (more precisely, $O(n \log(\textit{range}/\varepsilon))$ if the
precision $\varepsilon$ is treated as a variable rather than fixed).
\textbf{Space Complexity.} $O(n)$ to store the gaps.
\end{complexitybox}

\begin{takeawaybox}
Binary search on the answer is not limited to integer candidates. When the
problem allows real-valued placement, keep every other part of the
template (monotonic feasibility predicate, minimize/maximize narrowing
direction) and simply replace the integer loop and its $\pm 1$ adjustments
with a fixed-iteration or epsilon-tolerance loop that narrows
\textit{lo}/\textit{hi} directly to \textit{mid}.
\end{takeawaybox}

\section{Median of Two Sorted Arrays}
\label{sec:bs-median-two-arrays}

\problemstatement{We are given two sorted arrays $\textit{nums1}$ (size
$n_1$) and $\textit{nums2}$ (size $n_2$). We must find the median of the
combined sorted array of all $n_1 + n_2$ elements, \textbf{without}
actually merging the two arrays, in $O(\log(\min(n_1, n_2)))$ time.}

\begin{patternbox}
The keyword combination \emph{``median of two sorted arrays''} together
with a required complexity of $O(\log(\min(n_1,n_2)))$ -- far faster than
the $O(n_1+n_2)$ that an actual merge would cost -- is the signal for a
fundamentally different binary-search target than anything earlier in
this chapter: we will binary search not over values, and not over a
single array's indices, but over \textbf{how the combined sorted sequence
should be partitioned} between the two arrays. This ``binary search on a
partition point'' technique is one of the most intricate patterns on the
sheet, and is worth slowing down for.
\end{patternbox}

\subsection*{Understanding the problem carefully}

The median of a sorted sequence of length $T$ splits it into a left half
of size $\lceil T/2 \rceil$ and a right half of size $\lfloor T/2
\rfloor$, where every element in the left half is $\le$ every element in
the right half. Our goal is to find \emph{how many} elements of this left
half come from $\textit{nums1}$ (call this $\textit{cut1}$) and how many
come from $\textit{nums2}$ (call this $\textit{cut2} = \lceil T/2 \rceil -
\textit{cut1}$), \emph{without} ever materializing the merged array.

For a candidate $\textit{cut1} \in [0, n_1]$ (with $\textit{cut2}$
determined as above), define four boundary values:
\[
  l_1 = \textit{nums1}[\textit{cut1}-1] \ (\text{or } -\infty \text{ if }
  \textit{cut1}=0), \qquad
  r_1 = \textit{nums1}[\textit{cut1}] \ (\text{or } +\infty \text{ if }
  \textit{cut1}=n_1),
\]
\[
  l_2 = \textit{nums2}[\textit{cut2}-1] \ (\text{or } -\infty \text{ if }
  \textit{cut2}=0), \qquad
  r_2 = \textit{nums2}[\textit{cut2}] \ (\text{or } +\infty \text{ if }
  \textit{cut2}=n_2).
\]
A partition is \textbf{valid} exactly when $l_1 \le r_2$ \emph{and}
$l_2 \le r_1$: every element on the left side ($l_1$ and $l_2$) is at most
every element on the right side ($r_1$ and $r_2$), which is precisely the
condition needed for the left half (across both arrays combined) to
contain only values $\le$ every value in the right half.

\begin{ideabox}[Candidate Space and Predicate]
The candidate space is $\textit{cut1} \in [0, n_1]$ (binary searching on
the \emph{smaller} array keeps the range, and hence the iteration count,
as small as possible). The predicate is
$P(\textit{cut1}) = (l_1 \le r_2 \text{ and } l_2 \le r_1)$. Because
increasing $\textit{cut1}$ monotonically increases $l_1$ and $r_1$ while
monotonically decreasing $l_2$ and $r_2$ (as $\textit{cut2}$ shrinks
correspondingly), there is exactly one crossover region where the
partition is valid, and binary search finds it directly.
\end{ideabox}

\subsection*{Why exactly one valid partition exists, and how to search for it}

If $l_1 > r_2$, the left part of $\textit{nums1}$ is reaching too far
(taking values larger than what the right part of $\textit{nums2}$ can
absorb), so $\textit{cut1}$ must decrease: search left,
$\textit{hi} \gets \textit{cut1}-1$. If $l_2 > r_1$ (equivalently,
$\textit{nums1}$'s left part is not reaching far enough), $\textit{cut1}$
must increase: search right, $\textit{lo} \gets \textit{cut1}+1$. Because
$l_1, r_1$ increase weakly with $\textit{cut1}$ and $l_2, r_2$ decrease
weakly with $\textit{cut1}$ (equivalently increase weakly as
$\textit{cut1}$ decreases), exactly one of these two failure conditions
can hold at a time for an invalid partition, and the search converges to
the unique boundary where neither failure condition holds.

Once a valid partition is found, the median is computed from
$\max(l_1, l_2)$ and $\min(r_1, r_2)$: if $T$ is odd, the median is
$\max(l_1, l_2)$ (the largest element of the left half, which has one more
element than the right half); if $T$ is even, the median is the average of
$\max(l_1, l_2)$ and $\min(r_1, r_2)$.

\subsection*{Algorithm}

\begin{enumerate}
  \item Ensure $\textit{nums1}$ is the smaller array (swap if necessary,
        so the binary search range is $[0, n_1]$ with $n_1 \le n_2$).
  \item $T \gets n_1+n_2$, $\textit{half} \gets \lceil T/2 \rceil$.
  \item $\textit{lo} \gets 0$, $\textit{hi} \gets n_1$.
  \item While $\textit{lo} \le \textit{hi}$:
  \begin{enumerate}
    \item $\textit{cut1} \gets \textit{lo}+(\textit{hi}-\textit{lo})/2$,
          $\textit{cut2} \gets \textit{half}-\textit{cut1}$.
    \item Compute $l_1, r_1, l_2, r_2$ as above.
    \item If $l_1 \le r_2$ and $l_2 \le r_1$: valid partition found;
          compute and return the median as described above.
    \item Else if $l_1 > r_2$: $\textit{hi} \gets \textit{cut1}-1$.
    \item Else: $\textit{lo} \gets \textit{cut1}+1$.
  \end{enumerate}
\end{enumerate}

\begin{algorithm}[H]
\caption{Median of Two Sorted Arrays}
\begin{algorithmic}[1]
\Procedure{FindMedian}{$\textit{nums1}, \textit{nums2}$}
  \If{$n_1 > n_2$} \textbf{swap} $\textit{nums1}, \textit{nums2}$ \EndIf
  \State $T \gets n_1 + n_2$, $\textit{half} \gets \lceil T/2 \rceil$
  \State $\textit{lo} \gets 0$, $\textit{hi} \gets n_1$
  \While{$\textit{lo} \le \textit{hi}$}
    \State $\textit{cut1} \gets \textit{lo}+(\textit{hi}-\textit{lo})/2$
    \State $\textit{cut2} \gets \textit{half} - \textit{cut1}$
    \State $l_1 \gets (\textit{cut1}=0)\ ? -\infty : \textit{nums1}[\textit{cut1}-1]$
    \State $r_1 \gets (\textit{cut1}=n_1)\ ? +\infty : \textit{nums1}[\textit{cut1}]$
    \State $l_2 \gets (\textit{cut2}=0)\ ? -\infty : \textit{nums2}[\textit{cut2}-1]$
    \State $r_2 \gets (\textit{cut2}=n_2)\ ? +\infty : \textit{nums2}[\textit{cut2}]$
    \If{$l_1 \le r_2$ \textbf{and} $l_2 \le r_1$}
      \If{$T$ is odd} \Return $\max(l_1, l_2)$
      \Else \Return $(\max(l_1,l_2)+\min(r_1,r_2))/2$ \EndIf
    \ElsIf{$l_1 > r_2$}
      \State $\textit{hi} \gets \textit{cut1} - 1$
    \Else
      \State $\textit{lo} \gets \textit{cut1} + 1$
    \EndIf
  \EndWhile
\EndProcedure
\end{algorithmic}
\end{algorithm}

\subsection*{Dry run}

\dryrun{Take $\textit{nums1} = [1,3]$, $\textit{nums2} = [2]$.}

$n_1=2 \le n_2=1$? No -- swap so $\textit{nums1}=[2]$ ($n_1=1$),
$\textit{nums2}=[1,3]$ ($n_2=2$). $T=3$, $\textit{half}=\lceil 3/2
\rceil=2$.

\begin{itemize}
  \item $\textit{lo}=0,\textit{hi}=1$: $\textit{cut1}=0$,
        $\textit{cut2}=2-0=2$. $l_1=-\infty$ ($\textit{cut1}=0$),
        $r_1=\textit{nums1}[0]=2$. $l_2=\textit{nums2}[1]=3$
        ($\textit{cut2}=2=n_2$, so $r_2=+\infty$). Check: $l_1=-\infty \le
        r_2=+\infty$? Yes. $l_2=3 \le r_1=2$? No. So $l_2 > r_1$,
        $\textit{lo} \gets 1$.
  \item $\textit{lo}=1,\textit{hi}=1$: $\textit{cut1}=1$,
        $\textit{cut2}=2-1=1$. $l_1=\textit{nums1}[0]=2$
        ($\textit{cut1}=1=n_1$, so $r_1=+\infty$). $l_2=\textit{nums2}[0]=1$,
        $r_2=\textit{nums2}[1]=3$. Check: $l_1=2 \le r_2=3$? Yes.
        $l_2=1 \le r_1=+\infty$? Yes. Valid partition.
\end{itemize}

$T=3$ is odd, so the median is $\max(l_1,l_2) = \max(2,1) = 2$. Checking
by hand: the merged array is $[1,2,3]$, whose median is indeed $2$.

\subsection*{C++ implementation}

\begin{lstlisting}
#include <bits/stdc++.h>
using namespace std;

double findMedianSortedArrays(vector<int>& nums1, vector<int>& nums2) {
    if (nums1.size() > nums2.size()) swap(nums1, nums2);

    int n1 = nums1.size(), n2 = nums2.size();
    int total = n1 + n2;
    int half = (total + 1) / 2;

    int lo = 0, hi = n1;

    while (lo <= hi) {
        int cut1 = lo + (hi - lo) / 2;
        int cut2 = half - cut1;

        int l1 = (cut1 == 0) ? INT_MIN : nums1[cut1 - 1];
        int r1 = (cut1 == n1) ? INT_MAX : nums1[cut1];
        int l2 = (cut2 == 0) ? INT_MIN : nums2[cut2 - 1];
        int r2 = (cut2 == n2) ? INT_MAX : nums2[cut2];

        if (l1 <= r2 && l2 <= r1) {
            if (total % 2 == 1) {
                return max(l1, l2);
            } else {
                return (max(l1, l2) + min(r1, r2)) / 2.0;
            }
        } else if (l1 > r2) {
            hi = cut1 - 1;
        } else {
            lo = cut1 + 1;
        }
    }
    return -1.0; // unreachable for valid sorted input
}

int main() {
    vector<int> nums1 = {1, 3};
    vector<int> nums2 = {2};
    cout << "Median: "
         << findMedianSortedArrays(nums1, nums2) << endl;
    return 0;
}
\end{lstlisting}

\begin{complexitybox}
\textbf{Time Complexity.} The binary search runs over $\textit{cut1} \in
[0, n_1]$, where $n_1 = \min(n_1, n_2)$ (after the swap), giving
\[
  O(\log(\min(n_1, n_2))).
\]
\textbf{Space Complexity.} $O(1)$ auxiliary space.
\end{complexitybox}

\begin{takeawaybox}
Binary search does not need to search over array \emph{values} or even a
single array's \emph{indices} -- here, it searches over a \textbf{split
point} that simultaneously partitions two arrays, using a validity
condition ($l_1 \le r_2$ and $l_2 \le r_1$) as the monotonic predicate.
This ``binary search on a partition'' idea, once recognized, also unlocks
Section~\ref{sec:bs-kth-element}, the general $k$-th element version of
this same problem.
\end{takeawaybox}

\section{Kth Element of Two Sorted Arrays}
\label{sec:bs-kth-element}

\problemstatement{We are given two sorted arrays $\textit{nums1}$ (size
$n_1$) and $\textit{nums2}$ (size $n_2$), and an integer $k$. We must find
the $k$-th smallest element (1-indexed) of the combined sorted array,
without merging, in $O(\log(\min(n_1,n_2)))$ time.}

\begin{patternbox}
This is the direct generalization of
Section~\ref{sec:bs-median-two-arrays}: the median is simply the special
case $k = \lceil (n_1+n_2)/2 \rceil$ (for odd total length) or the average
of the $k$-th and $(k{+}1)$-th elements (for even total length). The same
``binary search on a partition split between two arrays'' technique
applies, with the partition sized to put exactly $k$ elements on the left
instead of exactly half.
\end{patternbox}

\subsection*{Understanding the problem carefully}

We again look for a split $(\textit{cut1}, \textit{cut2})$ with
$\textit{cut1} + \textit{cut2} = k$, $0 \le \textit{cut1} \le n_1$, and
$0 \le \textit{cut2} \le n_2$, such that the left side (the first
$\textit{cut1}$ elements of $\textit{nums1}$ and first $\textit{cut2}$
elements of $\textit{nums2}$) contains exactly the $k$ smallest elements
overall. The same four boundary values $l_1, r_1, l_2, r_2$ and the same
validity condition $l_1 \le r_2$ and $l_2 \le r_1$ from
Section~\ref{sec:bs-median-two-arrays} apply unchanged.

\begin{ideabox}[Candidate Space and Predicate]
The candidate space is $\textit{cut1} \in [\max(0, k-n_2), \min(k, n_1)]$
(the lower bound ensures $\textit{cut2} = k - \textit{cut1} \le n_2$; the
upper bound ensures $\textit{cut1} \le n_1$ and $\textit{cut2} \ge 0$).
The predicate is the same validity check as
Section~\ref{sec:bs-median-two-arrays}: $P(\textit{cut1}) = (l_1 \le r_2
\text{ and } l_2 \le r_1)$.
\end{ideabox}

\subsection*{Why only the target split size changes}

Nothing about the underlying correctness argument changes from
Section~\ref{sec:bs-median-two-arrays} -- we are still looking for the
unique partition where the left side's largest element does not exceed
the right side's smallest element, guaranteeing the left side is exactly
the $k$ smallest elements of the combined array. Once a valid partition is
found, the $k$-th smallest element is simply $\max(l_1, l_2)$ -- the
largest element among the $k$ elements placed on the left.

\subsection*{Algorithm}

\begin{enumerate}
  \item $\textit{lo} \gets \max(0, k-n_2)$,
        $\textit{hi} \gets \min(k, n_1)$.
  \item While $\textit{lo} \le \textit{hi}$:
  \begin{enumerate}
    \item $\textit{cut1} \gets \textit{lo}+(\textit{hi}-\textit{lo})/2$,
          $\textit{cut2} \gets k-\textit{cut1}$.
    \item Compute $l_1, r_1, l_2, r_2$ exactly as in
          Section~\ref{sec:bs-median-two-arrays}.
    \item If $l_1 \le r_2$ and $l_2 \le r_1$: return $\max(l_1,l_2)$.
    \item Else if $l_1 > r_2$: $\textit{hi} \gets \textit{cut1}-1$.
    \item Else: $\textit{lo} \gets \textit{cut1}+1$.
  \end{enumerate}
\end{enumerate}

\begin{algorithm}[H]
\caption{Kth Element of Two Sorted Arrays}
\begin{algorithmic}[1]
\Procedure{KthElement}{$\textit{nums1}, \textit{nums2}, k$}
  \State $\textit{lo} \gets \max(0, k-n_2)$, $\textit{hi} \gets \min(k, n_1)$
  \While{$\textit{lo} \le \textit{hi}$}
    \State $\textit{cut1} \gets \textit{lo}+(\textit{hi}-\textit{lo})/2$
    \State $\textit{cut2} \gets k - \textit{cut1}$
    \State $l_1 \gets (\textit{cut1}=0)\ ? -\infty : \textit{nums1}[\textit{cut1}-1]$
    \State $r_1 \gets (\textit{cut1}=n_1)\ ? +\infty : \textit{nums1}[\textit{cut1}]$
    \State $l_2 \gets (\textit{cut2}=0)\ ? -\infty : \textit{nums2}[\textit{cut2}-1]$
    \State $r_2 \gets (\textit{cut2}=n_2)\ ? +\infty : \textit{nums2}[\textit{cut2}]$
    \If{$l_1 \le r_2$ \textbf{and} $l_2 \le r_1$}
      \State \Return $\max(l_1, l_2)$
    \ElsIf{$l_1 > r_2$}
      \State $\textit{hi} \gets \textit{cut1} - 1$
    \Else
      \State $\textit{lo} \gets \textit{cut1} + 1$
    \EndIf
  \EndWhile
\EndProcedure
\end{algorithmic}
\end{algorithm}

\subsection*{Dry run}

\dryrun{Take $\textit{nums1} = [2,3,6,7,9]$, $\textit{nums2} = [1,4,8,10]$,
$k = 5$.}

$n_1=5$, $n_2=4$. $\textit{lo}=\max(0,5-4)=1$,
$\textit{hi}=\min(5,5)=5$.

\begin{itemize}
  \item $\textit{lo}=1,\textit{hi}=5$: $\textit{cut1}=3$,
        $\textit{cut2}=5-3=2$. $l_1=\textit{nums1}[2]=6$,
        $r_1=\textit{nums1}[3]=7$. $l_2=\textit{nums2}[1]=4$,
        $r_2=\textit{nums2}[2]=8$. Check: $l_1=6 \le r_2=8$? Yes.
        $l_2=4 \le r_1=7$? Yes. Valid partition.
\end{itemize}

The answer is $\max(l_1,l_2) = \max(6,4) = 6$. Checking by hand: the
merged array is $[1,2,3,4,6,7,8,9,10]$, whose $5$-th smallest element
(1-indexed) is indeed $6$.

\subsection*{C++ implementation}

\begin{lstlisting}
#include <bits/stdc++.h>
using namespace std;

int kthElement(vector<int>& nums1, vector<int>& nums2, int k) {
    int n1 = nums1.size(), n2 = nums2.size();

    int lo = max(0, k - n2);
    int hi = min(k, n1);

    while (lo <= hi) {
        int cut1 = lo + (hi - lo) / 2;
        int cut2 = k - cut1;

        int l1 = (cut1 == 0) ? INT_MIN : nums1[cut1 - 1];
        int r1 = (cut1 == n1) ? INT_MAX : nums1[cut1];
        int l2 = (cut2 == 0) ? INT_MIN : nums2[cut2 - 1];
        int r2 = (cut2 == n2) ? INT_MAX : nums2[cut2];

        if (l1 <= r2 && l2 <= r1) {
            return max(l1, l2);
        } else if (l1 > r2) {
            hi = cut1 - 1;
        } else {
            lo = cut1 + 1;
        }
    }
    return -1; // unreachable for valid input with 1 <= k <= n1+n2
}

int main() {
    vector<int> nums1 = {2, 3, 6, 7, 9};
    vector<int> nums2 = {1, 4, 8, 10};
    cout << "5th smallest element: "
         << kthElement(nums1, nums2, 5) << endl;
    return 0;
}
\end{lstlisting}

\begin{complexitybox}
\textbf{Time Complexity.} The binary search range has width at most
$\min(k, n_1) - \max(0, k-n_2) \le \min(n_1, n_2)$, giving
\[
  O(\log(\min(n_1, n_2))).
\]
\textbf{Space Complexity.} $O(1)$ auxiliary space.
\end{complexitybox}

\begin{takeawaybox}
Median-of-two-sorted-arrays and $k$-th-element-of-two-sorted-arrays are the
same partition-search technique with only the target split size changed
from ``half'' to ``$k$.'' Whenever you solve one binary-search-on-a-
partition problem, immediately check whether nearby sheet problems are
just different target split sizes of the same idea.
\end{takeawaybox}

\section{Row with Maximum Number of 1s}
\label{sec:bs-row-max-ones}

\noindent We now begin the third and final group of problems in this
chapter: \textbf{binary search on a 2D array}. The recurring theme is that
each \emph{row} (or column) of the matrix is independently sorted, which
lets us reuse the 1D binary search templates from
Sections~\ref{sec:bs-find-x}--\ref{sec:bs-peak-1d} \emph{per row or
column}, or combine row-sortedness and column-sortedness to search the
whole 2D index space in a single sweep.

\problemstatement{We are given an $n \times m$ binary matrix in which
\textbf{every row is sorted} (all $0$s appear before all $1$s within that
row). We must find the index of the row containing the \textbf{maximum
number of $1$s}.}

\begin{patternbox}
\emph{``Each row is sorted''} (even though the matrix as a whole is not
globally ordered in any simple way) is the keyword that tells us we can
binary search \emph{within each row independently}, rather than needing
any 2D-aware search technique. Whenever a 2D-array problem's sortedness
guarantee is row-by-row (or column-by-column) rather than global,
default first to ``run a 1D binary search on each row/column,'' since it
is almost always the simplest correct approach, even if not always the
asymptotically fastest one.
\end{patternbox}

\subsection*{Understanding the problem carefully}

Because each row is sorted with $0$s before $1$s, the number of $1$s in a
row equals $m - \textit{lowerBound}(\text{row}, 1)$, where
$\textit{lowerBound}(\text{row}, 1)$ is the first index in that row where
the value is $\ge 1$ (Section~\ref{sec:bs-lower-bound}). Computing this for
every row and tracking the maximum gives the answer.

\begin{ideabox}[Reuse, Don't Re-derive]
Apply Algorithm~\ref{sec:bs-lower-bound}'s \textsc{LowerBound} procedure
to each row independently (searching for the first $1$), convert each
result to a count of $1$s, and track the row index achieving the maximum
count.
\end{ideabox}

\subsection*{A note on the alternative $O(n+m)$ staircase approach}

A different, even faster technique exists for this specific problem:
starting from the top-right corner of the matrix, move left whenever the
current cell is $1$ (recording this row as the best-so-far and shrinking
the column boundary), and move down whenever the current cell is $0$. This
``staircase search'' exploits row-and-column-sortedness jointly and runs
in $O(n+m)$ total, faster than the $O(n \log m)$ of the per-row binary
search. We mention it here because it foreshadows the 2D search technique
used in Section~\ref{sec:bs-search-2d-matrix-2}, but the per-row binary
search remains the more directly transferable technique for this
chapter's purposes, since it requires no new derivation beyond what
Section~\ref{sec:bs-lower-bound} already established.

\subsection*{Algorithm (Per-Row Binary Search)}

\begin{enumerate}
  \item Initialize $\textit{bestRow} \gets -1$, $\textit{bestCount}
        \gets 0$.
  \item For each row $i$ from $0$ to $n-1$:
  \begin{enumerate}
    \item $\textit{idx} \gets \textit{lowerBound}(\text{row } i, 1)$.
    \item $\textit{count} \gets m - \textit{idx}$.
    \item If $\textit{count} > \textit{bestCount}$: update
          $\textit{bestCount} \gets \textit{count}$,
          $\textit{bestRow} \gets i$.
  \end{enumerate}
  \item Return \textit{bestRow}.
\end{enumerate}

\begin{algorithm}[H]
\caption{Row with Maximum Number of 1s}
\begin{algorithmic}[1]
\Procedure{RowWithMaxOnes}{\textit{matrix}}
  \State $\textit{bestRow} \gets -1$, $\textit{bestCount} \gets 0$
  \For{$i \gets 0$ to $n-1$}
    \State $\textit{idx} \gets \Call{LowerBound}{\textit{matrix}[i], 1}$
    \State $\textit{count} \gets m - \textit{idx}$
    \If{$\textit{count} > \textit{bestCount}$}
      \State $\textit{bestCount} \gets \textit{count}$, $\textit{bestRow} \gets i$
    \EndIf
  \EndFor
  \State \Return \textit{bestRow}
\EndProcedure
\end{algorithmic}
\end{algorithm}

\subsection*{Dry run}

\dryrun{Take the matrix
$\begin{bmatrix} 0&0&0 \\ 0&1&1 \\ 1&1&1 \\ 0&0&0 \end{bmatrix}$.}

\begin{itemize}
  \item Row $0$: $[0,0,0]$. Lower bound of $1$ is index $3$ (no $1$ found).
        Count $= 3-3=0$.
  \item Row $1$: $[0,1,1]$. Lower bound of $1$ is index $1$. Count
        $=3-1=2$. Best so far: row $1$, count $2$.
  \item Row $2$: $[1,1,1]$. Lower bound of $1$ is index $0$. Count
        $=3-0=3$. Best so far: row $2$, count $3$.
  \item Row $3$: $[0,0,0]$. Count $=0$.
\end{itemize}

The answer is row $2$, with $3$ ones.

\subsection*{C++ implementation}

\begin{lstlisting}
#include <bits/stdc++.h>
using namespace std;

int lowerBoundRow(vector<int>& row, int X) {
    int lo = 0, hi = row.size() - 1, ans = row.size();
    while (lo <= hi) {
        int mid = lo + (hi - lo) / 2;
        if (row[mid] >= X) { ans = mid; hi = mid - 1; }
        else lo = mid + 1;
    }
    return ans;
}

int rowWithMaxOnes(vector<vector<int>>& matrix) {
    int n = matrix.size(), m = matrix[0].size();
    int bestRow = -1, bestCount = 0;

    for (int i = 0; i < n; i++) {
        int idx = lowerBoundRow(matrix[i], 1);
        int count = m - idx;
        if (count > bestCount) {
            bestCount = count;
            bestRow = i;
        }
    }
    return bestRow;
}

int main() {
    vector<vector<int>> matrix = {
        {0,0,0},
        {0,1,1},
        {1,1,1},
        {0,0,0}
    };
    cout << "Row with max ones: " << rowWithMaxOnes(matrix) << endl;
    return 0;
}
\end{lstlisting}

\begin{complexitybox}
\textbf{Time Complexity.} For each of the $n$ rows, a binary search costs
$O(\log m)$, giving an overall time complexity of
\[
  O(n \log m).
\]
\textbf{Space Complexity.} $O(1)$ auxiliary space.
\end{complexitybox}

\begin{takeawaybox}
The simplest way to handle a 2D problem where sortedness is guaranteed
only \emph{per row} (or per column) is to apply a 1D binary search
template independently to each row, accepting an extra factor of $n$ (or
$m$) in the time complexity in exchange for needing no new derivation.
Faster joint row-and-column techniques exist for some such problems (the
staircase search above, and Section~\ref{sec:bs-search-2d-matrix-2}), but
they require the matrix to be sorted in \emph{both} directions
simultaneously.
\end{takeawaybox}

\section{Search in a 2D Matrix}
\label{sec:bs-search-2d-matrix}

\problemstatement{We are given an $n \times m$ matrix with the following
guarantees: every row is sorted in increasing order, and the first element
of every row (except the first) is strictly greater than the last element
of the previous row. Given a target $X$, determine whether $X$ exists in
the matrix, in $O(\log(nm))$ time.}

\begin{patternbox}
The extra guarantee here -- \emph{each row's first element exceeds the
previous row's last element} -- is what elevates this beyond ``each row is
sorted independently'' (Section~\ref{sec:bs-row-max-ones}): it means that
if you read the matrix row by row, left to right, the values you encounter
are \textbf{globally non-decreasing}, exactly as if the matrix were a
single flattened, sorted 1D array. Whenever a 2D matrix problem gives you
this row-chaining guarantee, immediately think ``flatten conceptually and
run ordinary binary search,'' rather than reaching for a more complex 2D
technique.
\end{patternbox}

\subsection*{Understanding the problem carefully}

Because the matrix reads as a single sorted sequence when traversed
row-major (row $0$ left to right, then row $1$ left to right, and so on),
we can treat it as a virtual 1D array of length $nm$, where virtual index
$\textit{idx}$ corresponds to matrix cell
$(\textit{idx} / m,\ \textit{idx} \bmod m)$ (integer division for the row,
remainder for the column). No actual flattening or copying is needed --
we simply convert each binary search midpoint index into a
$(\textit{row}, \textit{col})$ pair on the fly.

\begin{ideabox}[Candidate Space and Predicate]
The candidate space is virtual indices $0, \dots, nm-1$. The predicate is
exactly the same as ordinary binary search
(Section~\ref{sec:bs-find-x}): compare the value at the virtual index to
$X$ and discard the appropriate half.
\end{ideabox}

\subsection*{Why the row-chaining guarantee makes flattening valid}

Ordinary binary search needs the entire candidate sequence (values, in
virtual-index order) to be sorted. The row-chaining guarantee is precisely
what makes this true here: within a row, sortedness gives increasing
values left to right; across the row boundary, the guarantee that the
next row's first element exceeds the previous row's last element ensures
no drop in value occurs when virtual index crosses from the end of one row
to the start of the next. So the entire virtual sequence is non-decreasing,
and ordinary binary search's discard argument
(Section~\ref{sec:bs-find-x}) applies without any modification beyond the
index-mapping arithmetic.

\subsection*{Algorithm}

\begin{enumerate}
  \item Initialize $\textit{lo} \gets 0$, $\textit{hi} \gets nm-1$.
  \item While $\textit{lo} \le \textit{hi}$:
  \begin{enumerate}
    \item $\textit{mid} \gets \textit{lo}+(\textit{hi}-\textit{lo})/2$.
    \item $\textit{row} \gets \textit{mid}/m$, $\textit{col} \gets
          \textit{mid} \bmod m$.
    \item If $\textit{matrix}[\textit{row}][\textit{col}] = X$: return
          \texttt{true}.
    \item If $\textit{matrix}[\textit{row}][\textit{col}] < X$:
          $\textit{lo} \gets \textit{mid}+1$.
    \item Else: $\textit{hi} \gets \textit{mid}-1$.
  \end{enumerate}
  \item If the loop ends without returning, return \texttt{false}.
\end{enumerate}

\begin{algorithm}[H]
\caption{Search in a 2D Matrix}
\begin{algorithmic}[1]
\Procedure{SearchMatrix}{\textit{matrix}, $X$}
  \State $\textit{lo} \gets 0$, $\textit{hi} \gets nm-1$
  \While{$\textit{lo} \le \textit{hi}$}
    \State $\textit{mid} \gets \textit{lo}+(\textit{hi}-\textit{lo})/2$
    \State $\textit{row} \gets \textit{mid} / m$, $\textit{col} \gets \textit{mid} \bmod m$
    \If{$\textit{matrix}[\textit{row}][\textit{col}] = X$} \Return \textbf{true} \EndIf
    \If{$\textit{matrix}[\textit{row}][\textit{col}] < X$}
      \State $\textit{lo} \gets \textit{mid} + 1$
    \Else
      \State $\textit{hi} \gets \textit{mid} - 1$
    \EndIf
  \EndWhile
  \State \Return \textbf{false}
\EndProcedure
\end{algorithmic}
\end{algorithm}

\subsection*{Dry run}

\dryrun{Take $\textit{matrix} = \begin{bmatrix} 1&3&5&7 \\ 10&11&16&20 \\
23&30&34&60 \end{bmatrix}$, $X = 3$. Here $n=3$, $m=4$, $nm=12$.}

\begin{itemize}
  \item $\textit{lo}=0,\textit{hi}=11$: $\textit{mid}=5$,
        $\textit{row}=1,\textit{col}=1$,
        $\textit{matrix}[1][1]=11 > 3$. $\textit{hi}=4$.
  \item $\textit{lo}=0,\textit{hi}=4$: $\textit{mid}=2$,
        $\textit{row}=0,\textit{col}=2$,
        $\textit{matrix}[0][2]=5 > 3$. $\textit{hi}=1$.
  \item $\textit{lo}=0,\textit{hi}=1$: $\textit{mid}=0$,
        $\textit{row}=0,\textit{col}=0$,
        $\textit{matrix}[0][0]=1 < 3$. $\textit{lo}=1$.
  \item $\textit{lo}=1,\textit{hi}=1$: $\textit{mid}=1$,
        $\textit{row}=0,\textit{col}=1$,
        $\textit{matrix}[0][1]=3 = 3$. Return \texttt{true}.
\end{itemize}

\subsection*{C++ implementation}

\begin{lstlisting}
#include <bits/stdc++.h>
using namespace std;

bool searchMatrix(vector<vector<int>>& matrix, int X) {
    int n = matrix.size(), m = matrix[0].size();
    int lo = 0, hi = n * m - 1;

    while (lo <= hi) {
        int mid = lo + (hi - lo) / 2;
        int row = mid / m, col = mid % m;

        if (matrix[row][col] == X) return true;
        else if (matrix[row][col] < X) lo = mid + 1;
        else hi = mid - 1;
    }
    return false;
}

int main() {
    vector<vector<int>> matrix = {
        {1, 3, 5, 7},
        {10, 11, 16, 20},
        {23, 30, 34, 60}
    };
    cout << (searchMatrix(matrix, 3) ? "true" : "false") << endl;
    return 0;
}
\end{lstlisting}

\begin{complexitybox}
\textbf{Time Complexity.} The virtual candidate range of size $nm$ halves
each iteration, giving
\[
  O(\log(nm)).
\]
\textbf{Space Complexity.} $O(1)$ auxiliary space.
\end{complexitybox}

\begin{takeawaybox}
A 2D structure with a strong enough chaining guarantee between rows can be
searched with completely unmodified 1D binary search, as long as you
translate a single virtual index into a $(\textit{row}, \textit{col})$
pair via division and modulo. Always check first whether a 2D problem's
sortedness guarantee is this strong before reaching for a more specialized
2D technique like the one needed in
Section~\ref{sec:bs-search-2d-matrix-2}.
\end{takeawaybox}

\section{Search in a 2D Matrix II}
\label{sec:bs-search-2d-matrix-2}

\problemstatement{We are given an $n \times m$ matrix in which every row
is sorted in increasing order \textbf{left to right}, and every column is
sorted in increasing order \textbf{top to bottom} -- but, unlike
Section~\ref{sec:bs-search-2d-matrix}, there is \textbf{no} guarantee
chaining one row to the next (the last element of a row may be larger
than the first element of the next row). Given a target $X$, determine
whether it exists in the matrix.}

\begin{patternbox}
The missing row-chaining guarantee is the keyword to notice by its
\emph{absence}: it means the matrix can \textbf{not} be treated as a
flattened sorted 1D array (Section~\ref{sec:bs-search-2d-matrix}'s
virtual-index trick would give wrong answers here, since the virtual
sequence is not globally sorted). Instead, the two \emph{independent}
sortedness guarantees -- rows increasing rightward, columns increasing
downward -- call for a different technique: starting from a corner where
one direction increases and the other decreases the candidate's relation
to $X$, and walking a single staircase path.
\end{patternbox}

\subsection*{Understanding the problem carefully}

Start at the \textbf{top-right} corner, $(\textit{row}, \textit{col}) =
(0, m-1)$. From this corner: moving \textbf{left} (decreasing
\textit{col}) only ever decreases the value (row sortedness), and moving
\textbf{down} (increasing \textit{row}) only ever increases the value
(column sortedness). This gives us a clean three-way comparison at every
step:

\begin{itemize}
  \item If $\textit{matrix}[\textit{row}][\textit{col}] = X$: found it.
  \item If $\textit{matrix}[\textit{row}][\textit{col}] > X$: every entry
        further right in this row is even larger (row sortedness), so $X$
        cannot be in the rest of this row; move left:
        $\textit{col} \gets \textit{col}-1$.
  \item If $\textit{matrix}[\textit{row}][\textit{col}] < X$: every entry
        further up in this column is even smaller (column sortedness), so
        $X$ cannot be above in this column; move down:
        $\textit{row} \gets \textit{row}+1$.
\end{itemize}

\begin{ideabox}[Why the Top-Right Corner Is the Only Safe Starting Point]
The top-left corner is the smallest element in the entire matrix (both
row and column sortedness make every other element $\ge$ it), so both
``move right'' and ``move down'' would increase the value -- a comparison
there gives no way to discard a direction safely. The bottom-right corner
is symmetric (the largest element, both directions only decrease).
\emph{Only} the top-right (or, symmetrically, bottom-left) corner has the
property that one direction strictly increases the value while the other
strictly decreases it, which is exactly what is needed to safely discard
an entire row or column on each comparison.
\end{ideabox}

\subsection*{Why this staircase walk is correct and terminates in $O(n+m)$}

At each step, the walk either eliminates an entire remaining row's worth
of values (when it moves left, the eliminated value and everything to its
right in that row are confirmed too large) or an entire remaining column's
worth (when it moves down, the eliminated value and everything above it in
that column are confirmed too small) -- but more precisely, each step
moves the pointer one position left or one position down, and since
\textit{col} can decrease at most $m$ times and \textit{row} can increase
at most $n$ times before going out of bounds, the walk terminates after at
most $n+m$ steps, either finding $X$ or running off the matrix entirely
(confirming $X$ is absent).

\subsection*{Algorithm}

\begin{enumerate}
  \item Initialize $\textit{row} \gets 0$, $\textit{col} \gets m-1$.
  \item While $\textit{row} < n$ and $\textit{col} \ge 0$:
  \begin{enumerate}
    \item If $\textit{matrix}[\textit{row}][\textit{col}] = X$: return
          \texttt{true}.
    \item If $\textit{matrix}[\textit{row}][\textit{col}] > X$:
          $\textit{col} \gets \textit{col}-1$.
    \item Else: $\textit{row} \gets \textit{row}+1$.
  \end{enumerate}
  \item Return \texttt{false}.
\end{enumerate}

\begin{algorithm}[H]
\caption{Search in a 2D Matrix II (Staircase Search)}
\begin{algorithmic}[1]
\Procedure{SearchMatrixII}{\textit{matrix}, $X$}
  \State $\textit{row} \gets 0$, $\textit{col} \gets m-1$
  \While{$\textit{row} < n$ \textbf{and} $\textit{col} \ge 0$}
    \If{$\textit{matrix}[\textit{row}][\textit{col}] = X$} \Return \textbf{true} \EndIf
    \If{$\textit{matrix}[\textit{row}][\textit{col}] > X$}
      \State $\textit{col} \gets \textit{col} - 1$
    \Else
      \State $\textit{row} \gets \textit{row} + 1$
    \EndIf
  \EndWhile
  \State \Return \textbf{false}
\EndProcedure
\end{algorithmic}
\end{algorithm}

\subsection*{Dry run}

\dryrun{Take $\textit{matrix} = \begin{bmatrix} 1&4&7&11 \\ 2&5&8&12 \\
3&6&9&16 \\ 10&13&14&17 \end{bmatrix}$, $X = 5$. ($n=m=4$.)}

\begin{itemize}
  \item $\textit{row}=0,\textit{col}=3$: $\textit{matrix}[0][3]=11 > 5$.
        $\textit{col}=2$.
  \item $\textit{row}=0,\textit{col}=2$: $\textit{matrix}[0][2]=7 > 5$.
        $\textit{col}=1$.
  \item $\textit{row}=0,\textit{col}=1$: $\textit{matrix}[0][1]=4 < 5$.
        $\textit{row}=1$.
  \item $\textit{row}=1,\textit{col}=1$: $\textit{matrix}[1][1]=5 = 5$.
        Return \texttt{true}.
\end{itemize}

\subsection*{C++ implementation}

\begin{lstlisting}
#include <bits/stdc++.h>
using namespace std;

bool searchMatrixII(vector<vector<int>>& matrix, int X) {
    int n = matrix.size(), m = matrix[0].size();
    int row = 0, col = m - 1;

    while (row < n && col >= 0) {
        if (matrix[row][col] == X) return true;
        else if (matrix[row][col] > X) col--;
        else row++;
    }
    return false;
}

int main() {
    vector<vector<int>> matrix = {
        {1, 4, 7, 11},
        {2, 5, 8, 12},
        {3, 6, 9, 16},
        {10, 13, 14, 17}
    };
    cout << (searchMatrixII(matrix, 5) ? "true" : "false") << endl;
    return 0;
}
\end{lstlisting}

\begin{complexitybox}
\textbf{Time Complexity.} Each step decreases \textit{col} or increases
\textit{row}, and there are at most $m$ decreases and $n$ increases before
the pointer exits the matrix, giving
\[
  O(n+m).
\]
Note this is not $O(\log(nm))$ -- the lack of row-chaining genuinely costs
us the logarithmic complexity available in
Section~\ref{sec:bs-search-2d-matrix}, even though this is still a vast
improvement over the $O(nm)$ brute-force scan.
\textbf{Space Complexity.} $O(1)$ auxiliary space.
\end{complexitybox}

\begin{takeawaybox}
When a 2D matrix is sorted along both rows and columns independently but
not chained into one global order, look for a corner where the two
sortedness directions act as \emph{opposite} comparisons (one direction
increasing, the other decreasing) relative to the target -- that corner is
always the unique safe starting point for a linear ``staircase'' walk,
even though it does not give the $O(\log(nm))$ that full chaining would
allow.
\end{takeawaybox}

\section{Find a Peak Element in a 2D Matrix}
\label{sec:bs-peak-2d}

\problemstatement{We are given an $n \times m$ matrix of distinct
integers, where no two horizontally or vertically adjacent cells are
equal. A cell is a \textbf{peak} if it is strictly greater than all of its
(up to four) horizontal and vertical neighbors; cells outside the matrix
are treated as $-\infty$. Find the position of any one peak.}

\begin{patternbox}
This is the 2D generalization of Section~\ref{sec:bs-peak-1d}, and the
keyword pattern is identical: no sortedness, just a local-comparison rule
plus an implicit $-\infty$ border. The new idea needed is how to binary
search over \emph{one} dimension (say, columns) while using an $O(n)$
linear scan over the \emph{other} dimension (rows) as part of each
predicate evaluation -- a ``binary search over one axis, brute-force scan
over the other'' technique that is common whenever a 2D problem reduces
to a 1D binary search plus a per-step linear helper.
\end{patternbox}

\subsection*{Understanding the problem carefully}

Fix a candidate middle column $\textit{midCol}$. Scan down that column to
find the row $\textit{maxRow}$ containing the largest value in that
column -- call it $v = \textit{matrix}[\textit{maxRow}][\textit{midCol}]$.
Compare $v$ to its left neighbor $\textit{matrix}[\textit{maxRow}]
[\textit{midCol}-1]$ and right neighbor
$\textit{matrix}[\textit{maxRow}][\textit{midCol}+1]$ (treating
out-of-range neighbors as $-\infty$).

\begin{itemize}
  \item If $v$ is greater than both neighbors: $v$ is greater than
        everything above and below it in its own column (by choice of
        $\textit{maxRow}$ as the column's maximum) \emph{and} greater than
        its left and right neighbors -- so $v$ is a peak. Done.
  \item If the left neighbor is greater than $v$: since that left
        neighbor is, in particular, greater than the column maximum at
        $\textit{midCol}$, the column $\textit{midCol}-1$ (or further
        left) must contain an even larger local structure -- a peak is
        guaranteed to exist somewhere in the columns
        $[\textit{lo}, \textit{midCol}-1]$, by the same $-\infty$-border
        existence argument as Section~\ref{sec:bs-peak-1d}, applied along
        the column axis.
  \item Symmetrically, if the right neighbor is greater, a peak is
        guaranteed in columns $[\textit{midCol}+1, \textit{hi}]$.
\end{itemize}

\begin{ideabox}[Candidate Space and Predicate]
The candidate space is columns $0, \dots, m-1$. At each candidate column,
an $O(n)$ scan finds that column's maximum value's row, and a constant-time
comparison to its horizontal neighbors decides whether to stop (peak
found) or which half of the remaining columns to search next -- the same
``climbing toward a guaranteed peak'' argument as
Section~\ref{sec:bs-peak-1d}, just with an entire column's maximum
standing in for a single array element.
\end{ideabox}

\subsection*{Why restricting attention to each column's maximum is valid}

The reason we only need to examine the \emph{maximum} of each candidate
column, rather than every cell in it, is that if a peak exists in this
column, it is impossible for it to be at a position other than the
column's maximum: a peak must be greater than its upper and lower
neighbors within the same column, which already forces it to be at least
as large as both of those, and chaining this reasoning shows the only cell
in a column that could possibly satisfy ``greater than both vertical
neighbors'' even before checking horizontal neighbors is necessarily a
local maximum along that column. Examining the column's global maximum is
therefore sufficient to either confirm a peak or correctly decide which
direction to continue the binary search.

\subsection*{Algorithm}

\begin{enumerate}
  \item Initialize $\textit{lo} \gets 0$, $\textit{hi} \gets m-1$.
  \item While $\textit{lo} \le \textit{hi}$:
  \begin{enumerate}
    \item $\textit{midCol} \gets \textit{lo}+(\textit{hi}-\textit{lo})/2$.
    \item Scan column $\textit{midCol}$ to find $\textit{maxRow}$, the row
          of its maximum value $v$.
    \item Compare $v$ to its left and right neighbors.
    \item If $v$ exceeds both: return $(\textit{maxRow},
          \textit{midCol})$.
    \item If the left neighbor exceeds $v$: $\textit{hi} \gets
          \textit{midCol}-1$.
    \item Else (right neighbor exceeds $v$): $\textit{lo} \gets
          \textit{midCol}+1$.
  \end{enumerate}
\end{enumerate}

\begin{algorithm}[H]
\caption{Find Peak Element in a 2D Matrix}
\begin{algorithmic}[1]
\Procedure{FindPeak2D}{\textit{matrix}}
  \State $\textit{lo} \gets 0$, $\textit{hi} \gets m-1$
  \While{$\textit{lo} \le \textit{hi}$}
    \State $\textit{midCol} \gets \textit{lo}+(\textit{hi}-\textit{lo})/2$
    \State $\textit{maxRow} \gets \arg\max_i \textit{matrix}[i][\textit{midCol}]$
    \State $v \gets \textit{matrix}[\textit{maxRow}][\textit{midCol}]$
    \State $\textit{left} \gets (\textit{midCol}=0)\ ? -\infty : \textit{matrix}[\textit{maxRow}][\textit{midCol}-1]$
    \State $\textit{right} \gets (\textit{midCol}=m-1)\ ? -\infty : \textit{matrix}[\textit{maxRow}][\textit{midCol}+1]$
    \If{$v > \textit{left}$ \textbf{and} $v > \textit{right}$}
      \State \Return $(\textit{maxRow}, \textit{midCol})$
    \ElsIf{$\textit{left} > v$}
      \State $\textit{hi} \gets \textit{midCol} - 1$
    \Else
      \State $\textit{lo} \gets \textit{midCol} + 1$
    \EndIf
  \EndWhile
\EndProcedure
\end{algorithmic}
\end{algorithm}

\subsection*{Dry run}

\dryrun{Take $\textit{matrix} = \begin{bmatrix} 10&8&10&10 \\ 14&13&12&11
\\ 15&9&11&21 \\ 16&17&19&20 \end{bmatrix}$ ($n=m=4$).}

\begin{itemize}
  \item $\textit{lo}=0,\textit{hi}=3$: $\textit{midCol}=1$. Column $1$:
        $[8,13,9,17]$; max is $17$ at row $3$. Left neighbor
        $\textit{matrix}[3][0]=16$, right neighbor
        $\textit{matrix}[3][2]=19$. Is $17 > 16$ and $17 > 19$? No
        ($17 < 19$). Right neighbor is greater: $\textit{lo}=2$.
  \item $\textit{lo}=2,\textit{hi}=3$: $\textit{midCol}=2$. Column $2$:
        $[10,12,11,19]$; max is $19$ at row $3$. Left neighbor
        $\textit{matrix}[3][1]=17$, right neighbor
        $\textit{matrix}[3][3]=20$. Is $19>17$ and $19>20$? No
        ($19<20$). Right neighbor greater: $\textit{lo}=3$.
  \item $\textit{lo}=3,\textit{hi}=3$: $\textit{midCol}=3$. Column $3$:
        $[10,11,21,20]$; max is $21$ at row $2$. Left neighbor
        $\textit{matrix}[2][2]=11$. No right neighbor (out of range,
        $-\infty$). Is $21>11$ and $21>-\infty$? Yes. Peak found at
        $(2,3)$.
\end{itemize}

\subsection*{C++ implementation}

\begin{lstlisting}
#include <bits/stdc++.h>
using namespace std;

pair<int,int> findPeak2D(vector<vector<int>>& matrix) {
    int n = matrix.size(), m = matrix[0].size();
    int lo = 0, hi = m - 1;

    while (lo <= hi) {
        int midCol = lo + (hi - lo) / 2;
        int maxRow = 0;
        for (int i = 1; i < n; i++) {
            if (matrix[i][midCol] > matrix[maxRow][midCol]) maxRow = i;
        }

        int v = matrix[maxRow][midCol];
        int left  = (midCol == 0)     ? INT_MIN : matrix[maxRow][midCol - 1];
        int right = (midCol == m - 1) ? INT_MIN : matrix[maxRow][midCol + 1];

        if (v > left && v > right) {
            return {maxRow, midCol};
        } else if (left > v) {
            hi = midCol - 1;
        } else {
            lo = midCol + 1;
        }
    }
    return {-1, -1}; // unreachable for a valid matrix
}

int main() {
    vector<vector<int>> matrix = {
        {10, 8, 10, 10},
        {14, 13, 12, 11},
        {15, 9, 11, 21},
        {16, 17, 19, 20}
    };
    auto [r, c] = findPeak2D(matrix);
    cout << "Peak at (" << r << ", " << c << ")" << endl;
    return 0;
}
\end{lstlisting}

\begin{complexitybox}
\textbf{Time Complexity.} The column range halves each iteration
($O(\log m)$ iterations), and each iteration scans a full column in
$O(n)$, giving an overall time complexity of
\[
  O(n \log m).
\]
\textbf{Space Complexity.} $O(1)$ auxiliary space.
\end{complexitybox}

\begin{takeawaybox}
When a 2D problem's local-comparison structure generalizes a 1D pattern
you already know (here, peak-finding from
Section~\ref{sec:bs-peak-1d}), look for a way to binary search one axis
while reducing each candidate to a single representative value via a
linear scan along the other axis -- the representative value's own local
comparisons then drive the binary search exactly as in the 1D version.
\end{takeawaybox}

\section{Median of a Row-wise Sorted Matrix}
\label{sec:bs-matrix-median}

\problemstatement{We are given an $n \times m$ matrix in which every row
is sorted in increasing order (with no chaining guarantee between rows),
and $nm$ is odd. Find the median of all $nm$ elements.}

\begin{patternbox}
This closing problem of the chapter combines two ideas from earlier
sections: \emph{binary search on the answer}
(Sections~\ref{sec:bs-sqrt}--\ref{sec:bs-kth-element}), where the
candidate space is a range of \emph{values} rather than indices, and
\emph{per-row binary search} (Section~\ref{sec:bs-row-max-ones}), where
row-sortedness is exploited independently per row inside a helper
function. The keyword combination \emph{``median''} plus \emph{``each row
sorted''} plus a request for better than $O(nm \log(nm))$ (which a full
sort-and-pick would cost) should trigger exactly this combination.
\end{patternbox}

\subsection*{Understanding the problem carefully}

The median of $nm$ elements (with $nm$ odd) is the value $v$ such that
exactly $(nm-1)/2$ elements are strictly smaller than $v$, and $v$ itself
is present in the matrix. Equivalently, $v$ is the smallest value for
which $\textit{countLessOrEqual}(v) \ge \lceil nm/2 \rceil$, where
$\textit{countLessOrEqual}(v)$ counts how many matrix elements are $\le
v$. We do not need to merge or sort the matrix to compute this count: for
each row, since the row is independently sorted, the count of elements
$\le v$ in that row is exactly $\textit{upperBound}(\text{row}, v)$
(Section~\ref{sec:bs-upper-bound}, the first index strictly greater than
$v$, which equals the count of elements $\le v$). Summing this across all
$n$ rows gives $\textit{countLessOrEqual}(v)$ for the whole matrix in
$O(n \log m)$.

\begin{ideabox}[Candidate Space and Predicate]
The candidate space is integer values in
$[\min(\textit{matrix}), \max(\textit{matrix})]$ (the minimum and maximum
single elements across the whole matrix, found by checking the first and
last elements of each row, since rows are sorted). The predicate is
$P(v) = (\textit{countLessOrEqual}(v) \ge \lceil nm/2 \rceil)$, and we want
the smallest $v$ where $P$ holds.
\end{ideabox}

\subsection*{Why this nested binary search correctly finds the median}

As $v$ increases, $\textit{countLessOrEqual}(v)$ can only increase or stay
the same (a larger threshold never excludes an element that a smaller
threshold included), so $P$ is monotonic, justifying the outer binary
search over $v$. The smallest $v$ satisfying $P$ is, by definition, the
smallest value such that at least half of the matrix's elements do not
exceed it -- which is exactly the definition of the median for an odd
total count. Crucially, $v$ itself does not need to be searched for
separately within the matrix: because $\textit{countLessOrEqual}$ jumps by
at least $1$ at $v$ if and only if $v$ is actually present in the matrix
(if $v$ were absent, the count at $v$ and at $v-1$ would be identical, so
the crossover could not happen exactly at an absent value), the binary
search automatically lands on a value that is genuinely in the matrix.

\subsection*{Algorithm}

\begin{enumerate}
  \item $\textit{lo} \gets \min$ over all rows' first elements;
        $\textit{hi} \gets \max$ over all rows' last elements.
  \item While $\textit{lo} \le \textit{hi}$:
  \begin{enumerate}
    \item $\textit{mid} \gets \textit{lo}+(\textit{hi}-\textit{lo})/2$.
    \item For each row, add $\textit{upperBound}(\text{row},
          \textit{mid})$ to get $\textit{count}$.
    \item If $\textit{count} < \lceil nm/2 \rceil$: $\textit{lo} \gets
          \textit{mid}+1$.
    \item Else: $\textit{hi} \gets \textit{mid}-1$.
  \end{enumerate}
  \item Return \textit{lo}.
\end{enumerate}

\begin{algorithm}[H]
\caption{Median of a Row-wise Sorted Matrix}
\begin{algorithmic}[1]
\Procedure{MatrixMedian}{\textit{matrix}}
  \State $\textit{lo} \gets \min_i \textit{matrix}[i][0]$, $\textit{hi} \gets \max_i \textit{matrix}[i][m-1]$
  \State $\textit{half} \gets \lceil nm/2 \rceil$
  \While{$\textit{lo} \le \textit{hi}$}
    \State $\textit{mid} \gets \textit{lo}+(\textit{hi}-\textit{lo})/2$
    \State $\textit{count} \gets \sum_i \Call{UpperBound}{\textit{matrix}[i], \textit{mid}}$
    \If{$\textit{count} < \textit{half}$}
      \State $\textit{lo} \gets \textit{mid} + 1$
    \Else
      \State $\textit{hi} \gets \textit{mid} - 1$
    \EndIf
  \EndWhile
  \State \Return \textit{lo}
\EndProcedure
\end{algorithmic}
\end{algorithm}

\subsection*{Dry run}

\dryrun{Take $\textit{matrix} = \begin{bmatrix} 1&3&5 \\ 2&6&9 \\
3&6&9 \end{bmatrix}$ ($n=m=3$, $nm=9$).}

$\textit{half} = \lceil 9/2 \rceil = 5$. $\textit{lo}=\min(1,2,3)=1$,
$\textit{hi}=\max(5,9,9)=9$.

\begin{itemize}
  \item $\textit{lo}=1,\textit{hi}=9$: $\textit{mid}=5$. Counts $\le 5$:
        row $1$: $[1,3,5]$, upper bound of $5$ is index $3$ (all $3$
        qualify). Row $2$: $[2,6,9]$, upper bound of $5$ is index $1$
        (just the $2$). Row $3$: $[3,6,9]$, upper bound of $5$ is index
        $1$ (just the $3$). Total $=3+1+1=5 \ge 5$. $\textit{hi}=4$.
  \item $\textit{lo}=1,\textit{hi}=4$: $\textit{mid}=2$. Counts $\le 2$:
        row $1$: upper bound of $2$ is index $1$ (just the $1$). Row $2$:
        upper bound of $2$ is index $1$ (just the $2$). Row $3$: upper
        bound of $2$ is index $0$ (none). Total $=1+1+0=2 < 5$.
        $\textit{lo}=3$.
  \item $\textit{lo}=3,\textit{hi}=4$: $\textit{mid}=3$. Counts $\le 3$:
        row $1$: index $2$ (the $1,3$). Row $2$: index $1$ (just the
        $2$). Row $3$: index $1$ (just the $3$). Total $=2+1+1=4 < 5$.
        $\textit{lo}=4$.
  \item $\textit{lo}=4,\textit{hi}=4$: $\textit{mid}=4$. Counts $\le 4$:
        same as $\le 3$ for these rows (no element equals $4$), so total
        $=4 < 5$. $\textit{lo}=5$.
  \item $\textit{lo}=5 > \textit{hi}=4$: loop ends.
\end{itemize}

The answer is $\textit{lo} = 5$. Checking by hand: sorting all $9$
elements gives $[1,2,3,3,5,6,6,9,9]$, and the middle ($5$th) element is
indeed $5$.

\subsection*{C++ implementation}

\begin{lstlisting}
#include <bits/stdc++.h>
using namespace std;

int upperBoundRow(vector<int>& row, int X) {
    int lo = 0, hi = row.size() - 1, ans = row.size();
    while (lo <= hi) {
        int mid = lo + (hi - lo) / 2;
        if (row[mid] > X) { ans = mid; hi = mid - 1; }
        else lo = mid + 1;
    }
    return ans;
}

int matrixMedian(vector<vector<int>>& matrix) {
    int n = matrix.size(), m = matrix[0].size();
    int lo = INT_MAX, hi = INT_MIN;
    for (auto& row : matrix) {
        lo = min(lo, row[0]);
        hi = max(hi, row[m - 1]);
    }

    int half = (n * m) / 2 + 1; // ceil(nm/2) for odd nm

    while (lo <= hi) {
        int mid = lo + (hi - lo) / 2;
        int count = 0;
        for (auto& row : matrix) {
            count += upperBoundRow(row, mid);
        }

        if (count < half) {
            lo = mid + 1;
        } else {
            hi = mid - 1;
        }
    }
    return lo;
}

int main() {
    vector<vector<int>> matrix = {
        {1, 3, 5},
        {2, 6, 9},
        {3, 6, 9}
    };
    cout << "Median: " << matrixMedian(matrix) << endl;
    return 0;
}
\end{lstlisting}

\begin{complexitybox}
\textbf{Time Complexity.} The candidate value range halves each iteration
($O(\log(\max - \min))$ iterations), and each feasibility check sums a
per-row binary search costing $O(\log m)$ across $n$ rows, giving an
overall time complexity of
\[
  O(n \log m \log(\max-\min)),
\]
substantially better than sorting all $nm$ elements directly.
\textbf{Space Complexity.} $O(1)$ auxiliary space.
\end{complexitybox}

\begin{takeawaybox}
This closing problem is a deliberate capstone: it requires recognizing
that ``binary search on the answer'' (search over a value range using a
monotonic count predicate) and ``per-row binary search'' (exploiting
row-sortedness inside the predicate) are not mutually exclusive
techniques -- they nest naturally, with the outer search handling
\emph{which} value to test and the inner per-row searches handling
\emph{how cheaply} to test it. Recognizing when two patterns from
different parts of this chapter can be combined, rather than always
reaching for a single isolated template, is the most advanced skill this
chapter has to offer.
\end{takeawaybox}

